\documentclass[11pt]{amsart}

\usepackage{fontspec}
\input{komorebi-font-paths.tex}
\setmainfont{P052-Roman.otf}[Path=\KomorebiFontDirectory,BoldFont=P052-Bold.otf,ItalicFont=P052-Italic.otf,BoldItalicFont=P052-BoldItalic.otf,SmallCapsFont=P052-Roman.otf]
\setsansfont{DejaVu Sans}
\setmonofont{DejaVu Sans Mono}[Scale=MatchLowercase]
\usepackage[letterpaper,margin=1in]{geometry}
\usepackage{hyperref}
\usepackage{xurl}
\usepackage{newunicodechar}
\newfontfamily{\mathSymbolFont}{DejaVu Sans}
\newunicodechar{≅}{{\mathSymbolFont ≅}}
\newunicodechar{↦}{{\mathSymbolFont ↦}}
\newunicodechar{∘}{{\mathSymbolFont ∘}}
\newunicodechar{∤}{{\mathSymbolFont ∤}}
\usepackage{fvextra}
\DefineVerbatimEnvironment{verbatim}{Verbatim}{fontsize=\small,breaklines=true,breakanywhere=true}
\newcommand{\tightlist}{\setlength{\itemsep}{0pt}\setlength{\parskip}{0pt}}
\newcommand{\passthrough}[1]{#1}
\newlength{\pandocwidth}
\setlength{\pandocwidth}{\linewidth}
\usepackage[nopatch=footnote]{microtype}
\usepackage{amssymb,mathtools}
\usepackage{bm}
\usepackage{aliascnt}
\usepackage{enumitem}
\usepackage{graphicx}
\usepackage{booktabs}
\usepackage{array}
\usepackage{longtable}
\usepackage{needspace}
\usepackage[nameinlink,noabbrev]{cleveref}

\allowdisplaybreaks
\hyphenpenalty=10000
\exhyphenpenalty=10000
\tolerance=2600
\emergencystretch=4em



\setlist[enumerate]{topsep=4pt,itemsep=2pt}
\setlist[itemize]{topsep=4pt,itemsep=2pt}

\newtheorem{theorem}{Theorem}[section]
\newaliascnt{proposition}{theorem}
\newtheorem{proposition}[proposition]{Proposition}
\aliascntresetthe{proposition}
\newaliascnt{lemma}{theorem}
\newtheorem{lemma}[lemma]{Lemma}
\aliascntresetthe{lemma}
\newaliascnt{corollary}{theorem}
\newtheorem{corollary}[corollary]{Corollary}
\aliascntresetthe{corollary}
\theoremstyle{definition}
\newaliascnt{definition}{theorem}
\newtheorem{definition}[definition]{Definition}
\aliascntresetthe{definition}
\newaliascnt{example}{theorem}
\newtheorem{example}[example]{Example}
\aliascntresetthe{example}
\newaliascnt{standardfact}{theorem}
\newtheorem{standardfact}[standardfact]{Standard Fact}
\aliascntresetthe{standardfact}
\newaliascnt{conjecture}{theorem}
\newtheorem{conjecture}[conjecture]{Conjecture}
\aliascntresetthe{conjecture}
\theoremstyle{remark}
\newaliascnt{remark}{theorem}
\newtheorem{remark}[remark]{Remark}
\aliascntresetthe{remark}

\crefname{theorem}{theorem}{theorems}
\crefname{proposition}{proposition}{propositions}
\crefname{lemma}{lemma}{lemmas}
\crefname{corollary}{corollary}{corollaries}
\crefname{definition}{definition}{definitions}
\crefname{example}{example}{examples}
\crefname{standardfact}{standard fact}{standard facts}
\crefname{remark}{remark}{remarks}
\crefname{conjecture}{conjecture}{conjectures}

\newcommand{\CC}{\mathbb C}
\newcommand{\FF}{\mathbb F}
\newcommand{\ZZ}{\mathbb Z}
\newcommand{\QQ}{\mathbb Q}
\newcommand{\OO}{\mathcal O}
\newcommand{\GL}{\operatorname{GL}}
\newcommand{\SL}{\operatorname{SL}}
\newcommand{\Sp}{\operatorname{Sp}}
\newcommand{\Ind}{\operatorname{Ind}}
\newcommand{\cInd}{\operatorname{cInd}}
\newcommand{\Res}{\operatorname{Res}}
\newcommand{\Tr}{\operatorname{Tr}}
\newcommand{\Frob}{\operatorname{Frob}}
\newcommand{\Gal}{\operatorname{Gal}}
\newcommand{\Hom}{\operatorname{Hom}}
\newcommand{\End}{\operatorname{End}}
\newcommand{\Irr}{\operatorname{Irr}}
\newcommand{\rad}{\operatorname{rad}}
\newcommand{\rank}{\operatorname{rank}}
\newcommand{\ord}{\operatorname{ord}}
\newcommand{\lcm}{\operatorname{lcm}}
\newcommand{\Hol}{\operatorname{Hol}}
\newcommand{\im}{\operatorname{im}}
\newcommand{\coker}{\operatorname{coker}}
\newcommand{\Art}{\operatorname{Art}}
\newcommand{\Swan}{\operatorname{Swan}}
\newcommand{\Arf}{\operatorname{Arf}}
\newcommand{\one}{\bm 1}
\newcommand{\e}{\mathrm e}
\newcommand{\abs}[1]{\left\lvert#1\right\rvert}
\newcommand{\norm}[1]{\left\lVert#1\right\rVert}
\newcommand{\calH}{\mathcal H}
\newcommand{\calV}{\mathcal V}
\newcommand{\calW}{\mathcal W}
\newcommand{\calP}{\mathcal P}
\newcommand{\calA}{\mathcal A}
\newcommand{\calC}{\mathcal C}
\newcommand{\calG}{\mathcal G}
\newcommand{\calX}{\mathcal X}
\newcommand{\calT}{\mathcal T}
\newcommand{\frakH}{\mathfrak H}
\newcommand{\frakA}{\mathfrak A}
\newcommand{\frakC}{\mathfrak C}
\newcommand{\perpbig}{\mathop{\perp}\limits}

\graphicspath{{figures/}}

\title[Tame and wild packet geometry]{Tame and Wild Packet Geometry for Local Langlands}
\author[Komorebi P.]{Komorebi P.}
\subjclass[2020]{11F70, 20C33, 22E50, 11S37, 20C25}
\keywords{local Langlands correspondence, tame induction, Heisenberg representations,
wild inertia, quadratic refinements, holonomy, arithmetic packet realization}
\date{}

\pdfstringdefDisableCommands{\let\unskip\relax}
\hypersetup{
  pdftitle={Tame and Wild Packet Geometry for Local Langlands},
  pdfauthor={Komorebi P.},
  pdfsubject={Tame Frobenius packets, wild Heisenberg envelopes, finite Weil phase, matrix holonomy, and restricted local bridges}
}


% Operation Komorebi: shared publication profile.
\usepackage{komorebi-publication}
\renewcommand{\KomorebiTitle}{Tame and Wild Packet Geometry}
\renewcommand{\KomorebiSubtitle}{For local Langlands}
\renewcommand{\KomorebiSeries}{Chapter XI}
\renewcommand{\KomorebiKind}{Mathematical text}
\renewcommand{\KomorebiEdition}{}
\renewcommand{\KomorebiShortTitle}{Paper XI}
\renewcommand{\KomorebiTitleSize}{38}
\renewcommand{\KomorebiTitleLeading}{43}
\renewcommand{\KomorebiProfile}{classical}
\renewcommand{\KomorebiOrientation}{portrait}
\KomorebiSetup
\usepackage{newunicodechar}
\newunicodechar{∩}{\ensuremath{\cap}}
\newunicodechar{⊂}{\ensuremath{\subset}}
\newunicodechar{∈}{\ensuremath{\in}}
\newunicodechar{⊥}{\ensuremath{\perp}}
% Standard numberline TOC entries are written by titlesec. Give them a real number box.
\makeatletter
\providecommand\@dotsep{4.5}
\renewcommand*\l@section[2]{\@dottedtocline{1}{0em}{2.8em}{#1}{#2}}
\makeatother
\raggedbottom
\input{core-monograph-layout.tex}
\begin{document}
\KomorebiCover

\special{pdf:trailerid [<0e9d864536f570d7480471e0f99695cb><0e9d864536f570d7480471e0f99695cb>]}
\begin{abstract}
Let \(K\) be a nonarchimedean local field. We resolve several finite layers that occur inside the local Langlands correspondence for general linear groups. Frobenius semisimple Weil Deligne parameters first split into irreducible Weil atoms and finite monodromy chains. In the tame case, character projectors show that every irreducible atom is one transitive Frobenius orbit, and Mackey theory gives an exact greatest common divisor and least common multiple rule for tensor packets. Regular finite field characters produce primitive cuspidal projectors and finite Gauss phases. For wild data, projective intertwiners recover an alternating commutator form. Its radical free quotient determines a finite Heisenberg representation, and maximal isotropic averaging gives an exact selector. We prove a restricted one break bridge: an equivariant symplectic isometry between a declared Galois envelope and a type side quotient induces the Heisenberg comparison uniquely up to scalar, while one cyclic phase condition removes the remaining ambiguity. Odd residue characteristic is governed by finite Weil phases. Residue characteristic two requires a quadratic refinement and its Arf invariant. A matrix holonomy theorem then gives the exact gluing criterion on a finite transport graph. The general arithmetic packet realization is stated as Conjecture 66. Detailed ordinary proofs give restricted arithmetic realizations in all odd prime-power ranks and in all dyadic ranks for the stated increasing-break families. The complete comparison of pair local factors remains a conditional input for the same arithmetic candidate.
\end{abstract}
\makeatletter\@setabstract\makeatother
\thispagestyle{plain}
\clearpage
\setcounter{tocdepth}{1}
\begingroup\setlength{\parskip}{0pt}
\tableofcontents
\endgroup
\clearpage
\begingroup\raggedright
\section*{A route through Chapter XI}
\phantomsection
\addcontentsline{toc}{section}{A route through Chapter XI}
This chapter presents the finite packet theory and retains the detailed restricted arithmetic constructions. Two separate research companions explain the later progress toward Conjecture 66. They do not replace this chapter's proved statements or turn its conditional comparison into an existence theorem.

\subsection*{The mathematical text and its companions}
\begin{tabular}{@{}>{\raggedright\arraybackslash}p{.28\linewidth}>{\raggedright\arraybackslash}p{.67\linewidth}@{}}\toprule
Reader & Route and purpose\\\midrule
Mathematical text & Finite tame and wild theory, restricted arithmetic constructions, and the exact realization/comparison boundary.\\[7pt]
Core visual guide & \href{https://komorebi.invalid/b07/core_visual}{Geometric route through the mathematical text}. The original arithmetic examples retain their stated scope.\\[7pt]
Formal companion & \href{https://komorebi.invalid/b07/core_lean}{Formal statements, literal excerpts and evidence lineage}. The guide distinguishes the selected source census, compiler replay and arithmetic hypotheses.\\[7pt]
Progress paper & \href{https://komorebi.invalid/b07/research_start}{Toward Arithmetic Packet Realization}. Connected proofs, failed routes, conductor conventions and native-marking conditions.\\[7pt]
Research visual atlas & \href{https://komorebi.invalid/b07/atlas_start}{Arithmetic Packet Realization}. Thirty-two worked plates accompany the progress paper.\\\bottomrule
\end{tabular}

\subsection*{The endpoint is a quantified construction problem}
Conjecture 66 asks for an actual arithmetic witness for each admitted source with its supplied H1--H4 datum. The comparison proposition then uses H5 for that same realized candidate. The formal companion distinguishes this conjecture from the differently numbered statements in the source registers.

The \href{https://komorebi.invalid/b07/contract}{research contract in Section 1} and \href{https://komorebi.invalid/b07/contract_atlas}{Plate XI.1} make the quantifiers explicit. Supplying finite operators does not construct an arithmetic representation. Likewise, proving a factor identity for a different candidate does not identify the packet's realizer.

\subsection*{How to use the links}
The reading links connect the mathematical text, visual studies, research companions and formal companion. Section and result numbers remain local to the named treatment. Plate references include the chapter number, so a reference such as Plate XI.32 identifies a visual plate rather than a source version.
\clearpage
\section*{Scope and research routes}
\phantomsection
\addcontentsline{toc}{section}{Scope and research routes}
\begin{tabular}{@{}>{\raggedright\arraybackslash}p{.32\linewidth}>{\raggedright\arraybackslash}p{.63\linewidth}@{}}\toprule
Question & Route beyond the core statement\\\midrule
Parameter, full character and recovery & \href{https://komorebi.invalid/b07/arithmetic}{Progress Sections 7--10} and \href{https://komorebi.invalid/b07/arithmetic_atlas}{Plates XI.6--XI.10}. The full character, central character and noncompact completion remain fixed.\\[8pt]
Monomial, Clifford and occurrence routes & \href{https://komorebi.invalid/b07/monomial}{Progress Part II}. The original character route is excluded under its hypotheses. A different corrected character passes a scalar test, but occurrence in the same candidate is still open.\\[8pt]
Conductor and capacity & \href{https://komorebi.invalid/b07/conductor}{Progress Sections 22--25} and \href{https://komorebi.invalid/b07/conductor_atlas}{Plates XI.11--XI.19}. Swan, ordinary conductor and additive-character epsilon exponent are named separately. Equal numerical invariants need not give equal quotient capacity.\\[8pt]
Native markings and allowed repairs & \href{https://komorebi.invalid/b07/markings}{Progress Sections 26--29} and \href{https://komorebi.invalid/b07/markings_atlas}{Plates XI.20--XI.29}. A freely transported marking is different from a prescribed native marking. Finite corrections require their stated inputs and permission to change the relevant map.\\[8pt]
Remaining arithmetic certificate & \href{https://komorebi.invalid/b07/remaining}{Progress Section 30} and \href{https://komorebi.invalid/b07/remaining_atlas}{Plates XI.30--XI.32}. General realization, H5, corrected-character occurrence and the stronger native certificate remain separate tasks.\\\bottomrule
\end{tabular}

\subsection*{The construction and its limits}
The breakwise phase product uses its moving symplectic spaces. The type-side
pairing fixes a prime-field character, and the tame comparison retains the
rank-one conductor and trivial-character Gauss-sum boundaries. Interaction
carriers are distinguished from minimal character orbits.

The detailed construction appendices give ordinary proofs under their stated
type-theory hypotheses. The selected Henniart--Lemaire induction inputs
\cite{HenniartLemaire2011} retain their characteristic, rank and normalization
conditions. They do not establish H5 or corrected-character occurrence.
The formal companion identifies the actual compiler and axiom coverage.

\endgroup

\clearpage

\section{Introduction}\label{sec:introduction}

The local Langlands correspondence for \(\GL_n(K)\) is an established theorem. It matches irreducible smooth representations of \(\GL_n(K)\) with Frobenius semisimple Weil Deligne representations of dimension \(n\). Its proof and normalization draw on global methods, harmonic analysis, types, local factors, and the geometry of Shimura varieties \cite{Henniart2000,HarrisTaylor2001,BushnellHenniart2006}. This chapter asks a narrower structural question. Which parts of the local correspondence admit finite packet coordinates whose projectors, transfers, phases, and gluing laws can be proved directly?

The answer has a sharp division. The tame layers are explicit. A tame irreducible Weil representation is induced from a character of an unramified extension. Its inertia restriction is one Frobenius orbit. Regular characters of elliptic finite field tori give cuspidal blocks through the Deligne and Lusztig construction \cite{DeligneLusztig1976}. Their tensor products are controlled by a greatest common divisor and least common multiple decomposition. The corresponding cycle orientation and finite Gauss phases can be computed exactly. These features are closely related to the tame lifting and essentially tame correspondence of Bushnell and Henniart \cite{BushnellHenniart1996,BushnellHenniart1999,BushnellHenniart2005ET1,BushnellHenniart2005ET2,BushnellHenniart2010ET3}.

Wild data are different. The restriction to wild inertia need not split into cyclic characters. On the type side, simple character theory produces finite symplectic quotients and Heisenberg representations \cite{BushnellKutzko1993}. Canonical beta extensions fix a phase ambiguity in the extension to the normalizer \cite{Blondel2012,Dotto2022}. On the Galois side, higher ramification gives a canonical filtration and strong numerical constraints \cite{BushnellHenniart2017}. The missing construction is not another finite matrix calculation. It is a choice free passage from an arbitrary wild inertia atom to a finite symplectic envelope that agrees with the type side quotient.

We therefore separate three levels of assertion. First, this chapter proves finite algebraic theorems about tame orbits, Heisenberg groups, twisted group algebras, finite Weil phases, and transport graphs. Second, it proves a restricted one break bridge once a normal wild datum and an equivariant symplectic identification are declared. Third, it records the exact additional input needed to pass from that restricted bridge to an independent construction of all wild atoms.

The finite Heisenberg theory used here is classical. We use the Stone von Neumann theorem, finite Weil representations, and standard projective representation theory \cite{Gerardin1977,Neuhauser2002,Prasad2009,Karpilovsky1985,GurevichHadani2007}. The contribution is the way these results are arranged into one comparison diagram and the matrix holonomy theorem that controls finite gluing. In residue characteristic two, the alternating form does not determine the phase. A quadratic refinement and its Arf invariant are required. The underlying quadratic form theory is standard \cite{MilnorHusemoller1973}.

Local factors provide the final test. Henniart showed that pair epsilon factors characterize the local correspondence in the required range \cite{Henniart1993}. The tame packet calculations below recover exact local polynomials and finite Gauss phases for the declared tame atoms. They do not supply the complete pair factor comparison for arbitrary wild atoms. This boundary is retained throughout this chapter.

\begin{figure}[t]
\centering
\includegraphics[width=.98\textwidth]{01_architecture.pdf}
\caption{The finite theorem layers and the remaining wild boundary.}
\label{fig:architecture}
\end{figure}

Arithmetic realization and identification of an actual candidate are
separate problems. The realization endpoint is
\cref{conj:arithmetic_realization}. The restricted constructions are developed
in \cref{sec:arithmetic_progress} and the detailed construction appendices.
They have ordinary mathematical proofs with retained internal reviews.
Their formal coverage is stated separately in the formal companion.

The accompanying sources include detailed proofs, construction notes,
proof reviews and research records.
These documents distinguish established results, conditional results, and
the remaining obligations associated with Conjecture 66.

This chapter is organized as follows. Weil Deligne chains appear in \cref{sec:wd}. Tame Frobenius packets, finite field cuspidal projectors, tensor refinement, and Gauss phases are developed in \cref{sec:tame,sec:finitefield,sec:tensor,sec:gauss}. Finite Heisenberg groups and projective commutator envelopes appear in \cref{sec:heisenberg,sec:twisted}. The restricted one break Galois envelope and its comparison with the type side are treated in \cref{sec:onebreak,sec:typeside,sec:bridge}. Odd and dyadic phase data are separated in \cref{sec:oddphase,sec:dyadic}. Matrix holonomy and breakwise reconstruction appear in \cref{sec:holonomy,sec:breakwise}. The local-factor boundary appears in \cref{sec:boundary}. Examples and scope
follow in \cref{sec:examples,sec:scope}. Section~\ref{sec:arithmetic_progress}
and the appendices develop the restricted arithmetic constructions. The
separate progress paper reconciles subsequent routes, corrections and
remaining obligations without replacing these source records.

\section{Weil Deligne chains}\label{sec:wd}

Let \(W_K\) be the Weil group of \(K\), let \(I_K\) be inertia, let \(P_K\) be wild inertia, and choose a geometric Frobenius lift \(\Phi\). A Weil Deligne representation is a pair \((r,N)\) with
\[
r(w)Nr(w)^{-1}=\abs{w}_K N.
\]

\begin{standardfact}[Indecomposable monodromy blocks]\label{fact:spdecomp}
Every Frobenius semisimple Weil Deligne representation is a direct sum of blocks
\[
\operatorname{Sp}_m(\tau),
\]
where \(\tau\) is an irreducible Weil representation and the nilpotent operator is one Jordan chain of length \(m\).
\end{standardfact}

This standard decomposition reduces the construction of a local correspondence to two tasks. One must understand irreducible Weil atoms, then one must attach the correct generalized Steinberg chain. The second task is formal once the first is known.

\begin{definition}[Monodromy packet chain]\label{def:monodromy_chain}
Let \(V_0,\ldots,V_{m-1}\) be copies of the space of \(\tau\). Put
\[
V=\bigoplus_{j=0}^{m-1}V_j.
\]
A monodromy packet chain is a nilpotent operator \(N\) with
\[
N(V_j)\subseteq V_{j-1}
\]
for \(j>0\), with \(N(V_0)=0\), together with a Frobenius operator whose conjugation scales \(N\) by \(q^{-1}\).
\end{definition}

\begin{proposition}[Exact block reconstruction]\label{prop:monodromy_reconstruction}
Let \(P_j\) be the projector onto \(V_j\). Then
\[
N=\sum_{j=1}^{m-1}P_{j-1}NP_j,
\qquad
N^m=0.
\]
Every power \(N^r\) is the sum of the packet paths of length \(r\) along the chain.
\end{proposition}

\begin{proof}
The projectors resolve the identity. Since \(N\) maps only from \(V_j\) to \(V_{j-1}\), every block except the displayed nearest neighbor blocks vanishes. A product of \(r\) nonzero blocks moves down exactly \(r\) levels. No path of length \(m\) exists, so \(N^m=0\).
\end{proof}

\begin{proposition}[Scalar covariance model]\label{prop:scalar_covariance}
Let
\[
F=\operatorname{diag}(1,q,q^2,\ldots,q^{m-1})
\]
and let \(N_{j-1,j}=1\). Then
\[
FNF^{-1}=q^{-1}N.
\]
\end{proposition}

\begin{proof}
The only nonzero entry in column \(j\) and row \(j-1\) is multiplied by \(q^{j-1}q^{-j}=q^{-1}\) under conjugation.
\end{proof}

\begin{figure}[t]
\centering
\includegraphics[width=.9\textwidth]{02_weil_deligne_chain.pdf}
\caption{A Weil Deligne chain is reconstructed from one atom and finitely many nearest neighbor transfers.}
\label{fig:wdchain}
\end{figure}

\begin{theorem}[Formal atomic reconstruction]\label{thm:formal_atomic}
Assume that every irreducible Weil representation \(\tau\) has been assigned a supercuspidal representation \(\sigma_\tau\), compatibly with isomorphism and unramified twist. Then the rules
\[
\operatorname{Sp}_m(\tau)
\longmapsto
\operatorname{St}_m(\sigma_\tau)
\]
and normalized parabolic induction determine a candidate correspondence on all Frobenius semisimple Weil Deligne parameters.
\end{theorem}

\begin{proof}
Decompose the parameter into indecomposable monodromy blocks. Replace each Weil atom by its assigned supercuspidal atom and each monodromy chain by the corresponding generalized Steinberg block. Order the resulting essentially square integrable factors in Langlands order and take the unique irreducible quotient. The multiset of indecomposable Weil Deligne blocks is intrinsic, so the result does not depend on the chosen display.
\end{proof}

\begin{remark}\label{rem:atomic_boundary}
\Cref{thm:formal_atomic} is a reconstruction theorem. It does not produce the atomic assignment \(\tau\mapsto\sigma_\tau\), and it does not prove the required local factor identities.
\end{remark}

\section{Tame Frobenius packets}\label{sec:tame}

Assume that \(P_K\) acts trivially. After passing to a finite quotient, tame inertia is cyclic. Let \(t\) generate a quotient of order \(M\), and put \(\zeta_M=\e^{2\pi i/M}\). For \(j\in\ZZ/M\ZZ\), define
\[
P_j
=
\frac1M
\sum_{a=0}^{M-1}
\zeta_M^{-ja}r(t)^a.
\]

\begin{proposition}[Tame character projectors]\label{prop:tame_projectors}
The operators \(P_j\) are pairwise orthogonal idempotents whose sum is the identity. Frobenius permutes them by
\[
\Phi P_j\Phi^{-1}=P_{qj}.
\]
\end{proposition}

\begin{proof}
The first assertion is the Fourier idempotent calculation for a cyclic group. For the geometric Frobenius convention of \cref{sec:wd}, the relation is \(\Phi^{-1}t\Phi=t^q\). Equivalently, \(\Phi t\Phi^{-1}=t^{q^{-1}}\), with inverse taken modulo \(M\) since \(\gcd(q,M)=1\). Conjugating the Fourier sum and changing its summation index therefore gives \(\Phi P_j\Phi^{-1}=P_{qj}\). The inverse Frobenius convention reverses the label direction while preserving the same orbits.
\end{proof}

\begin{theorem}[Tame orbit induction]\label{thm:tame_induction}
Let \(\tau\) be an irreducible complex representation of \(W_K\) with trivial wild inertia. Then the nonzero tame character packets form one transitive Frobenius orbit of length \(f\). If \(K_f/K\) is the unramified extension of degree \(f\), there is a character \(\xi\) of \(W_{K_f}\) such that
\[
\boxed{
\tau\cong\Ind_{W_{K_f}}^{W_K}\xi.
}
\]
\end{theorem}

\begin{proof}
The nonzero tame packets are permuted by Frobenius. A union of Frobenius orbits is stable under \(W_K\), so irreducibility forces one orbit. Choose one packet \(V_\theta\). Its stabilizer is \(W_{K_f}=\langle I_K,\Phi^f\rangle\). A nonzero proper stable subspace of \(V_\theta\) would generate a nonzero proper \(W_K\) stable subspace after summing its Frobenius translates. Thus \(V_\theta\) is irreducible for the stabilizer. Inertia acts on it by the scalar character \(\theta\), and the remaining action is generated by \(\Phi^f\). The stabilizer action on this packet is therefore one dimensional. Induction reconstructs the complete orbit.
\end{proof}

\begin{corollary}[Dimension and regularity]\label{cor:tame_dimension}
The dimension of \(\tau\) equals the Frobenius orbit length of its tame inertia character. The inducing character is regular, meaning that no smaller positive power of Frobenius fixes it.
\end{corollary}

\begin{proof}
The induced representation has one line for every Frobenius translate of the chosen character. A shorter orbit would produce a proper stabilizer and contradict the minimality of \(f\).
\end{proof}


\begin{proposition}[Dual and self dual tame packets]\label{prop:tame_dual}
For a tame atom
\[
\tau_\xi=\Ind_{W_{K_f}}^{W_K}\xi,
\]
we have
\[
\tau_\xi^\vee
\cong
\Ind_{W_{K_f}}^{W_K}\xi^{-1}.
\]
The atom is self dual exactly when
\[
\xi^{-1}=\xi^{\Phi^j}
\]
for some \(j\in\ZZ/f\ZZ\).
\end{proposition}

\begin{proof}
Duality commutes with finite index induction, giving the first statement. Two regular induced tame atoms are isomorphic exactly when their inducing characters lie in the same Frobenius orbit. Applying this to \(\xi\) and \(\xi^{-1}\) gives the criterion.
\end{proof}

\begin{proposition}[Unramified twist]\label{prop:tame_twist}
Let \(\nu\) be an unramified character of \(W_K\). Then
\[
\tau_\xi\otimes\nu
\cong
\Ind_{W_{K_f}}^{W_K}
\left(
\xi\,\nu|_{W_{K_f}}
\right).
\]
The tame inertia orbit is unchanged. If the original Frobenius cycle has closing scalar \(a\), the twisted cycle has closing scalar
\[
a\nu(\Phi)^f.
\]
\end{proposition}

\begin{proof}
The tensor identity follows from the projection formula for induction. Since \(\nu\) is trivial on inertia, it does not change any tame character packet. It multiplies every Frobenius step by \(\nu(\Phi)\), so a closed cycle of length \(f\) gains the factor \(\nu(\Phi)^f\).
\end{proof}

\begin{figure}[t]
\centering
\includegraphics[width=.62\textwidth]{03_tame_frobenius_orbit.pdf}
\caption{A tame irreducible atom is one transitive Frobenius orbit.}
\label{fig:tameorbit}
\end{figure}

\subsection{Regular orbit counts}

Characters of \(\FF_{q^n}^{\times}\) are indexed by exponents modulo \(q^n-1\), and Frobenius acts by multiplication by \(q\).

\begin{proposition}[Möbius orbit count]\label{prop:mobius_count}
The number of regular Frobenius orbits of length \(n\) is
\[
\boxed{
c_n(q)
=
\frac1n
\sum_{d\mid n}
\mu(d)
\bigl(q^{n/d}-1\bigr).
}
\]
\end{proposition}

\begin{proof}
For every divisor \(r\mid n\), the characters fixed by \(\Frob^r\) are exactly those descending to \(\FF_{q^r}^{\times}\), and there are \(q^r-1\) of them. Möbius inversion gives the number of characters of exact orbit length \(n\). Division by \(n\) gives the number of orbits.
\end{proof}

\begin{example}[Degree four over \(\FF_2\)]\label{ex:degree_four}
The regular exponent orbits modulo fifteen are
\[
(1,2,4,8),
\qquad
(3,6,12,9),
\qquad
(7,14,13,11).
\]
The formula gives
\[
\frac14\bigl((2^4-1)-(2^2-1)\bigr)=3.
\]
\end{example}

\begin{figure}[t]
\centering
\includegraphics[width=.76\textwidth]{04_regular_orbit_counts.pdf}
\caption{Regular Frobenius orbit counts for three residue field sizes.}
\label{fig:orbitcounts}
\end{figure}

\section{Finite field cuspidal projectors}\label{sec:finitefield}

Let
\[
G_n=\GL_n(\FF_q),
\qquad
T_n\cong\FF_{q^n}^{\times}
\]
with \(T_n\) an elliptic maximal torus.

\begin{standardfact}[Regular elliptic torus characters]\label{fact:deligne_lusztig}
A regular character \(\theta:T_n\to\CC^{\times}\) determines an irreducible cuspidal representation \(\kappa_\theta\) of \(G_n\). Frobenius conjugate characters determine isomorphic representations, and every irreducible cuspidal representation of \(G_n\) occurs in this way.
\end{standardfact}

This is the finite field input from Deligne and Lusztig theory \cite{DeligneLusztig1976}. We use standard finite group character theory for the central idempotents \cite{Serre1977,Isaacs1976}. Once \(\kappa_\theta\) is known, its block is selected by a central idempotent.

\begin{proposition}[Primitive cuspidal idempotent]\label{prop:cuspidal_idempotent}
Let \(\chi_\theta\) be the irreducible character of \(\kappa_\theta\), and let \(d_\theta=\chi_\theta(1)\). Then
\[
e_\theta
=
\frac{d_\theta}{\abs{G_n}}
\sum_{g\in G_n}
\chi_\theta(g^{-1})g
\]
is the primitive central idempotent selecting the \(\kappa_\theta\) block in \(\CC[G_n]\).
\end{proposition}

\begin{proof}
On an irreducible representation of character \(\psi\), the central element acts by a scalar. Taking traces shows that the scalar is
\[
\frac{d_\theta}{\abs{G_n}\psi(1)}
\sum_{g\in G_n}
\chi_\theta(g^{-1})\psi(g).
\]
Character orthogonality makes this scalar one for \(\psi=\chi_\theta\) and zero otherwise. The element is therefore central, idempotent, and primitive.
\end{proof}

\begin{proposition}[Parabolic annihilation]\label{prop:parabolic_annihilation}
Let \(P=LU\) be a proper parabolic subgroup of \(G_n\), and put
\[
e_U=\frac1{\abs U}\sum_{u\in U}u.
\]
If \(\kappa_\theta\) is cuspidal, then
\[
\boxed{
e_Ue_\theta=e_\theta e_U=0.
}
\]
\end{proposition}

\begin{proof}
The operator \(e_U\) projects onto the subspace of \(U\) fixed vectors. A cuspidal representation has no nonzero fixed vector under the unipotent radical of a proper parabolic. Thus \(e_U\) acts by zero on the \(\kappa_\theta\) block. Since \(e_\theta\) is central, both products vanish.
\end{proof}

\begin{figure}[t]
\centering
\includegraphics[width=.92\textwidth]{05_cuspidal_projector.pdf}
\caption{A regular torus character determines a cuspidal block and its primitive central projector.}
\label{fig:cuspidal}
\end{figure}

\subsection{Depth zero reconstruction and cycle orientation}

The depth zero construction inflates a finite cuspidal representation to \(\GL_n(\OO_K)\), extends it to the standard compact modulo center subgroup, and then compactly induces it to \(\GL_n(K)\). This is classical type theory \cite{BushnellKutzko1993,BushnellHenniart2006}.

The induced tame Weil atom has an elementary cycle invariant. In an orbit basis \(e_0,\ldots,e_{n-1}\), write
\[
\tau(\Phi)e_j=e_{j+1}
\]
for \(j<n-1\), and
\[
\tau(\Phi)e_{n-1}=ae_0.
\]

\begin{proposition}[Frobenius cycle polynomial]\label{prop:frobenius_cycle}
The Frobenius transfer matrix has characteristic polynomial
\[
X^n-a
\]
and determinant
\[
\boxed{
\det\tau(\Phi)=(-1)^{n-1}a.
}
\]
\end{proposition}

\begin{proof}
The matrix is the companion matrix of \(X^n-a\). Its only nonzero determinant term is the closed \(n\) cycle. That cycle has sign \((-1)^{n-1}\), and the closing edge contributes \(a\).
\end{proof}

\begin{corollary}[Unramified orientation character]\label{cor:orientation}
The sign \((-1)^{n-1}\) is the orientation of the closed Frobenius packet cycle. It is unchanged by rescaling the individual orbit basis vectors.
\end{corollary}

\begin{proof}
A basis rescaling conjugates the transfer matrix by a diagonal matrix, so the determinant and characteristic polynomial remain unchanged.
\end{proof}

\section{Tensor refinement of tame atoms}\label{sec:tensor}

Let
\[
\tau_\xi=\Ind_{W_{K_n}}^{W_K}\xi,
\qquad
\tau_\eta=\Ind_{W_{K_m}}^{W_K}\eta,
\]
where \(K_n/K\) and \(K_m/K\) are unramified. Put
\[
d=\gcd(n,m),
\qquad
L=\lcm(n,m).
\]

\begin{theorem}[Greatest common divisor and least common multiple packets]\label{thm:gcd_lcm}
There is a decomposition
\[
\boxed{
\tau_\xi\otimes\tau_\eta
\cong
\bigoplus_{i=0}^{d-1}
\Ind_{W_{K_L}}^{W_K}
\left(
\xi_{L,n}\eta_{L,m}^{\Phi^i}
\right),
}
\]
where \(\xi_{L,n}\) and \(\eta_{L,m}\) are the norm lifts to \(K_L\).
\end{theorem}

\begin{proof}
Apply Mackey decomposition \cite{Mackey1952}. The double cosets
\[
W_{K_n}\backslash W_K/W_{K_m}
\]
are the orbits of simultaneous translation on \(\ZZ/n\ZZ\times\ZZ/m\ZZ\). There are \(d\) orbits, each of length \(L\). The stabilizer of a point in each orbit is \(W_{K_L}\). The second character is shifted by the chosen Frobenius representative. Summing the induced terms over the double cosets gives the decomposition.
\end{proof}

\begin{corollary}[Packet count and dimension]\label{cor:gcd_lcm_count}
The tensor product has \(d\) induced interaction packets, each with carrier dimension and Frobenius coset-cycle length \(L\). The interaction character may have a shorter inertia orbit, in which case a packet may split further into irreducible atoms. The total dimension is preserved because
\[
dL=nm.
\]
\end{corollary}
\begin{proof}
The double coset calculation in \cref{thm:gcd_lcm} gives exactly \(d\) summands, and every stabilizer is \(W_{K_L}\), so every induced packet has dimension \(L\). The identity \(dL=nm\) follows from the definitions of greatest common divisor and least common multiple.
\end{proof}


\begin{figure}[t]
\centering
\includegraphics[width=.82\textwidth]{06_gcd_lcm_refinement.pdf}
\caption{Tensor refinement gives as many packets as the greatest common divisor, each with carrier dimension equal to the least common multiple.}
\label{fig:gcdlcm}
\end{figure}

\subsection{Exponent coordinates}

Represent the residue characters by exponents
\[
a\pmod{q^n-1},
\qquad
b\pmod{q^m-1}.
\]

\begin{proposition}[Refined exponent formula]\label{prop:refined_exponent}
The \(i\)th interaction packet on \(\FF_{q^L}^{\times}\) has exponent
\[
\boxed{
c_i
=
a\frac{q^L-1}{q^n-1}
+bq^i\frac{q^L-1}{q^m-1}
\pmod{q^L-1}.
}
\]
It is trivial on tame inertia exactly when \(c_i=0\).
\end{proposition}

\begin{proof}
Norm lifting a character exponent from \(\FF_{q^n}^{\times}\) to \(\FF_{q^L}^{\times}\) multiplies the exponent by \((q^L-1)/(q^n-1)\). The Frobenius shift multiplies the second exponent by \(q^i\). Multiplication of characters adds exponents. Triviality is equivalent to the resulting exponent being zero modulo \(q^L-1\).
\end{proof}

\subsection{Tame local polynomials}

Let \(V\) be one interaction packet. If its inertia character is nontrivial, then \(V^{I_K}=0\). If the inertia character is trivial, Frobenius acts on the invariant packet by a companion cycle.

\begin{proposition}[Local polynomial of one tame packet]\label{prop:tame_local_polynomial}
Suppose an interaction packet has orbit length \(r\) and closing scalar \(a\). Its local polynomial is
\[
\det\left(I-u\Phi\mid V^{I_K}\right)
=
\begin{cases}
1-a u^r,&\text{if the inertia character is trivial},\\
1,&\text{otherwise}.
\end{cases}
\]
\end{proposition}

\begin{proof}
For nontrivial inertia, the invariant space is zero. In the trivial case, Frobenius on the invariant packet is the companion matrix of \(X^r-a\). Substituting \(X=u^{-1}\) and multiplying by \(u^r\) gives \(1-a u^r\).
\end{proof}

\begin{corollary}[Tame tensor local factor]\label{cor:tame_tensor_factor}
The local factor of the tame tensor product is the product of
\[
(1-a_i u^{r_i})^{-1}
\]
over precisely those refined packets whose inertia exponent vanishes.
\end{corollary}

\begin{proof}
The invariant space is the direct sum of the invariant refined packets, and determinant is multiplicative on direct sums.
\end{proof}

\begin{proposition}[Trivial tensor packet criterion]\label{prop:trivial_tensor_packet}
For every refined interaction packet, let \(r_i\) be its cycle length and \(a_i\) its closing scalar. The dimension of the \(W_K\) fixed space in \(\tau_\xi\otimes\tau_\eta\) equals the number of packets satisfying both
\[
c_i=0
\qquad\text{and}\qquad
a_i=1.
\]
\end{proposition}

\begin{proof}
A nonzero fixed vector must lie in an inertia trivial packet, so \(c_i=0\). On that packet, Frobenius is a companion cycle with characteristic polynomial \(X^{r_i}-a_i\). Its fixed space is the eigenspace for eigenvalue one. This eigenspace is one dimensional exactly when \(a_i=1\), and it is zero otherwise. Summing over the refined packets proves the formula.
\end{proof}

\begin{corollary}[Tame Hom test]\label{cor:tame_hom}
For irreducible tame atoms,
\[
\dim\Hom_{W_K}(\tau_\xi^\vee,\tau_\eta)
\]
is zero or one, and it is one exactly when the refined tensor packet test in \cref{prop:trivial_tensor_packet} finds one trivial packet.
\end{corollary}

\begin{proof}
The Hom space is the fixed space of \(\tau_\xi\otimes\tau_\eta\). Schur's lemma gives dimension zero or one, and \cref{prop:trivial_tensor_packet} gives the packet criterion.
\end{proof}

\subsection{Tame conductor}

\begin{proposition}[Conductor of a nontrivial tame atom]\label{prop:tame_conductor}
Let \(\tau\) be an irreducible tame Weil representation of dimension \(n>1\). Then
\[
\Swan(\tau)=0,
\qquad
\dim\tau^{I_K}=0,
\qquad
\Art(\tau)=n.
\]
\end{proposition}

\begin{proof}
Wild inertia acts trivially, so the Swan conductor is zero. An inertia fixed vector would generate a nonzero Frobenius stable sum of trivial character packets. Irreducibility would then make inertia act trivially on all of \(\tau\). An irreducible representation of the quotient generated by Frobenius would be one dimensional, contrary to \(n>1\). The Artin conductor is therefore
\[
\dim\tau-\dim\tau^{I_K}+\Swan(\tau)=n.
\]
\end{proof}

\section{Finite Gauss phases}\label{sec:gauss}

Let \(k=\FF_q\), let \(\psi:k\to\CC^{\times}\) be a nontrivial additive character, and let \(\chi:k^{\times}\to\CC^{\times}\) be a multiplicative character. Define
\[
G(\chi,\psi)
=
\sum_{x\in k^{\times}}\chi(x)\psi(x).
\]
The standard theory of Gauss and Jacobi sums will be used in the form recorded in \cite{DavenportHasse1935,BerndtEvansWilliams1998}.

\begin{proposition}[Gauss magnitude]\label{prop:gauss_magnitude}
If \(\chi\) is nontrivial, then
\[
\boxed{
\abs{G(\chi,\psi)}=q^{1/2}.
}
\]
\end{proposition}

\begin{proof}
Compute
\[
\abs{G(\chi,\psi)}^2
=
\sum_{x,y\in k^{\times}}
\chi(xy^{-1})\psi(x-y).
\]
Write \(x=ty\). The inner sum over \(y\) equals \(q-1\) when \(t=1\) and equals \(-1\) otherwise. Since \(\sum_{t\in k^{\times}}\chi(t)=0\), the result is \(q\).
\end{proof}

\begin{definition}[Normalized Gauss phase]\label{def:gauss_phase}
For nontrivial \(\chi\), define
\[
\gamma(\chi,\psi)
=
q^{-1/2}G(\chi,\psi).
\]
Then \(\abs{\gamma(\chi,\psi)}=1\).
\end{definition}

\begin{proposition}[Duality of Gauss phases]\label{prop:gauss_duality}
For nontrivial \(\chi\),
\[
\boxed{
\gamma(\chi,\psi)\gamma(\chi^{-1},\psi)
=
\chi(-1).
}
\]
\end{proposition}

\begin{proof}
The standard Gauss identity gives
\[
G(\chi,\psi)G(\chi^{-1},\psi)=\chi(-1)q.
\]
Divide by \(q\).
\end{proof}

\begin{standardfact}[Hasse and Davenport lifting]\label{fact:hasse_davenport}
Let \(k_r/k\) be the extension of degree \(r\). Lift \(\chi\) by the norm and \(\psi\) by the trace. Then the lifted Gauss sum satisfies the Hasse and Davenport product relation, with a sign determined by the chosen additive character normalization.
\end{standardfact}

\begin{figure}[t]
\centering
\includegraphics[width=.53\textwidth]{07_gauss_phase_circle.pdf}
\caption{Normalized finite Gauss sums lie on the unit circle.}
\label{fig:gausscircle}
\end{figure}

\begin{remark}\label{rem:gauss_boundary}
Finite Gauss phases reproduce the finite field part of tame and depth zero normalization. They do not by themselves give the full epsilon factor of a general wild pair.
\end{remark}

\section{Finite Heisenberg packets}\label{sec:heisenberg}

Let \(p\) be odd. Let \(X\) be an \(r\) dimensional vector space over \(\FF_p\), and put
\[
V=X\oplus X^*.
\]
Define the alternating form
\[
b\bigl((x,\xi),(y,\eta)\bigr)
=
\eta(x)-\xi(y).
\]
Let \(Z=\FF_p\), and define the Heisenberg group
\[
\frakH(V)=V\times Z
\]
with multiplication
\[
(v,z)(w,t)
=
\left(v+w,z+t+\frac12 b(v,w)\right).
\]
Fix the additive character
\[
\psi(z)=\e^{2\pi iz/p}.
\]

\begin{definition}[Schrödinger operators]\label{def:schrodinger}
On \(\calH=\CC^X\), define
\[
\bigl(\rho(x,\xi,z)f\bigr)(u)
=
\psi\left(z+\xi(u)+\frac12\xi(x)\right)f(u+x).
\]
\end{definition}

\begin{theorem}[Finite Schrödinger representation]\label{thm:schrodinger_rep}
The map \(\rho:\frakH(V)\to\GL(\calH)\) is a representation. It is irreducible, its center acts by \(\psi\), and
\[
\dim\calH=p^r=\sqrt{\abs V}.
\]
\end{theorem}

\begin{proof}
A direct calculation gives
\[
\rho(v,z)\rho(w,t)
=
\rho\left(v+w,z+t+\frac12b(v,w)\right).
\]
Translations by \(X\) and multiplications by characters from \(X^*\) generate the full matrix algebra on \(\CC^X\). Indeed, the character multipliers give all diagonal projectors by finite Fourier inversion, and translations move each coordinate line to every other coordinate line. The commutant is therefore scalar, so the representation is irreducible. The dimension is \(\abs X=p^r\).
\end{proof}

\begin{proposition}[Character support]\label{prop:heisenberg_trace}
The character of the Schrödinger representation satisfies
\[
\boxed{
\Tr\rho(v,z)
=
\begin{cases}
p^r\psi(z),&v=0,\\
0,&v\ne0.
\end{cases}
}
\]
\end{proposition}

\begin{proof}
If the translation component \(x\) is nonzero, the operator permutes the standard basis without a fixed point, so its trace is zero. If \(x=0\) but \(\xi\ne0\), the trace is
\[
\psi(z)\sum_{u\in X}\psi(\xi(u))=0.
\]
When \(x=\xi=0\), the operator is the scalar \(\psi(z)\).
\end{proof}

\begin{standardfact}[Finite Stone von Neumann theorem]\label{fact:stone_von_neumann}
For a finite Heisenberg group with nondegenerate commutator form and a fixed nontrivial central character, there is one irreducible representation up to isomorphism. Its dimension is the square root of the order of the symplectic quotient.
\end{standardfact}

The standard fact follows from the twisted group algebra argument in \cref{sec:twisted}. It is also the finite form used in the classical Weil representation literature \cite{Gerardin1977,Prasad2009}.

\subsection{Maximal isotropic selectors}

Let \(L\subseteq V\) be maximal isotropic. Choose a subgroup lift \(\widetilde L\subseteq\frakH(V)\) on which the central coordinate is zero.

\begin{proposition}[Isotropic averaging projector]\label{prop:isotropic_selector}
The operator
\[
P_L
=
\frac1{\abs L}
\sum_{\ell\in\widetilde L}
\rho(\ell)
\]
is an orthogonal projector. In the Schrödinger model with \(L=X\), it has rank one and selects the constant functions.
\end{proposition}

\begin{proof}
The lifted isotropic subgroup is abelian, so group averaging is an idempotent. Every operator in the average is unitary, and the subgroup is closed under inversion, so the average is self adjoint. For \(L=X\), the average is the translation average
\[
(P_Lf)(u)
=
\frac1{\abs X}
\sum_{x\in X}f(u+x),
\]
which is the constant function with value equal to the mean of \(f\). Its range is one dimensional.
\end{proof}

\begin{corollary}[Primitive line selector]\label{cor:primitive_selector}
Every maximal isotropic subgroup selects a one dimensional fixed line in the irreducible Heisenberg representation.
\end{corollary}

\begin{proof}
The projector rank equals its trace. The identity element contributes \(p^r\), while every nonidentity element of the lifted isotropic subgroup is noncentral and has trace zero by \cref{prop:heisenberg_trace}. Since \(\abs L=p^r\), we get
\[
\Tr P_L
=
\frac{p^r}{p^r}=1.
\]
Thus the projector has rank one.
\end{proof}

\begin{figure}[t]
\centering
\includegraphics[width=.9\textwidth]{08_heisenberg_geometry.pdf}
\caption{The finite Heisenberg packet and its maximal isotropic selector.}
\label{fig:heisenberg}
\end{figure}

\section{Twisted group algebras and commutator envelopes}\label{sec:twisted}

Let \(A\) be a finite abelian group, and let
\[
c:A\times A\to\CC^{\times}
\]
be a normalized two cocycle. The twisted group algebra \(\CC^c[A]\) has basis \((u_a)_{a\in A}\) and multiplication
\[
u_au_b=c(a,b)u_{a+b}.
\]
Define the commutator bicharacter
\[
b_c(a,b)=c(a,b)c(b,a)^{-1}.
\]

\begin{proposition}[Scalar choice invariance]\label{prop:cocycle_choice}
Replacing \(u_a\) by \(\lambda_au_a\), with \(\lambda_a\in\CC^{\times}\), changes \(c\) by a coboundary and leaves \(b_c\) unchanged.
\end{proposition}

\begin{proof}
The new cocycle is
\[
c'(a,b)
=
\lambda_a\lambda_b\lambda_{a+b}^{-1}c(a,b).
\]
The scalar coboundary factor is symmetric in \(a\) and \(b\), so it cancels from \(c'(a,b)c'(b,a)^{-1}\).
\end{proof}

\begin{proposition}[Alternating commutator form]\label{prop:commutator_form}
The map \(b_c\) is an alternating bicharacter. Its radical
\[
R=\{a\in A:b_c(a,x)=1\text{ for every }x\in A\}
\]
is a subgroup, and \(b_c\) descends to a nondegenerate alternating bicharacter on \(A/R\).
\end{proposition}

\begin{proof}
Associativity of the twisted product gives multiplicativity of scalar commutators in each variable. We have \(b_c(a,a)=1\) and \(b_c(b,a)=b_c(a,b)^{-1}\). The radical is the kernel of the homomorphism from \(A\) to its character group defined by \(a\mapsto b_c(a,\cdot)\). The induced pairing on the quotient is nondegenerate by construction.
\end{proof}

\begin{theorem}[Nondegenerate twisted group algebra]\label{thm:twisted_matrix}
Assume that \(b_c\) is nondegenerate. Then \(\abs A\) is a square, and
\[
\boxed{
\CC^c[A]
\cong
M_{\sqrt{\abs A}}(\CC).
}
\]
Consequently, \(\CC^c[A]\) has one irreducible module up to isomorphism, of dimension \(\sqrt{\abs A}\).
\end{theorem}

\begin{proof}
Let
\[
z=\sum_{a\in A}z_au_a
\]
be central. Commuting with \(u_x\) gives
\[
z_ab_c(x,a)=z_a
\]
for every \(x\). Nondegeneracy forces \(z_a=0\) for every nonzero \(a\). Thus the center is \(\CC\). To prove semisimplicity directly, choose a linear projection \(P\) onto a submodule of a finite dimensional \(\CC^c[A]\)-module and form
\[
\overline P=\frac{1}{\abs A}\sum_{a\in A}u_a P u_a^{-1}.
\]
Every summand maps into the submodule and is the identity there. Conjugation by any \(u_x\) permutes the summands: the cocycle scalars cancel between each operator and its inverse. Hence \(\overline P\) is an equivariant projection, and its kernel is an invariant complement. The algebra is therefore semisimple. Its center \(\CC\) permits only one simple factor. A finite dimensional central simple algebra over \(\CC\) is a full matrix algebra. Its dimension is \(\abs A\), so the matrix size is \(\sqrt{\abs A}\).
\end{proof}

\begin{corollary}[Radical free envelope]\label{cor:radical_free}
For an arbitrary scalar projective action of \(A\), the quotient
\[
\boxed{
A^{\mathrm{symp}}=A/\rad b_c
}
\]
is the maximal quotient detected by noncommuting projective intertwiners. If \(A\) is an elementary abelian \(p\) group, then \(A^{\mathrm{symp}}\) is a symplectic \(\FF_p\) space.
\end{corollary}

\begin{proof}
Every radical element commutes with all projective generators, so it contributes only central data. Quotienting by the radical leaves the nondegenerate commutator form. In the elementary abelian case, multiplicative values in the group of \(p\)th roots can be written as an alternating bilinear form over \(\FF_p\).
\end{proof}

\begin{figure}[t]
\centering
\includegraphics[width=.94\textwidth]{09_commutator_envelope.pdf}
\caption{Projective intertwiners determine a radical free symplectic envelope.}
\label{fig:commutator}
\end{figure}

\begin{remark}\label{rem:projective_scope}
The envelope depends only on the projective representation. It does not by itself prove that a given wild inertia representation carries a canonical finite quotient with the required projective action.
\end{remark}

\section{A restricted one break Galois envelope}\label{sec:onebreak}

We now isolate a finite Clifford setting in which the quotient projective action is defined on a multiplicity space. This avoids assigning a quotient action directly to arbitrary representatives of \(G/N\).

\begin{definition}[One break Clifford datum]\label{def:one_break_datum}
A one break Clifford datum consists of finite groups and representations
\[
N\triangleleft G,
\qquad
A=G/N,
\qquad
\Pi:G\to\GL(\calX),
\qquad
\rho:N\to\GL(W),
\]
with the following properties.

\begin{enumerate}
\item The quotient \(A\) is a finite elementary abelian \(p\) group.
\item The representations \(\Pi\) and \(\rho\) are irreducible.
\item The isomorphism class of \(\rho\) is fixed by \(G\).
\item The restriction of \(\Pi\) to \(N\) is \(\rho\) isotypic.
\item The quotient records one declared ramification break.
\end{enumerate}

The multiplicity space is
\[
M_\rho=\Hom_N(W,\calX).
\]
\end{definition}

The isotypic condition gives an \(N\) equivariant evaluation isomorphism
\[
W\otimes M_\rho
\longrightarrow
\calX.
\]
Choose a section \(s:A\to G\). For every \(a\in A\), choose an intertwiner \(I_a\) from \(\rho\) to its conjugate by \(s(a)\). Under the evaluation isomorphism, there is a unique operator \(T_a\) on \(M_\rho\) such that
\[
\Pi(s(a))=I_a\otimes T_a.
\]
The operators \(T_a\) form a projective representation of \(A\).

\begin{proposition}[Clifford multiplicity action]\label{prop:clifford_action}
The projective equivalence class of the action \((T_a)_{a\in A}\) is independent of the section, the intertwiners, and the chosen isotypic evaluation isomorphism. The projective action is irreducible.
\end{proposition}

\begin{proof}
Changing an intertwiner \(I_a\) by a scalar changes \(T_a\) by the inverse scalar. Changing the isotypic evaluation isomorphism conjugates all \(T_a\) by one element of \(\GL(M_\rho)\). Changing the section multiplies \(s(a)\) by an element of \(N\). That factor is absorbed by the \(\rho\) action on \(W\), leaving a scalar change on the multiplicity operator by Schur's lemma. Thus all choices give projectively equivalent actions.

If a nonzero proper subspace of \(M_\rho\) were stable under all projective operators \(T_a\), then its tensor product with \(W\) would give a nonzero proper \(G\) stable subspace of \(\calX\). This contradicts irreducibility of \(\Pi\).
\end{proof}

\begin{definition}[Clifford cocycle class]\label{def:clifford_class}
Choose scalars \(c(a,b)\) such that
\[
T_aT_b=c(a,b)T_{a+b}.
\]
The cohomology class
\[
[c_\rho]\in H^2(A,\CC^{\times})
\]
is the Clifford cocycle class of the datum.
\end{definition}

\begin{proposition}[Intrinsic commutator pairing]\label{prop:one_break_form}
The pairing
\[
b_\rho(a,b)
=
c(a,b)c(b,a)^{-1}
\]
is independent of every choice in \cref{def:one_break_datum}. It is an alternating bicharacter
\[
b_\rho:A\times A\to\mu_p.
\]
\end{proposition}

\begin{proof}
The projective equivalence class is intrinsic by \cref{prop:clifford_action}. Conjugation does not change scalar commutators, and scalar changes alter the cocycle by a coboundary. The commutator pairing is therefore unchanged by \cref{prop:cocycle_choice}. Its alternating and bilinear properties follow from \cref{prop:commutator_form}. Since every element of \(A\) has order \(p\), every value has order dividing \(p\).
\end{proof}

\begin{definition}[One break symplectic envelope]\label{def:one_break_envelope}
The one break symplectic envelope is
\[
\boxed{
V_\rho=A/\rad b_\rho.
}
\]
\end{definition}

\begin{theorem}[Nondegenerate one break packet]\label{thm:one_break_matrix}
If \(b_\rho\) is nondegenerate, then
\[
\CC^{c_\rho}[A]
\cong
\End(M_\rho),
\qquad
\dim M_\rho=\sqrt{\abs A}.
\]
\end{theorem}

\begin{proof}
The projective action gives a nonzero representation of the twisted group algebra on \(M_\rho\). The algebra is a full matrix algebra by \cref{thm:twisted_matrix}. Its unique irreducible module has dimension \(\sqrt{\abs A}\). The projective multiplicity action is irreducible by \cref{prop:clifford_action}, so it is that module.
\end{proof}

\begin{proposition}[Isomorphism invariance]\label{prop:one_break_invariance}
The cocycle class and the equivariant isometry class of \(V_\rho\) depend only on the isomorphism class of the one break Clifford datum.
\end{proposition}

\begin{proof}
An isomorphism of data transports the restriction, the multiplicity space, and its projective quotient action. It therefore transports the cocycle class, the commutator pairing, its radical, and the quotient.
\end{proof}

\begin{proposition}[Tame stabilizer action]\label{prop:tame_stabilizer}
Suppose a finite group \(C\) acts by automorphisms of the one break Clifford datum. Then \(C\) acts on \(V_\rho\) by symplectic automorphisms.
\end{proposition}

\begin{proof}
The automorphism action transports the intrinsic projective multiplicity representation. It therefore preserves the scalar commutator pairing and its radical. The induced action on the quotient is symplectic.
\end{proof}

\begin{figure}[t]
\centering
\includegraphics[width=.96\textwidth]{10_one_break_envelope.pdf}
\caption{The restricted one break construction is made on the Clifford multiplicity space.}
\label{fig:onebreak}
\end{figure}

\begin{corollary}[Heisenberg representation attached to one break]\label{cor:one_break_heisenberg}
Fix a nontrivial central character. The envelope \(V_\rho\) determines one irreducible finite Heisenberg representation up to isomorphism.
\end{corollary}

\begin{proof}
Apply the finite Stone von Neumann theorem to the nondegenerate symplectic quotient.
\end{proof}

\begin{remark}\label{rem:clifford_scope}
The group cohomology class records the extension obstruction for the multiplicity action. General background on projective representations and group cohomology can be found in \cite{Karpilovsky1985,Brown1982}. The theorem remains restricted because the finite quotient and the one break isotypic datum are assumed rather than constructed from arbitrary wild inertia.
\end{remark}

\section{The type side symplectic quotient}\label{sec:typeside}

Let \(\theta\) be a maximal simple character for \(\GL_n(K)\), with compact open subgroups
\[
H_\theta^1\subseteq J_\theta^1\subseteq J_\theta.
\]
The quotient
\[
V_\theta=J_\theta^1/H_\theta^1
\]
is regarded as a finite vector space over \(\FF_p\), where \(p\) is the residue characteristic. Fix a nontrivial additive character \(\psi_0:\FF_p\to\mu_p\subset\CC^\times\), where \(\mu_p\) is the group of \(p\)-th roots of unity. The simple-character commutator takes values in \(\mu_p\). Since \(\psi_0\) is injective, the equation
\[
\psi_0\bigl(h_\theta(\bar x,\bar y)\bigr)
=
\theta([x,y])
\]
uniquely defines the \(\FF_p\)-valued commutator pairing \(h_\theta\). We use the same identification of \(\mu_p\) with the additive group of \(\FF_p\) on the Galois side when comparing the two symplectic spaces.

\begin{standardfact}[Type side Heisenberg quotient]\label{fact:type_heisenberg}
The pairing \(h_\theta\) is nondegenerate and alternating. There is one irreducible representation \(\eta_\theta\) of \(J_\theta^1\) containing \(\theta\), and its dimension is
\[
\sqrt{\abs{J_\theta^1/H_\theta^1}}.
\]
\end{standardfact}

This is part of the simple character and Heisenberg representation theory used in the construction of types \cite{BushnellKutzko1993}. The extension of \(\eta_\theta\) from \(J_\theta^1\) to \(J_\theta\) is not unique before a normalization is chosen.

\begin{standardfact}[Canonical beta extension]\label{fact:canonical_beta}
The canonical beta extension is obtained from a primary extension together with finite symplectic and tame orientation corrections. The construction is functorial in the declared simple character data and removes the remaining extension ambiguity.
\end{standardfact}

The normalization and its relation to finite Weil representations are developed in \cite{Blondel2012,Dotto2022}. This chapter uses this theorem as type side input. It does not rederive simple character theory.

\begin{proposition}[Type side projective commutator]\label{prop:type_commutator}
The projective action of \(J_\theta/J_\theta^1\) on the Heisenberg representation preserves the commutator form on \(V_\theta\). Any two choices of projective lifts give the same alternating form.
\end{proposition}

\begin{proof}
Conjugation by \(J_\theta\) preserves \(\theta\), so it preserves the commutator pairing defined from \(\theta([x,y])\). Two choices of projective lifts differ by scalars. The scalar commutator is therefore unchanged by \cref{prop:cocycle_choice}.
\end{proof}

\section{The restricted Galois to type bridge}\label{sec:bridge}

The full bridge would construct an intrinsic Galois envelope for every wild atom and identify it canonically with \(V_\theta\). We prove the comparison after the required symplectic identification has been supplied.

\begin{definition}[Bridge datum]\label{def:bridge_datum}
A restricted bridge datum consists of

\begin{enumerate}
\item a one break Galois envelope \((V_\rho,b_\rho)\)
\item a type side quotient \((V_\theta,h_\theta)\)
\item an equivariant symplectic isometry
\[
\iota:V_\rho\longrightarrow V_\theta
\]
for a finite cyclic stabilizer \(C=\langle c\rangle\)
\item the same nontrivial central character on the two Heisenberg groups
\item when \(p=2\), compatible quadratic refinements and an isomorphism of the resulting Heisenberg central extensions lifting \(\iota\).
\end{enumerate}
\end{definition}

\begin{theorem}[Restricted Heisenberg bridge]\label{thm:restricted_bridge}
A restricted bridge datum determines an isomorphism between the two Heisenberg representations, unique up to a nonzero scalar. More precisely, there is an invertible map
\[
U:\calH_\rho\longrightarrow\calH_\theta
\]
such that
\[
U\omega_\rho(h)
=
\omega_\theta(\iota_H(h))U
\]
for every Heisenberg element \(h\), and every other such map is a scalar multiple of \(U\).
\end{theorem}

\begin{proof}
When \(p\) is odd, the symplectic isometry and the common central character induce the standard isomorphism \(\iota_H\) of the two Heisenberg groups. When \(p=2\), the required lift of \(\iota\) is part of the bridge datum because the alternating form alone does not determine the central extension. Pulling the type side representation back through \(\iota_H\) gives an irreducible Heisenberg representation with the same central character as the Galois side representation. The finite Stone von Neumann theorem gives an isomorphism. Schur's lemma gives uniqueness up to scalar.
\end{proof}

The stabilizer actions need not agree after choosing \(U\). Their discrepancy is scalar because both actions implement the same automorphisms of the irreducible Heisenberg representation.

\begin{proposition}[Cyclic discrepancy character]\label{prop:cyclic_discrepancy}
Assume that the two projective stabilizer actions have been linearized to honest representations of the cyclic group \(C\). Then there is a character
\[
\delta:C\to\CC^{\times}
\]
such that
\[
U\Omega_\rho(g)
=
\delta(g)\Omega_\theta(g)U
\]
for every \(g\in C\).
\end{proposition}

\begin{proof}
For each \(g\), the operator
\[
\Omega_\theta(g)^{-1}U\Omega_\rho(g)U^{-1}
\]
commutes with the type side Heisenberg representation. It is therefore scalar. Since both stabilizer actions are honest representations, multiplication of group elements shows that the scalars form a character.
\end{proof}

\begin{lemma}[Linearization over a cyclic stabilizer]\label{lem:cyclic_linearization}
Every finite dimensional projective representation of a finite cyclic group over \(\CC\) is projectively equivalent to an honest representation.
\end{lemma}

\begin{proof}
Let \(C=\langle c\rangle\) have order \(m\), and let \(T\) be a lift of the projective generator. Since \(c^m=1\), the operator \(T^m\) is scalar. Choose \(\lambda\in\CC^{\times}\) with \(\lambda^mT^m=I\). Replacing \(T\) by \(\lambda T\) gives an honest cyclic representation.
\end{proof}

\begin{theorem}[One phase closes the cyclic bridge]\label{thm:one_phase_bridge}
In a restricted bridge datum, choose linearizations of the cyclic stabilizer actions. If the two actions have the same scalar phase on the generator \(c\), then the Heisenberg intertwiner can be chosen to intertwine the full semidirect product action.
\end{theorem}

\begin{proof}
By \cref{prop:cyclic_discrepancy}, the discrepancy is a character \(\delta\). Equality of the generator phase gives \(\delta(c)=1\). Since \(c\) generates \(C\), the character is trivial. The chosen map \(U\) then intertwines every stabilizer element as well as the Heisenberg group.
\end{proof}

\begin{figure}[t]
\centering
\includegraphics[width=.94\textwidth]{11_restricted_bridge.pdf}
\caption{An equivariant symplectic isometry reduces the finite comparison to one scalar cyclic phase.}
\label{fig:restrictedbridge}
\end{figure}

\begin{corollary}[Bridge uniqueness]\label{cor:bridge_uniqueness}
When the semidirect product representation is irreducible, the completed bridge is unique up to scalar.
\end{corollary}

\begin{proof}
Any two completed bridges differ by an endomorphism commuting with the semidirect product representation. Schur's lemma applies.
\end{proof}

\begin{remark}\label{rem:bridge_scope}
The theorem does not construct \(\iota\). It proves that, once an equivariant symplectic isometry exists, the remaining finite comparison is one dimensional for a cyclic stabilizer.
\end{remark}

\section{Odd residue characteristic phases}\label{sec:oddphase}

Let \(p\) be odd, let \((V,b)\) be a finite symplectic \(\FF_p\) space of dimension \(2r\), and let
\[
g\in\Sp(V)
\]
have no fixed vector. Fix one of the standard finite Weil normalizations used in the literature \cite{Gerardin1977,Neuhauser2002,GurevichHadani2007}.

\begin{standardfact}[Fixed point free Weil trace]\label{fact:weil_trace}
Under the declared normalization,
\[
\boxed{
\Tr\omega(g)
=
\left(\frac{-1}{p}\right)^r
\left(\frac{\det(g-I)}{p}\right).
}
\]
Here \((\frac{\cdot}{p})\) denotes the quadratic character of \(\FF_p^{\times}\).
\end{standardfact}

\begin{definition}[Odd finite Weil phase]\label{def:odd_phase}
For fixed point free \(g\), define
\[
\Phi_p(g,V)
=
\left(\frac{-1}{p}\right)^r
\left(\frac{\det(g-I)}{p}\right).
\]
\end{definition}

\begin{proposition}[Symplectic invariance]\label{prop:phase_invariance}
If
\[
\iota:(V,b)\longrightarrow(V',b')
\]
is a symplectic isometry and \(g'=\iota g\iota^{-1}\), then
\[
\Phi_p(g,V)=\Phi_p(g',V').
\]
\end{proposition}

\begin{proof}
The dimensions agree, and conjugate matrices have the same determinant of \(g-I\).
\end{proof}

\begin{proposition}[Orthogonal multiplicativity]\label{prop:odd_multiplicativity}
Suppose
\[
V=V_1\perpbig\cdots\perpbig V_s
\]
and \(g\) preserves every summand. If no \(g|_{V_j}\) has a fixed vector, then
\[
\boxed{
\Phi_p(g,V)
=
\prod_{j=1}^{s}
\Phi_p(g|_{V_j},V_j).
}
\]
\end{proposition}

\begin{proof}
The half dimensions add, and
\[
\det(g-I\mid V)
=
\prod_j\det(g-I\mid V_j).
\]
The quadratic character is multiplicative.
\end{proof}

\begin{corollary}[Characteristic polynomial certificate]\label{cor:characteristic_phase}
Let
\[
P_g(X)=\det(XI-g\mid V).
\]
Then
\[
\boxed{
\Phi_p(g,V)
=
\left(\frac{-1}{p}\right)^r
\left(\frac{P_g(1)}{p}\right).
}
\]
\end{corollary}

\begin{proof}
Since \(\dim V\) is even,
\[
\det(g-I)=\det(I-g)=P_g(1).
\]
\end{proof}


\begin{proposition}[Fixed and moving decomposition]\label{prop:fixed_moving}
Assume that \(g\in\Sp(V)\) has order prime to \(p\). Then
\[
V=\ker(g-I)\perpbig\im(g-I).
\]
The restriction of the symplectic form to \(\im(g-I)\) is nondegenerate, and \(g-I\) is invertible on that moving summand.
\end{proposition}

\begin{proof}
The order hypothesis makes \(g\) semisimple, so
\[
V=\ker(g-I)\oplus\im(g-I).
\]
If \(x\in\ker(g-I)\) and \(y=(g-I)z\), then
\[
b(x,y)=b(x,gz)-b(x,z)=b(g^{-1}x,z)-b(x,z)=0.
\]
Thus the two summands are orthogonal. Nondegeneracy on each summand follows from nondegeneracy on \(V\). The last statement is immediate from the direct sum decomposition.
\end{proof}

\begin{definition}[Moving phase]\label{def:moving_phase}
Under the hypotheses of \cref{prop:fixed_moving}, define
\[
\Phi_p^{\mathrm{mov}}(g,V)
=
\Phi_p\left(g|_{\im(g-I)},\im(g-I)\right).
\]
\end{definition}

\begin{proposition}[Hyperbolic phase formula]\label{prop:hyperbolic_phase}
Let \(U\) be a finite dimensional \(\FF_p\) space, let \(a\in\GL(U)\) have no eigenvalue one, and let \(g\) act on
\[
H(U)=U\oplus U^*
\]
by
\[
g=(a,(a^{-1})^*).
\]
Then
\[
\boxed{
\Phi_p(g,H(U))
=
\left(\frac{\det a}{p}\right).
}
\]
\end{proposition}

\begin{proof}
Put \(d=\dim U\). We have
\[
\det\bigl((a^{-1})^*-I\bigr)
=
(-1)^d(\det a)^{-1}\det(a-I).
\]
Hence
\[
\det(g-I)
=
(-1)^d(\det a)^{-1}\det(a-I)^2.
\]
Multiplying its quadratic character by the normalization factor \((\frac{-1}{p})^d\) cancels the sign and the square. Since a quadratic character equals its inverse, the result is \((\frac{\det a}{p})\).
\end{proof}

\begin{corollary}[Phase only bridge test]\label{cor:phase_only_test}
Two odd characteristic moving packets have the same normalized phase whenever their moving half dimensions agree and the determinants of \(g-I\) lie in the same square class.
\end{corollary}

\begin{proof}
Apply the definition of \(\Phi_p\).
\end{proof}

\begin{figure}[t]
\centering
\includegraphics[width=.88\textwidth]{12_odd_phase_blocks.pdf}
\caption{Odd finite Weil phases multiply across orthogonal atomic blocks.}
\label{fig:oddphase}
\end{figure}

\begin{remark}\label{rem:fixed_space}
When \(g\) has fixed vectors, one first separates the fixed part and applies the formula to the moving quotient. Additional normalization data may enter the full trace on the fixed part.
\end{remark}

\section{Dyadic quadratic refinements}\label{sec:dyadic}

Let \((V,b)\) be a symplectic vector space over \(\FF_2\). The alternating form alone is not enough to determine the finite phase.

\begin{definition}[Quadratic refinement]\label{def:quadratic_refinement}
A quadratic refinement of \(b\) is a map
\[
q:V\to\FF_2
\]
satisfying
\[
q(x+y)-q(x)-q(y)=b(x,y).
\]
\end{definition}

\begin{proposition}[Affine space of refinements]\label{prop:refinement_affine}
If one quadratic refinement exists, then the set of all quadratic refinements is an affine space over \(V^*\).
\end{proposition}

\begin{proof}
The difference of two refinements has zero polar form, so it is additive and hence linear over \(\FF_2\). Conversely, adding a linear form does not change the polar form.
\end{proof}

\begin{definition}[Arf invariant]\label{def:arf}
Choose a symplectic basis \((e_i,f_i)_{i=1}^{r}\). Define
\[
\Arf(q)
=
\sum_{i=1}^{r}q(e_i)q(f_i)
\in\FF_2.
\]
\end{definition}

\begin{standardfact}[Arf invariance and classification]\label{fact:arf_classification}
The value in \cref{def:arf} is independent of the symplectic basis. Two nondegenerate quadratic refinements of one finite symplectic space are isometric exactly when their Arf invariants agree.
\end{standardfact}

This is the classical classification of nonsingular quadratic forms over \(\FF_2\) \cite{MilnorHusemoller1973}.

\begin{theorem}[Dyadic Gauss sign]\label{thm:arf_gauss}
Let \(\dim V=2r\), and let \(q\) be nondegenerate. Then
\[
\boxed{
\sum_{x\in V}(-1)^{q(x)}
=
(-1)^{\Arf(q)}2^r.
}
\]
\end{theorem}

\begin{proof}
Orthogonally decompose \((V,q)\) into two dimensional nonsingular planes. On a plane with Arf invariant zero, the four values of \((-1)^q\) sum to \(2\). On a plane with Arf invariant one, they sum to \(-2\). The total sum is the product of the plane sums. The parity of the number of anisotropic planes is the Arf invariant.
\end{proof}

\begin{corollary}[Orthogonal dyadic multiplicativity]\label{cor:arf_multiplicative}
For an orthogonal sum
\[
(V,q)=(V_1,q_1)\perpbig(V_2,q_2),
\]
we have
\[
\Arf(q)=\Arf(q_1)+\Arf(q_2)
\]
and the normalized Gauss sign is the product of the two normalized signs.
\end{corollary}

\begin{proof}
Take symplectic bases of the two summands and concatenate them. The Gauss sum factors as a product over the Cartesian decomposition.
\end{proof}

\begin{proposition}[The form alone does not determine the phase]\label{prop:dyadic_need_q}
On every nonzero symplectic \(\FF_2\) space, there are quadratic refinements of both Arf invariants. Their polar forms agree, but their normalized Gauss signs are opposite.
\end{proposition}

\begin{proof}
On a symplectic plane, the standard hyperbolic refinement has Arf invariant zero. Adding a suitable nonzero linear form produces a refinement with Arf invariant one. Orthogonally adjoining a fixed refinement on the complement gives the two choices in every positive even dimension. The signs differ by \cref{thm:arf_gauss}.
\end{proof}


\begin{proposition}[Translation of quadratic refinements]\label{prop:quadratic_translation}
For \(v\in V\), put
\[
q_v(x)=q(x)+b(v,x).
\]
Then
\[
\boxed{
\Arf(q_v)=\Arf(q)+q(v).
}
\]
\end{proposition}

\begin{proof}
Using the quadratic identity,
\[
q_v(x)=q(x+v)+q(v).
\]
Therefore
\[
\sum_x(-1)^{q_v(x)}
=
(-1)^{q(v)}
\sum_x(-1)^{q(x)}.
\]
Compare the signs in \cref{thm:arf_gauss}.
\end{proof}

\begin{theorem}[Number of dyadic refinements]\label{thm:quadratic_counts}
Let \(\dim V=2r\). The number of quadratic refinements with Arf invariant zero is
\[
\boxed{
2^{r-1}(2^r+1),
}
\]
and the number with Arf invariant one is
\[
\boxed{
2^{r-1}(2^r-1).
}
\]
\end{theorem}

\begin{proof}
Every refinement is \(q_v\) for a unique \(v\in V\), because \(b\) identifies \(V\) with \(V^*\). By \cref{prop:quadratic_translation}, the Arf invariant changes according to \(q(v)\). Let \(Z\) and \(O\) be the numbers of zeros and ones of \(q\). Then
\[
Z+O=2^{2r}
\]
and \cref{thm:arf_gauss} gives
\[
Z-O=(-1)^{\Arf(q)}2^r.
\]
Solving gives the two displayed numbers, with their roles exchanged when the starting refinement has Arf invariant one. Since translation runs through every refinement, the final counts are independent of the starting choice.
\end{proof}

\begin{figure}[t]
\centering
\includegraphics[width=.92\textwidth]{13_dyadic_quadratic_refinement.pdf}
\caption{At residue characteristic two, the quadratic refinement and its Arf invariant carry phase information not contained in the alternating form.}
\label{fig:dyadic}
\end{figure}

\section{Matrix holonomy and finite gluing}\label{sec:holonomy}

Let \(\Gamma\) be a connected finite graph. At every vertex \(v\), place finite dimensional complex vector spaces \(V_v\) and \(W_v\) of the same dimension. For every oriented edge \(e:u\to v\), let
\[
A_e:V_u\longrightarrow V_v,
\qquad
B_e:W_u\longrightarrow W_v
\]
be invertible. The reversed edge carries the inverse map.

\begin{definition}[Path transport]\label{def:path_transport}
For an oriented path
\[
\gamma=(v_0\to v_1\to\cdots\to v_m),
\]
define
\[
A_\gamma=A_{v_m v_{m-1}}\cdots A_{v_1v_0},
\qquad
B_\gamma=B_{v_m v_{m-1}}\cdots B_{v_1v_0}.
\]
For a loop based at \(r\), these are the based holonomies.
\end{definition}

\begin{theorem}[Matrix holonomy criterion]\label{thm:matrix_holonomy}
Fix a root \(r\) and an isomorphism
\[
U_r:V_r\longrightarrow W_r.
\]
The following conditions are equivalent.

\begin{enumerate}
\item There are isomorphisms \(U_v:V_v\to W_v\), with the declared \(U_r\), such that
\[
U_vA_e=B_eU_u
\]
for every edge \(e:u\to v\).
\item For one fundamental cycle basis based at \(r\),
\[
U_rA_\gamma=B_\gamma U_r.
\]
\item The same identity holds for every based loop \(\gamma\).
\end{enumerate}
\end{theorem}

\begin{proof}
If the vertex maps exist, multiply the edge equations around a loop. Every intermediate vertex map cancels, giving condition three. Condition three implies condition two.

Conversely, choose a spanning tree. Starting from \(U_r\), define \(U_v\) by transport along the unique tree path from \(r\) to \(v\). The edge equations hold on every tree edge by construction. Every non tree edge closes one fundamental cycle. The edge equation on that edge is equivalent to the root holonomy identity for its fundamental cycle. Condition two therefore gives all remaining edge equations.
\end{proof}

\begin{corollary}[Unbased gauge classification]\label{cor:unbased_holonomy}
There exists some family of vertex isomorphisms if and only if the two holonomy representations
\[
\pi_1(\Gamma,r)\longrightarrow\GL(V_r),
\qquad
\pi_1(\Gamma,r)\longrightarrow\GL(W_r)
\]
are conjugate.
\end{corollary}

\begin{proof}
Apply \cref{thm:matrix_holonomy} and allow \(U_r\) to vary.
\end{proof}

\begin{corollary}[Scalar and projective gluing]\label{cor:scalar_holonomy}
If every based loop acts by a scalar, finite gluing is equivalent to equality of the scalar holonomies on a fundamental cycle basis. The number of independent tests is
\[
\boxed{
\abs E-\abs V+1.
}
\]
\end{corollary}

\begin{proof}
A connected graph has a spanning tree with \(\abs V-1\) edges. Every remaining edge gives one fundamental cycle. In the scalar case, conjugacy has no effect.
\end{proof}

\begin{proposition}[Projective cocycle form]\label{prop:projective_holonomy}
Suppose the edge maps are defined only up to scalar, in the setting where all based holonomies of both systems are scalar. Choose representatives. The ratio of the two scalar holonomies on each loop is unchanged by vertex rescaling. It defines a character of the first homology group when the scalar transport is abelian.
\end{proposition}

\begin{proof}
A vertex rescaling appears once with positive exponent and once with negative exponent around a closed path, so it cancels. Multiplication of loops multiplies the scalar ratios. In the abelian case, the ratio factors through the abelianization of the fundamental group, which is the first homology group.
\end{proof}


\begin{proposition}[Exact edge residual from cycle defect]\label{prop:holonomy_residual}
Fix a spanning tree and construct the vertex maps by tree transport from \(U_r\). Let \(e:u\to v\) be a non tree edge. Let \(p_u\) and \(p_v\) be the tree paths from the root to \(u\) and \(v\). Define the fundamental cycle defect
\[
D_e
=
U_rA_{p_v}^{-1}A_eA_{p_u}
-
B_{p_v}^{-1}B_eB_{p_u}U_r.
\]
Then the edge residual is exactly
\[
\boxed{
U_vA_e-B_eU_u
=
B_{p_v}D_eA_{p_u}^{-1}.
}
\]
\end{proposition}

\begin{proof}
Substitute
\[
U_u=B_{p_u}U_rA_{p_u}^{-1},
\qquad
U_v=B_{p_v}U_rA_{p_v}^{-1}
\]
into the left side and factor \(B_{p_v}\) on the left and \(A_{p_u}^{-1}\) on the right.
\end{proof}

\begin{corollary}[Quantitative finite gluing]\label{cor:quantitative_holonomy}
With the notation of \cref{prop:holonomy_residual},
\[
\norm{U_vA_e-B_eU_u}
\leq
\norm{B_{p_v}}\norm{D_e}\norm{A_{p_u}^{-1}}.
\]
Thus small fundamental cycle defects give small non tree edge residuals whenever the tree transports are uniformly conditioned.
\end{corollary}

\begin{proof}
Take operator norms in the exact identity.
\end{proof}

\begin{figure}[t]
\centering
\includegraphics[width=.72\textwidth]{14_matrix_holonomy.pdf}
\caption{A spanning tree fixes the vertex gauges. The non tree edges test the remaining cycle holonomies.}
\label{fig:holonomy}
\end{figure}

\begin{remark}\label{rem:holonomy_boundary}
The theorem is an exact finite gluing result. It assumes that the packet spaces and edge transports have already been constructed on both sides.
\end{remark}

\section{Breakwise assembly}\label{sec:breakwise}

Wild inertia carries an upper ramification filtration. In favorable settings, a wild atom separates into finitely many break pieces. We state the finite algebra needed once such pieces are intrinsic.

\begin{definition}[Orthogonal break decomposition]\label{def:break_decomposition}
An orthogonal break decomposition is a symplectic direct sum
\[
V=\perpbig_{s\in S}V_s
\]
that is preserved by the declared tame stabilizer. Each \(V_s\) is the finite symplectic packet assigned to one break.
\end{definition}

\begin{proposition}[Heisenberg tensor factorization]\label{prop:heisenberg_tensor}
Fix compatible central characters. The irreducible Heisenberg representation of an orthogonal sum is the tensor product of the irreducible Heisenberg representations of its summands:
\[
\boxed{
\omega_V
\cong
\bigotimes_{s\in S}\omega_{V_s}.
}
\]
\end{proposition}

\begin{proof}
The Heisenberg group of an orthogonal sum is the central product of the Heisenberg groups of the summands. The tensor product representation has the declared central character and has dimension
\[
\prod_s\sqrt{\abs{V_s}}
=
\sqrt{\abs V}.
\]
It is irreducible by the finite Stone von Neumann theorem, so it is the required representation.
\end{proof}

\begin{corollary}[Breakwise phase product]\label{cor:break_phase}
In odd residue characteristic, let \(g\in\Sp(V)\) preserve every break packet and suppose its restriction to \(M=\im(g-I)\) is fixed point free. Put \(g_s=g|_{V_s}\) and \(M_s=\im(g_s-I)\). The moving restrictions are symplectic, and
\[
\boxed{
\Phi_p(g|_M,M)
=
\prod_{s\in S}\Phi_p(g_s|_{M_s},M_s).
}
\]
In residue characteristic two, the analogous normalized Gauss sign is the product of the signs determined by the breakwise quadratic refinements.
\end{corollary}

\begin{proof}
The invariant orthogonal decomposition gives \(M=\mathop{\perp}\limits_{s\in S}M_s\). Symplectic invariance gives \(M^\perp=\ker(g-I)\). Fixed point freeness on \(M\) therefore makes the restricted form nondegenerate on \(M\) and on each \(M_s\). Each restriction of \(g_s-I\) to \(M_s\) is invertible, so \cref{prop:odd_multiplicativity} applies to the moving spaces. A zero moving space contributes one. In characteristic two, apply \cref{cor:arf_multiplicative}.
\end{proof}

\begin{theorem}[Breakwise restricted bridge]\label{thm:breakwise_bridge}
Suppose every break has a restricted bridge datum, the symplectic isometries are compatible with the stabilizer action, and the scalar phases agree on a generating set of the stabilizer. Then the tensor product of the breakwise Heisenberg intertwiners gives a bridge for the full finite envelope. Any remaining gluing obstruction is exactly the matrix holonomy obstruction of \cref{thm:matrix_holonomy}.
\end{theorem}

\begin{proof}
Tensor the breakwise Heisenberg intertwiners from \cref{thm:restricted_bridge}. Phase agreement makes the tensor map equivariant on the chosen stabilizer generators. The breakwise direct sum and tensor factorization make the local comparisons compatible. If the packets are arranged over a transport graph, the only remaining condition is consistency of the local maps around cycles, which is exactly the holonomy criterion.
\end{proof}

\begin{figure}[t]
\centering
\includegraphics[width=.9\textwidth]{15_breakwise_decomposition.pdf}
\caption{Orthogonal break packets produce tensor factors of the finite Heisenberg representation.}
\label{fig:breakwise}
\end{figure}

\section{Local factor tests and the remaining bridge}\label{sec:boundary}

Local class field theory identifies one dimensional Weil characters with characters of \(K^{\times}\) \cite{Serre1967,Tate1979}. The tame induction theorem extends this explicit picture to higher dimensional tame atoms. The local polynomial and conductor calculations in \cref{prop:tame_local_polynomial,prop:tame_conductor} therefore give explicit finite invariants for the declared tame packets. Finite Gauss sums give the residue field phase. The full local factor, however, also includes ramification depth, Haar measure normalization, the additive character conductor, and wild interaction data.

\Needspace{13\baselineskip}
\begin{proposition}[Finite tame comparison package]\label{prop:tame_package}
For a tame induced atom, the following data are determined by its packet coordinates:

\begin{enumerate}
\item the dimension and inertia orbit
\item the Frobenius cycle polynomial
\item the Artin conductor
\item the local polynomial
\item the tensor refinement into interaction packets
\item the finite Gauss sum of each declared residue character, with unit-modulus phase normalization when that character is nontrivial.
\end{enumerate}
\end{proposition}

\begin{proof}
The dimension and orbit follow from \cref{thm:tame_induction}. The cycle polynomial is \cref{prop:frobenius_cycle}. For rank greater than one, the conductor is \cref{prop:tame_conductor}. In rank one, the Swan conductor is zero, and the Artin conductor is zero for trivial inertia and one otherwise, by the formula \(\operatorname{Art}=\dim V-\dim V^{I_K}+\operatorname{Sw}\). The local polynomial is \cref{prop:tame_local_polynomial}. Tensor refinement is \cref{thm:gcd_lcm}. The Gauss sum is its defining finite sum; for a nontrivial multiplicative character its unit phase is \cref{def:gauss_phase}. For the trivial character and a nontrivial additive character \(\psi\), the sum is \(G(1,\psi)=-1\), because the sum of \(\psi\) over the whole field is zero.
\end{proof}

\subsection{Arithmetic realization and the conditional comparison}

Fix this chapter's reciprocity, additive-character, local-factor, and phase
normalizations. Let \(\mathcal A\) be the declared class of irreducible wild
Weil atoms under consideration. For \(\rho\in\mathcal A\), the supplied
assembled finite datum \(d_\rho\) consists of the following inputs,
historically denoted H1--H4:
\begin{enumerate}
\item an intrinsic finite break decomposition on the Galois side;
\item the declared equivariant symplectic bridges to the supplied breakwise
type-side quotients;
\item compatible odd or dyadic central-extension and phase data on every
break, including every prescribed finite marking;
\item a compatible transport system with vanishing matrix holonomy
obstruction.
\end{enumerate}
These are finite comparison inputs. Supplying separate breakwise target
types does not itself supply one common arithmetic candidate. If a common
candidate is already supplied, that is a stronger restricted hypothesis.

An \emph{arithmetic realization witness}
\(\mathcal R(d_\rho,\pi)\) consists of actual arithmetic inducing data
for an admissible supercuspidal \(\pi\), together with explicit
comparison and recovery maps carrying the prescribed assembled finite
packet. The maps preserve its declared stabilizer action, all character values
explicitly included in the supplied datum, central-extension and phase
data, and compatible transports.
Recovery returns that same \(\pi\), with its fixed central character and
noncompact completion. An enlarged actual inducing carrier is permitted;
the whole packet is not required to be the original maximal-simple-type
quotient. Any stronger native-marking or numerical-layer requirement is
part of the declared datum and must be verified separately.

\begingroup
\renewcommand{\theconjecture}{66}
\begin{conjecture}[Arithmetic packet realization]
\label{conj:arithmetic_realization}\label{thm:conditional_atomic}
For every \(\rho\in\mathcal A\) with its supplied H1--H4 assembled
finite datum \(d_\rho\), there exists an actual admissible supercuspidal
representation \(\pi\) of \(\GL_{\dim\rho}(K)\) and an arithmetic
realization witness \(\mathcal R(d_\rho,\pi)\).
\end{conjecture}
\endgroup

\noindent This is the realization problem isolated from historical
Theorem 16.2 and original proof obligation 66 (row 95 of the historical
all-records register). It asks for one witness for each
admitted source, rather than for success of every proposed host or every
noncompact completion. No proof of the unrestricted statement is claimed.

\begin{proposition}[Comparison of one realized candidate]
\label{prop:realized_candidate_comparison}
Let \(\rho\in\mathcal A\), and let \(\pi\) be one actual admissible
supercuspidal representation of \(\GL_{\dim\rho}(K)\) with
\(\mathcal R(d_\rho,\pi)\). Suppose H5 holds for this same \(\pi\):
its required pair local factors equal those prescribed for \(\rho\)
on the stipulated characterizing family of twists, with the fixed
normalizations. Assume the stated local-factor characterization applies
to that family and class. Then \(\pi\) is isomorphic to the representation
assigned to \(\rho\) by the established local Langlands correspondence.
\end{proposition}

\begin{proof}
H5 identifies the tested local factors for this actual candidate. The
applicable local-factor characterization identifies its isomorphism
class with the established correspondence \cite{Henniart1993}. The
realization witness ensures that it is this same candidate which carries
the prescribed assembled packet. No existence conclusion is inferred
from the finite gluing conditions alone.
\end{proof}

\begin{remark}\label{rem:conditional_not_proof}
H5 remains a conditional premise about the same fixed candidate. It is
not a formal consequence of the finite symplectic bridge, conductor
profiles, or a choice of noncompact root. This proposition adds no
unconditional pair-factor obligation to the construction conjecture.
\end{remark}

\subsection{Exact open boundary}

The following issues delimit the unrestricted program. The first four construct and normalize the finite inputs for broader source classes. Conjecture 66 asks for arithmetic realization once H1--H4 are supplied. The fifth is the independent route to replacing the retained conditional H5 input.

\begin{enumerate}
\item Construct a choice free finite symplectic envelope from an arbitrary irreducible wild inertia atom, with its ramification filtration and tame stabilizer action.
\item Prove a canonical equivariant symplectic isometry from that envelope to the type side quotient \(J_\theta^1/H_\theta^1\).
\item Recover the canonical beta extension phase in every break, including the required dyadic quadratic refinement.
\item Prove matrix holonomy compatibility for the resulting packet transport system.
\item Compare the full family of pair \(L\), epsilon, and gamma factors without invoking the established correspondence.
\end{enumerate}

\begin{figure}[t]
\centering
\includegraphics[width=.78\textwidth]{16_boundary_ladder.pdf}
\caption{The proved finite layers and the remaining independent local correspondence steps.}
\label{fig:boundary}
\end{figure}

The first item is a construction problem. The second is a comparison problem. The third is a normalization problem. The fourth is a descent problem. The fifth concerns proving the local factor premise needed for identification. Conjecture 66 does not assert this fifth step. None is replaced by a larger finite computation.

\section{Examples and finite certificates}\label{sec:examples}

The examples in this section illustrate the exact finite formulas. They are not used as proofs of the general theorems.

\begin{example}[A degree four tame orbit]\label{ex:tame_four_again}
Over \(\FF_2\), the exponent orbit
\[
1\longmapsto2\longmapsto4\longmapsto8\longmapsto1
\]
has length four. It therefore gives a tame atom induced from the unramified extension of degree four. The Frobenius transfer has one four cycle, and the number of regular degree four orbits is three by \cref{prop:mobius_count}.
\end{example}

\begin{example}[Tensor packets of lengths four and six]\label{ex:tensor_46}
For \(n=4\) and \(m=6\),
\[
\gcd(4,6)=2,
\qquad
\lcm(4,6)=12.
\]
The tensor product splits into two interaction packets of length twelve. Their total dimension is \(2\cdot12=24\), equal to the product \(4\cdot6\).
\end{example}

\begin{example}[Rank one Heisenberg group]\label{ex:heisenberg_rank_one}
Take \(V=\FF_3^2\). The Heisenberg group has order \(3^3=27\), and its irreducible representation with a fixed nontrivial central character has dimension three. The trace vanishes on every noncentral element. Averaging the translation subgroup gives a rank one projector.
\end{example}

\begin{example}[Nondegenerate twisted algebra]\label{ex:twisted_plane}
Let \(A=\FF_p^2\) with its standard alternating form. A cocycle with that commutator form gives
\[
\CC^c[A]\cong M_p(\CC).
\]
The unique irreducible projective representation has dimension \(p\).
\end{example}

\begin{example}[Dyadic phase ambiguity]\label{ex:dyadic_plane}
On a symplectic plane over \(\FF_2\), choose refinements \(q_0\) and \(q_1\) with Arf invariants zero and one. Their polar forms are the same, while
\[
\sum_x(-1)^{q_0(x)}=2,
\qquad
\sum_x(-1)^{q_1(x)}=-2.
\]
Thus the alternating packet alone cannot select the dyadic sign.
\end{example}

\begin{example}[One cycle gluing obstruction]\label{ex:one_cycle}
Let a connected graph have one non tree edge. Tree transport determines every vertex map from one root map. The final edge equation holds exactly when the two holonomies around the unique fundamental cycle intertwine. Changing the scalar on that edge changes the cycle holonomy and produces a genuine obstruction.
\end{example}

\section{Scope of the finite reconstruction}\label{sec:scope}

This chapter proves the following statements under their displayed hypotheses.

\begin{enumerate}
\item Tame irreducible Weil atoms are transitive Frobenius packets induced from unramified extensions.
\item Regular tame packets have the stated Möbius count.
\item Finite cuspidal blocks have exact primitive central projectors.
\item Tame tensor products satisfy the greatest common divisor and least common multiple packet formula.
\item Tame local polynomials, conductors, and finite Gauss phases are explicit.
\item Nondegenerate scalar projective commutators give full matrix twisted group algebras and unique finite Heisenberg representations.
\item A declared one break Clifford datum has a choice independent radical free symplectic envelope.
\item An equivariant symplectic isometry gives a Heisenberg bridge unique up to scalar.
\item One generator phase removes the scalar discrepancy for a cyclic stabilizer.
\item Odd finite Weil phases and dyadic Arf signs are multiplicative across orthogonal break packets.
\item Matrix holonomy on a fundamental cycle basis is the exact finite gluing obstruction.
\end{enumerate}

This chapter does not independently prove the local Langlands correspondence, simple character theory, the canonical beta extension, or universal wild companion existence. It does not construct the Galois to type symplectic isometry for an arbitrary wild atom. It does not compare all pair local factors in the wild case.

\section{Restricted arithmetic realization theorems}
\label{sec:arithmetic_progress}

The later constructions provide actual arithmetic candidates and exact
finite packet recovery in substantial restricted families. They do not
establish \cref{conj:arithmetic_realization} for every admitted source.
Their detailed ordinary proofs are typeset in the appendices, and the
original reviewed Markdown files accompany this edition. In these
appendices \(F\) denotes the base local field, while \(K=F[\beta]\)
denotes a parameter field; these are local notation conventions.
All references to ``claim 66'' in the retained notes mean the general
endpoint now isolated as Conjecture 66.

\subsection{The odd-prime trace-zero construction}

\begin{theorem}[Restricted odd-prime arithmetic realization]
\label{thm:odd_arithmetic}
Let \(p\) be an odd prime, \(F=\FF_p((t))\), fix \(d\ge1\), and let
\[
0<h_1<\cdots<h_d,\qquad p\nmid h_i,
\qquad q_i=(p+1)h_i,\qquad N=p^d.
\]
There is an explicit maximal simple stratum in \(M_N(F)\) with totally
ramified full-degree parameter field, a literal defining sequence of
degrees \(p,p^2,\ldots,p^d\), and levels
\(n_j=p^{d-j}q_{d-j+1}\). Every native parameter and consecutive
correction has trace zero over its recursive base field.
The stratum admits a full simple character trivial on \(U_F^1\).
Its fixed canonical beta extension can be completed to an actual
irreducible admissible supercuspidal \(\pi\) with any prescribed smooth
central quasicharacter trivial on \(U_F^1\), when its completion choices
are free or already compatible.

Assume actually supplied labeled full-character extensions with whole
two-dimensional symplectic quotients whose balanced complex-scalar
central product has \(d\) planes, one vertex, and trivial stabilizer.
Then the same \(\pi\) admits a full scalar-pushout realization of that
packet and exact recovery of its unchanged inducing representation.
The prescribed central character must lie in the covered tame class.
No literal embedding of the old matrix groups, canonical native marking,
or prescribed numerical host-depth placement is inferred.
\end{theorem}

\begin{proof}
The explicit Artin--Schreier recursion and its normalized filtered lift
construct the maximal parameter and literal defining sequence; see
\cref{app:odd_tower}. Its trace transport identity makes every native
trace zero, rather than relying on amplification by \(p\). Native-rank
full-character construction, full beta-determinant transfer, trace-zero
correction, and maximal-field character extension give the scalar
triviality and central completion; see \cref{app:odd_character}.
The actual leading symplectic layer has dimension
\((p-1)p^{d-1}\ge2d\) and survives every later approximation
(\cref{app:odd_capacity}). A labeled symplectic embedding gives a
center-fixing full scalar-pushout isomorphism. Finite Stone--von Neumann
decomposition separates the source factor from its orthogonal
complement. Maximal-isotropic averaging on the complementary Heisenberg
factor selects a one-dimensional representation line and hence a
source-sized sector. The retained induction/idempotent recovery returns the same
\(\Lambda\) and \(\pi\), as proved in \cref{app:odd_synthesis}.
\end{proof}

For \(p>3\), a native minimal rank-\(p\) type quotient has dimension
\(p-1\), so it is not automatically a plane. The theorem requires
genuinely supplied whole-plane data. Independence of source tensor
factors is also an input. An arbitrary single-factor wild determinant
need not be tame.

Under additional actual-inertia hypotheses, genuine scalar phases on
the entire source carrier are realized by actual central matrices of a
finite cover and descend to the stipulated operators. A direct matrix
homomorphism from the stabilizer itself need not exist; the exact
statement and descent proof are in \cref{app:odd_scalar}.

\subsection{The weighted dyadic construction and enlarged carriers}

\begin{theorem}[Restricted dyadic marked assembly realization]
\label{thm:dyadic_arithmetic}
Let \(F=\FF_2((t))\), fix an additive character of conductor
\(t\OO_F\), fix \(d\ge1\), and let \(h_1<\cdots<h_d\) be positive odd integers.
Put \(q_i=3h_i\), \(N=2^d\), and
\(n_j=2^{d-j}q_{d-j+1}\). The explicit recursive matrix generates a
separable totally ramified degree-\(N\) maximal parameter field,
normalizes the complete lattice chain, and has the literal defining
sequence of degrees \(2,4,\ldots,2^d\). It admits a full simple character
trivial on \(U_F^1\) and an actual supercuspidal completion with any
prescribed central quasicharacter trivial on \(U_F^1\), when the
noncompact completion is free or already compatible.

For actually supplied full marked whole-plane scalar-pushout data with
one vertex, trivial stabilizer, and prescribed source central character
trivial on \(U_F^1\), an arithmetic assembly carrier exists
at every length and recovers the same candidate: inside its original
type for \(d\ge3\), and in explicit enlarged inducing carriers for
\(d=1,2\). Every supplied full character value and transported whole
finite subgroup, lift, square, and phase marking is retained at the
represented marking level. Any central-character kernel already factored
out remains recorded separately; scalar pushout does not recover its
invisible elements. This does
not identify the transported marking with an independently prescribed
canonical native marking.
\end{theorem}

\begin{proof}
The Artin--Schreier filtered lift has native trace one. Consequently
each native correction \(\delta\) satisfies
\[
v_F(\operatorname{tr}\delta)+s\ge1 \quad (s=1),
\]
which makes its additive scalar correction trivial on \(U_F^1\).
The actual field and defining-sequence proof is
\cref{app:dyadic_tower}. Full beta-determinant transfer and character
extension at each new maximal native field prove the full-character
claim, including the odd-critical endpoint
\(2(\lfloor r/2\rfloor+1)-r=1\); see \cref{app:dyadic_character}.
The exact original quotient has dimension \(2^d-2\), which accommodates
the required \(2d\) source dimensions precisely for \(d\ge3\)
(\cref{app:dyadic_capacity}).

At rank two, with \(q=2m+1\), enlarge
\(J=U_K^1U_A^{m+1}\) to \(P=U_K^1U_A^m\). The actual character
ratio \(\lambda=\theta^g/\theta\) for \(g\notin J\) is nontrivial
and sign-valued. With \(H_0=\ker\lambda\), the quotient is
\(P/H_0\simeq\FF_2^2\), and \(\Ind_J^P\theta\) is an
irreducible two-dimensional Heisenberg carrier
(\cref{app:rank2}). At rank four, the actual outer quadratic
centralizer gives \(P=J U_B^m\), \([P:J]=2\), and
\(P/H_0\simeq\FF_2^4\), with carrier \(\Ind_J^P\eta\) of
dimension four (\cref{app:rank4}). The finite scalar center can be
larger than \(\mu_2\); no literal extraspecial quotient is claimed.

At both ranks \(K^\times\) normalizes \(P\). Thus
\(\mathcal N=K^\times J\) is normal of index two in
\(\mathcal M=K^\times P\), and
\[
\Sigma=\Ind_{\mathcal N}^{\mathcal M}\Lambda,
\qquad
\pi\simeq\cInd_{\mathcal M}^{\GL_N(F)}\Sigma.
\]
Induction in stages preserves the full original \(\Lambda\), its
central character, and its noncompact root; see
\cref{app:rank4_stages,app:lowrank_contract}. The full marked
scalar-pushout comparison applies to these actual enlarged pairs as
well as to the original higher-rank types (\cref{app:dyadic_markings}).
The full synthesis and scope are in \cref{app:dyadic_synthesis}.
\end{proof}

The original low-rank maximal-type capacity obstruction remains true.
The new pair is a common arithmetic assembly carrier; it is not renamed
as the original maximal-simple-character pair. For an independently
fixed native quadratic marking, a preassigned embedding must satisfy
\(q_W(f(v))=q_{\rm src}(v)\). At equal capacity, a free embedding
requires equal Arf invariant. Spare capacity for \(d\ge4\) admits
either source Arf type when the embedding is free. Retaining only the
central \(\{\pm1\}\) would erase the distinction between whole marked
finite extensions such as \(D_8\) and \(Q_8\).

\subsection{Scope, numerical invariants, and evidence}

For the same candidates the constructions give
\[
\begin{array}{c|c|c}
 & p\text{ odd} & p=2\\ \hline
\dim_{\FF_p}(J^1/H^1) & p^d-1 & 2^d-2\\
\Swan(\pi)&p^{d-1}(p+1)h_d&3\cdot2^{d-1}h_d
\end{array}
\]
and torsion number one. Here \(\Swan\) denotes the Swan conductor,
and \(a\) below denotes the ordinary pair conductor. An epsilon exponent
for a specified additive character is a different quantity. The separate
progress paper, Sections 22--24, reconciles these conventions and the
actual-order hypotheses. Their self-pair conductor formulas are
\[
\begin{aligned}
a(\pi^\vee\times\pi)
 &=p^{2d}-1+\sum_{i=1}^d(p^{2i}-p^{2i-2})h_i
 &&(p\text{ odd}),\\
a(\pi^\vee\times\pi)
 &=N^2-1+\sum_{i=1}^d3\cdot4^{i-1}h_i
 &&(p=2).
\end{aligned}
\]
These identities do not prove H5, identify an endo-class, or construct
a canonical source-to-host marking. At rank eight, the retained
skipped-degree host has the same displayed numerical profiles as the
full tower and a smaller Heisenberg quotient. The explicit construction,
rather than numerical matching alone, supplies the arithmetic object.

The proofs use the stated Bushnell--Kutzko simple-type theorems
\cite{BushnellKutzko1993} and the retained conductor inputs. Exact source
locations and hypothesis checks are recorded in the component proofs
and accompanying audits. The finite rank-nine and rank-eight runs check
particular matrices, traces, normalized lifts, and filtration grades.
They are bounded checks, not arbitrary-length proofs. No new Lean build
of these arithmetic constructions is claimed.

Nontrivial prescribed symplectic stabilizer actions, independently fixed
canonical native markings, arbitrary source/transport diagrams, and any
required numerical layer placement remain outside the restricted results.
The current totally ramified host cannot implement nontrivial
prime-to-\(p\) symplectic actions by the stronger character-preserving
inner-normalizer mechanism; this obstruction is stated precisely in
\cref{app:stabilizer_contract}. That stronger mechanism is not silently
imposed on the original conjecture. H5 remains conditional for the same
fixed candidate.

The accompanying archive retains the mathematical reviews, computational evidence, and source records underlying the displayed scope.

\clearpage
\appendix
\section*{Detailed ordinary proofs and construction records}
\phantomsection
\addcontentsline{toc}{section}{Detailed ordinary proofs and construction records}

The following appendices typeset the complete selected current construction
proofs from Blocks 166 and 167. Their source notation, numbered source
subsections, exact hypotheses, and mathematical scope are retained.
The synthesis appendices state the final combined results; component notes
also document narrower lemmas, alternative routes, and failed interpretations.
In particular, the separate classical arithmetic-induction application in
the leading-capacity note is not an input to the direct constructor.

Historical research identifiers, filenames, and references to earlier
blocks are preserved to make the arguments traceable to the accompanying
Markdown files. Earlier ``pending'' review language is superseded by the
retained internal review records included with those source files, at their
recorded dates and scopes. No external peer review is claimed.
Negative capacity conclusions concern the original maximal-type carrier;
the later low-rank enlargement does not repeal them.
These appendices are ordinary mathematical proofs. Neither their typesetting
nor the inherited formalization certificates constitute a fresh Lean proof
of these arithmetic results or of Conjecture 66.

\clearpage
\section{Odd-prime synthesis and complete scope}\label{app:odd_synthesis}

Source record: \path{Paper_XI_Base_Realization165.zip}.

This is a synthesis of the independently reviewed proofs in this
continuation. The result is an ordinary mathematical proof, using the
stated classical simple-type theorems. It is not a Lean certificate.
Universal claim 66 remains open. Original H5 remains the original
conditional premise about the same fixed candidate; no unconditional
pair-factor comparison is asserted here.

\hypertarget{odd_synthesis-the-arithmetic-construction-now-proved}{%
\subsection{1. The arithmetic construction now
proved}\label{odd_synthesis-the-arithmetic-construction-now-proved}}

Let \(p\) be odd, \(F=\mathbb F_p((t))\), and \[
0<h_1<\cdots<h_d,\qquad p\nmid h_i,\qquad
q_i=(p+1)h_i,\qquad N=p^d.
\] For every such tuple there is an explicit matrix \(\beta\in M_N(F)\)
whose field \(K=F[\beta]\) is totally ramified of degree \(N\),
normalizes the complete lattice chain, and gives a maximal simple
stratum. It has a literal defining sequence of degrees
\(p,p^2,\ldots,p^d\), with leading level \(n_1\) and successive critical
levels \(n_2,\ldots,n_d\) given by \[
n_j=p^{d-j}q_{d-j+1}\quad(1\le j\le d).
\] Every native parameter in the construction, every amplified
approximation, and every consecutive correction has matrix trace zero
over its base field.

For this actual stratum there exists a full simple character \(\theta\)
with \(\theta|_{U_F^1}=1\). Retaining its canonical beta extension, the
residue character and noncompact extension can be selected so that the
resulting irreducible admissible supercuspidal \(\pi\) has any
prescribed smooth central quasicharacter \(\omega\) trivial on
\(U_F^1\). These completion choices must be free or already compatible;
incompatible immutably supplied choices are not replaced.

The same candidate has \[
\begin{aligned}
\operatorname{rank}(\pi)&=N, &t(\pi)&=1,\\
\operatorname{Sw}(\pi)&=p^{d-1}(p+1)h_d, &
\dim_{\mathbb F_p}(J^1/H^1)&=N-1,\\
\dim\eta&=p^{(N-1)/2},\\
a(\pi^\vee\times\pi)
 &=p^{2d}-1+\sum_{i=1}^d(p^{2i}-p^{2i-2})h_i.
\end{aligned}
\] These numerical identities do not identify a Weil representation, an
endo-class, or the pair epsilon factors required by H5.

\hypertarget{odd_synthesis-construction-and-proof}{%
\subsubsection{Construction and proof}\label{odd_synthesis-construction-and-proof}}

At a recursive base \((L,t_L)\), introduce \[
E=L(\alpha),\qquad \alpha^p-\alpha=t_L^{-q_d}.
\] This is a separable totally ramified extension of degree \(p\), with
\(v_E(\alpha)=-q_d\) and \(\operatorname{Tr}_{E/L}\alpha=0\). Choose
\(1\le b<p\) with \(q_db\equiv-1\pmod p\), set \(a=(q_db+1)/p\), and put
\(z=t_L^a\alpha^b\). Then \(z\) is an explicit uniformizer. Write its
Eisenstein polynomial as \(f(X)=X^p+\cdots+a_0\).

Let \(P_0\) extract the constant coefficient in the basis
\(1,z,\ldots,z^{p-1}\). The normalized corestriction
\(s=\kappa\operatorname{pr}\), where \(\kappa=-f'(z)/z^{p-1}\),
satisfies \[
u_0=s(P_0)=a_0/z^p\in\mathcal O_E^\times,
\qquad
\operatorname{Lift}_E(C)=u_0^{-1}CP_0,
\qquad s(\operatorname{Lift}_E(C))=C.
\] The lift preserves the exact filtration valuation. Its crucial
additional identity is \[
\operatorname{tr}_L(\operatorname{Lift}_E(C))
 =\operatorname{Tr}_{E/L}
   \bigl(\kappa^{-1}\operatorname{tr}_E(C)\bigr).
\] Define \[
B(L,t_L;h_1,\ldots,h_d)
 =\alpha I_{p^{d-1}}+
   \operatorname{Lift}_E\bigl(B(E,z;h_1,\ldots,h_{d-1})\bigr),
\] with the degree-\(p\) multiplication matrix as the base case.
Induction proves trace zero at every native rank. This is stronger than
the automatic zero trace obtained by multiplying any matrix multiplicity
by \(p\) in characteristic \(p\).

For the inductive field step \(d\ge2\), put \(n=p^{d-1}q_d\) and
\(R=p^{d-2}q_{d-1}\). The derived matrix is the recursively constructed
maximal-field parameter, with \(k_0=-q_1\). The inequalities \(n>pR>2R\)
imply \(-2R>-n-q_1\). Bushnell--Kutzko Remark 2.2.7 therefore applies to
the actual displayed matrix. It supplies its maximal field,
normalization, and critical exponent. Valuation preservation and literal
truncation of the recursion supply the whole defining sequence. The full
details, quotient count, and conductor dependencies are in
\path{tower/TRACE_ZERO_AS_TOWER166.md} and
\path{tower/RECURSIVE_TOWER166.md}.

The full-character argument uses a specific order of operations. First
construct a full character at native rank \(p^j\), trivial on \(U_F^1\).
Transfer its full beta-determinant unit character to multiplicity \(p\),
using BK 3.6.13(ii) and 3.6.14. On a scalar \(u\in U_F^1\), this gives
\(\theta_j(u^p)=1\). Next add the actual trace-zero correction, using BK
3.3.18. It remains trivial on the scalar intersection of its domain.
Finally extend at the next \textbf{maximal native parameter field}: the
forced deep character and the trivial character of \(U_F^1\) agree on
their intersection, so they extend through a finite abelian unit
quotient. BK 3.2.5 constructs the full simple character.

This procedure never substitutes an unjustified extension through a
nonmaximal centralizer. It proves full-character compatibility,
including shallow units, rather than merely cancellation of a leading
trace. The residue factor fixes \(\omega\) on \(\mathbb F_p^\times\); an
\(N\)-th-root choice for the noncompact generator fixes \(\omega(t)\).
Every choice is recorded before H5. See
\path{central/TRACE_ZERO_CHARACTER_INDUCTION166.md} for the
complete induction.

\hypertarget{odd_synthesis-the-finite-packet-that-this-same-candidate-realizes}{%
\subsection{2. The finite packet that this same candidate
realizes}\label{odd_synthesis-the-finite-packet-that-this-same-candidate-realizes}}

Assume the one-vertex scalar-pushout contract: there are \(d\) labeled,
actually supplied full character extensions with whole symplectic
quotients \(V_i\) of dimension two over \(\mathbb F_p\). Their balanced
\(\mathbb C^\times\)-central product has quotient
\(V_{\mathrm{src}}=\bigoplus_iV_i\) and Heisenberg carrier dimension
\(p^d=N\). Initially take the declared stabilizer \(C=\{1\}\).

For \(p=3\), this applies directly to the retained cubic family with
whole rank-three break-type quotients. For \(p>3\), it applies only when
the supplied plane data genuinely satisfy the stated original bridge
contract. A native minimal rank-\(p\) type quotient has dimension
\(p-1\), so it must not automatically be called a plane.

The actual candidate from Section 1 has a nondegenerate leading subspace
\[
S\subset J^1/H^1,\qquad
\dim_{\mathbb F_p}S=(p-1)p^{d-1}\ge2d.
\] Here \(S\) is the image of the concrete congruence subgroup at half
the leading even level. BK's high-depth equality preserves this subgroup
and its commutator pairing through every later approximation; see
\path{capacity/LEADING_LAYER_CAPACITY166.md}, Section 1. This part
of that note applies to the direct constructor and does not require its
separate classical arithmetic-induction application.

Choose a labeled symplectic embedding \(V_{\mathrm{src}}\to S\). For
each original full character form \[
\mathcal P_i=(\mathbb C^\times\times J_i^1)/
 \{(\theta_i(h)^{-1},h):h\in H_i^1\}.
\] The balanced product and the corresponding host restriction have the
same alternating commutator. The center-fixing extension isomorphism,
obtained with matched Darboux lifts and complex roots, therefore
preserves every supplied value \(\theta_i(h)\), including the shallow
values.

Write \(W=V_{\mathrm{src}}\perp V_{\mathrm{src}}^\perp\). Finite
Stone--von Neumann gives
\(\eta_W\simeq\eta_{\mathrm{src}}\otimes\eta_\perp\). Choose an
isotropic eigenspace giving a line in the complementary factor. Its
actual Heisenberg image-algebra projector \(e\) has rank \(N\). The
retained induction/idempotent theorem recovers the \textbf{same}
inducing representation \(\Lambda\) from \(e\Lambda\), and then the same
\(\pi\) by compact induction. No candidate is changed at recovery.

This attachment preserves the full source characters as scalar-pushout
data. It does not identify the old matrix groups with literal subgroups
of the new type, retain the old lower summand conjugacy classes, or
place each source label at its old numerical host depth. Those stronger
requirements have not been proved.

When the retained independent extraspecial tensor source has \(d\ge2\),
its full determinant is trivial on wild inertia: each factor determinant
has order dividing \(p\), and tensor determinants raise it to
\(p^{d-1}\). Thus its central character belongs to the class in Section
1. For \(d=1\), wild triviality of the source determinant must be
checked separately. Independence of the tensor factors is an input, not
a consequence of distinct \(h_i\).

\hypertarget{odd_synthesis-a-proved-extension-to-genuine-scalar-inertia-phases}{%
\subsection{3. A proved extension to genuine scalar inertia
phases}\label{odd_synthesis-a-proved-extension-to-genuine-scalar-inertia-phases}}

There is also a positive result beyond \(C=\{1\}\). Suppose the finite
prime-to-\(p\) stabilizer acts trivially on the source scalar extension,
and its prescribed honest operators on the \textbf{entire} rank-\(N\)
Weil carrier are \(\chi(c)I_N\), each lying in the actual inertia image.
Assume the candidate's central character is the determinant of that same
source under the fixed rank-one reciprocity convention.

Then a finite cover \(\widetilde C\twoheadrightarrow C\) acts by actual
central matrices in the unchanged candidate. Its represented operators
descend to exactly \(\chi\) on \(C\). Hence the attachment and recovery
above are equivariant for that action. A direct matrix homomorphism from
\(C\) itself is not asserted and can fail.

Indeed, for residue cardinality \(q\), let
\(B=\omega(\operatorname{Teich}(k^\times))\). The actual inertia
hypothesis gives \(\chi(C)^N\subset B\). Since \(N\) is a power of \(p\)
and \(\chi(C)\) has order prime to \(p\), this implies
\(\chi(C)\subset B\). Choose \(k_0\) inverse to \(N\) modulo \(q-1\),
and set \[
\delta(u)=\omega(\operatorname{Teich}(u))^{k_0},\qquad
\widetilde C=\{(c,u):\chi(c)=\delta(u)\}.
\] The matrix \(\operatorname{Teich}(u^{k_0})I_N\) acts by \(\chi(c)\).
The kernel acts trivially on the entire sector induction and ambient
compact induction, because these matrices are central. This preserves
the fixed \(\Lambda\), its canonical beta-extension factor, and all
completion choices. It does not assert that the canonical beta extension
alone has scalar character \(\chi\).

The full theorem and the nonsplitting example are in
\path{stabilizer/SCALAR_NATIVE_ATTACHMENT166.md}, with an
independent review in \path{audit/SCALAR_NATIVE_REVIEW166.md}. An
arbitrary abstract phase or a phase on a proper multiplicity space does
not satisfy its hypotheses merely because it is scalar.

\hypertarget{odd_synthesis-what-was-repaired-and-what-remains}{%
\subsection{4. What was repaired, and what
remains}\label{odd_synthesis-what-was-repaired-and-what-remains}}

The old fixed-field central argument genuinely failed. At \(p=3\),
\((h_1,h_2)=(1,5)\), its naturally transferred full character takes the
value \(\zeta_3\ne1\) on \(1+t^2\), although the desired wild center is
trivial. A shallow character choice cannot repair a forced value. The
two-stage changing-field repair is retained in
\path{central/CENTRAL_REPAIR166.md}; the trace-zero tower above
resolves this central issue uniformly in every length.

The construction does not use an inverse established correspondence or
the classical arithmetic-induction/descent candidate to select \(\pi\).
It does use the named BK construction and type theorems, the retained
type/formal-degree conductor formulas, and exact sector recovery. The
capacity note separately retains an alternative consequence for a
candidate supplied by the substantial classical arithmetic constructor;
that is a different route and is not an input to Section 1.

The unresolved universal obligations are precise:

\begin{enumerate}
\def\labelenumi{\arabic{enumi}.}
\tightlist
\item
  \textbf{Nontrivial prescribed symplectic stabilizer actions.} For this
  totally ramified full-degree host, the stronger inner-normalizer
  interpretation forces prime-to-\(p\) actions on the quotient to be
  trivial. It therefore cannot implement a genuinely nontrivial supplied
  symplectic action. The raw original contract does not automatically
  require that stronger interpretation; it still needs an actual valid
  arithmetic comparison, not relabeling of external operators.
\item
  \textbf{Arbitrary admissible source packets and transport diagrams.}
  The proved attachment has the whole-plane, one-vertex, finite
  scalar-pushout hypotheses above. It does not establish all source
  classes, arbitrary graphs, or prescribed numerical layer placement.
\item
  \textbf{The dyadic construction with all markings.} A separate
  field-level recursion and some capacity statements exist at \(p=2\).
  The odd-prime trace-zero/full-character theorem here does not settle
  that case.
\end{enumerate}

Original H5 stays exactly where it was: if supplied for this fixed
candidate and the stipulated tests, the retained conditional endpoint
applies. No claim that the displayed profiles imply H5, or that one of
the noncompact root choices must satisfy H5, is made. The unresolved
construction issues above are not renamed epsilon-factor obligations.

\hypertarget{odd_synthesis-evidence-and-reproducibility}{%
\subsection{5. Evidence and
reproducibility}\label{odd_synthesis-evidence-and-reproducibility}}

The arbitrary-length theorem is proved symbolically and independently
reviewed. The exact \(p=3,d=2,(h_1,h_2)=(1,2)\) computation in
\texttt{tower/as\_rank9check/} verifies the displayed matrix, trace
identities, corestriction, and levels; it is a finite check, not the
universal proof. The separate rank-27 run in \texttt{tower/run\_rank27/}
checks the earlier non-trace-zero recursion and must not be described as
an AS central completion example. No Lean build was run.

\path{STATUS166.json} records the closed and open obligations. The
component audits identify the exact reviewed files and hashes. The
dependency file pins prior checkpoints and the authorized primary-source
copy without redistributing that publication.

\clearpage
\section{The trace-zero Artin--Schreier tower}\label{app:odd_tower}


\hypertarget{odd_tower-exact-statement}{%
\subsection{Exact statement}\label{odd_tower-exact-statement}}

Let p be odd, F=F\_p((t)), and
0\textless h\_1\textless\ldots\textless h\_d with p not dividing any
h\_i. Put q\_i=(p+1)h\_i. The constructor below gives an actual matrix
B\_d over F of size p\^{}d with exactly the field, simplicity,
approximation-level, odd-prime quotient, and conductor conclusions of
\path{RECURSIVE_TOWER166.md}. In addition, its matrix trace is exactly zero.
Every native shorter constructor and every amplified arithmetic
approximation also has matrix trace zero. Therefore every difference
between consecutive arithmetic approximations has matrix trace zero.

The trace statement is an identity over the current coefficient field,
not merely a congruence, a field trace inferred from characteristic
zero, or a division by the matrix size. All coefficient fields have
characteristic p, so division by p\^{}j would be invalid.

\hypertarget{odd_tower-an-explicit-separable-minimal-extension-with-a-traceless-parameter}{%
\subsection{1. An explicit separable minimal extension with a traceless
parameter}\label{odd_tower-an-explicit-separable-minimal-extension-with-a-traceless-parameter}}

At an arbitrary recursive step let L have residue field F\_p and
uniformizer t\_L. Given q\textgreater0 prime to p, set

\begin{Verbatim}[fontsize=\small,breaklines=true,breakanywhere=true]
E=L(alpha),  alpha^p-alpha=t_L^(-q).
\end{Verbatim}

The right side is not an Artin--Schreier coboundary in L: the negative
valuation of x\^{}p-x, if x has negative valuation, is divisible by p.
Thus the polynomial is irreducible of degree p.~The extension is
separable, and the valuation equation forces ramification index p and

\begin{Verbatim}[fontsize=\small,breaklines=true,breakanywhere=true]
v_E(alpha)=-q.
\end{Verbatim}

It is therefore totally ramified and alpha is minimal over L. For odd p
its field trace is

\begin{Verbatim}[fontsize=\small,breaklines=true,breakanywhere=true]
Tr_(E/L)(alpha)=sum_(c in F_p)(alpha+c)=0.
\end{Verbatim}

Choose the unique b in \{1,\ldots,p-1\} with qb congruent to -1 modulo
p, put a=(qb+1)/p, and define

\begin{Verbatim}[fontsize=\small,breaklines=true,breakanywhere=true]
z=t_L^a alpha^b.
\end{Verbatim}

Then v\_E(z)=pa-qb=1, so z is a uniformizer and E=L(z). Its monic
minimal polynomial \(f(X)=X^p+a_{p-1}X^{p-1}+\cdots+a_0\)
is Eisenstein. It is an explicit polynomial: compute the product over
the p conjugates

\begin{Verbatim}[fontsize=\small,breaklines=true,breakanywhere=true]
f(X)=product_(c in F_p) (X-t_L^a(alpha+c)^b)
\end{Verbatim}

and reduce coefficients using \(\alpha^p-\alpha=t_L^{-q}\).
The invariant coefficients belong to L. This is a finite algebraic
instruction, with no choice of an unknown extension or matrix.

\hypertarget{odd_tower-normalized-corestriction-and-an-explicit-filtered-right-inverse}{%
\subsection{2. Normalized corestriction and an explicit filtered right
inverse}\label{odd_tower-normalized-corestriction-and-an-explicit-filtered-right-inverse}}

Use the basis 1,z,\ldots,z\^{}(p-1) of E over L. If T\_ij is its
coefficient matrix unit and pr is the trace corestriction, then

\begin{Verbatim}[fontsize=\small,breaklines=true,breakanywhere=true]
pr(T_ij)=z^i z_j^*,
z_j^*=sum_(k=j+1)^p a_k z^(k-j-1) / f'(z),  a_p=1.
\end{Verbatim}

Set

\begin{Verbatim}[fontsize=\small,breaklines=true,breakanywhere=true]
kappa=-f'(z)/z^(p-1),  s=kappa pr.
\end{Verbatim}

The displayed dual-basis formula gives

\begin{Verbatim}[fontsize=\small,breaklines=true,breakanywhere=true]
s(T_ij)=-z^(i-j) + terms of strictly larger E-valuation.
\end{Verbatim}

Indeed v\_E(a\_k)\textgreater=p for 1\textless=k\textless p, so the
contribution indexed by a\_k has valuation at least k greater than the
first contribution. Applying s entrywise on E\^{}m gives a B-bimodule
corestriction from End\_L(E\^{}m) to B=End\_E(E\^{}m). With the complete
interleaved period-pm L-chain and period-m E-chain, it preserves grades
and satisfies \(s(P^r)=Q^r\).

Let P0 extract the constant coefficient on every E-coordinate. The
constant dual coefficient satisfies

\begin{Verbatim}[fontsize=\small,breaklines=true,breakanywhere=true]
z_0^*=-a_0/(z f'(z)),
u_0=s(P0)=kappa z_0^*=a_0/z^p in O_E^times.
\end{Verbatim}

For any C in M\_m(E), define the actual L-linear matrix

\begin{Verbatim}[fontsize=\small,breaklines=true,breakanywhere=true]
Lift_E(C)=u_0^(-1) C P0.
\end{Verbatim}

Then

\begin{Verbatim}[fontsize=\small,breaklines=true,breakanywhere=true]
s(Lift_E(C))=C.
\end{Verbatim}

This lift preserves the exact filtration valuation. To see this, first
multiply C by the unit u\_0\^{}(-1), which does not change its valuation
in the E-chain. Expand each entry in the basis 1,z,\ldots,z\^{}(p-1).
The only nonzero columns of Lift are its constant-coefficient columns.
The L-grade of a coefficient c\_i z\^{}i E\_uv is

\begin{Verbatim}[fontsize=\small,breaklines=true,breakanywhere=true]
pm*v_L(c_i)+mi+u-v,
\end{Verbatim}

equal to its E-grade, and distinct i have valuations distinct modulo p.
Thus the minimum cannot cancel. This also proves the stated surjectivity
\(s(P^r)=Q^r\) without a compactness or lifting-existence
assumption.

\hypertarget{odd_tower-exact-trace-identity}{%
\subsection{3. Exact trace identity}\label{odd_tower-exact-trace-identity}}

For y in E, extraction of its constant coefficient is

\begin{Verbatim}[fontsize=\small,breaklines=true,breakanywhere=true]
coeff_0(y)=Tr_(E/L)(z_0^* y).
\end{Verbatim}

The ordinary L-matrix trace of the rank-one operator yP0 is coeff\_0(y).
Consequently, for every C in M\_m(E),

\begin{Verbatim}[fontsize=\small,breaklines=true,breakanywhere=true]
tr_L(Lift_E(C))
   =Tr_(E/L)(z_0^* u_0^(-1) tr_E(C))
   =Tr_(E/L)(kappa^(-1) tr_E(C)).             (T)
\end{Verbatim}

Here tr\_E(C) is the ordinary m-by-m matrix trace over E. In particular,
if tr\_E(C)=0, then the lift has ordinary L-matrix trace zero exactly.

\hypertarget{odd_tower-the-recursive-actual-matrix}{%
\subsection{4. The recursive actual
matrix}\label{odd_tower-the-recursive-actual-matrix}}

Define B(L,t\_L;h\_1,\ldots,h\_d) by induction. For d=1 use the
multiplication matrix of the alpha from Section~1 with q=q\_1. For
d\textgreater=2 use Section~1 with q=q\_d, put m=p\^{}(d-1), recursively
construct

\begin{Verbatim}[fontsize=\small,breaklines=true,breakanywhere=true]
C=B(E,z;h_1,...,h_(d-1)),
\end{Verbatim}

and set

\begin{Verbatim}[fontsize=\small,breaklines=true,breakanywhere=true]
B(L,t_L;h_1,...,h_d)=alpha I_m+Lift_E(C).       (ASR)
\end{Verbatim}

Everything in (ASR) is explicit in finite field operations and the
displayed degree-p polynomial extensions. It differs from the old
conjugacy-preserving two-summand construction. Preserving old lower
summand conjugacy classes is not claimed.

Inductively tr\_E(C)=0. Equation (T) and the Artin--Schreier trace give

\begin{Verbatim}[fontsize=\small,breaklines=true,breakanywhere=true]
tr_L(B)=m Tr_(E/L)(alpha)+tr_L(Lift_E(C))=0.
\end{Verbatim}

For the induction base the trace is Tr\_(E/L)(alpha)=0. Thus this
argument proves tracelessness in the current base field for every native
rank, not merely tracelessness after an amplification whose size is
divisible by p.

\hypertarget{odd_tower-actual-fields-and-profiles}{%
\subsection{5. Actual fields and
profiles}\label{odd_tower-actual-fields-and-profiles}}

Set gamma=alpha I\_m and b=Lift\_E(C). The minimal outer stratum has

\begin{Verbatim}[fontsize=\small,breaklines=true,breakanywhere=true]
n=p^(d-1)q_d,  k_0(gamma)=-n.
\end{Verbatim}

The inner actual maximal field supplied by induction has

\begin{Verbatim}[fontsize=\small,breaklines=true,breakanywhere=true]
R=p^(d-2)q_(d-1),  v(C)=-R,  k_0(C)=-q_1.
\end{Verbatim}

We have s(b)=C, v(b)=-R, and n\textgreater pR\textgreater2R because p is
odd and the h\_i increase. Hence

\begin{Verbatim}[fontsize=\small,breaklines=true,breakanywhere=true]
-2R > -n-q_1 = k_0(gamma)+k_0(C).
\end{Verbatim}

BK Remark 2.2.7 applies to this pure derived maximal field. It gives the
field generated by the actual displayed gamma+b, with degree and
ramification p\^{}d, residue degree one, normalization of the complete
chain, and k\_0=-q\_1. Its valuation remains -n.~Thus {[}A,n,0,B{]} is
simple.

This repeats the precise argument of \path{RECURSIVE_TOWER166.md} with a new
outer minimal extension and normalized right inverse. No step of that
argument required the old polynomial e\^{}p+t e+t.

\hypertarget{odd_tower-native-approximations-amplification-and-traceless-corrections}{%
\subsection{6. Native approximations, amplification, and traceless
corrections}\label{odd_tower-native-approximations-amplification-and-traceless-corrections}}

For the full rank p\^{}d matrix put, for 1\textless=j\textless=d,

\begin{Verbatim}[fontsize=\small,breaklines=true,breakanywhere=true]
B^[j]=B(F,t;h_(d-j+1),...,h_d) tensor I_(p^(d-j)).
\end{Verbatim}

These are literal truncations of (ASR). At the outer level the same
largest q\_d selects the same Artin--Schreier extension and uniformizer
in all truncations. At the next level the same q\_(d-1) does the same,
and so on. Thus no hidden change of bases or incompatible parameter
fields enters the truncations. Lift\_E commutes with amplification by an
identity because P0 acts separately on each coordinate.

Every native shorter B has trace zero by Section~4. Its amplification
therefore has trace zero. In particular

\begin{Verbatim}[fontsize=\small,breaklines=true,breakanywhere=true]
tr_F(B^[j])=0,
tr_F(B^[j+1]-B^[j])=0.                        (Z)
\end{Verbatim}

The same is true at native rank p\^{}(j+1): the correction between
B(F,t;h\_(d-j),\ldots,h\_d) and the identity-p amplification of
B(F,t;h\_(d-j+1),\ldots,h\_d) has trace zero. This native assertion is
the one relevant when constructing characters before each transfer.

The field degrees and critical levels of \(B^{[j]}\) are exactly

\begin{Verbatim}[fontsize=\small,breaklines=true,breakanywhere=true]
degree p^j,  n_j=p^(d-j)q_(d-j+1).
\end{Verbatim}

Valuation-preservation of every Lift gives

\begin{Verbatim}[fontsize=\small,breaklines=true,breakanywhere=true]
v_A(B^[j+1]-B^[j])=-n_(j+1).
\end{Verbatim}

Thus the approximation and defining-sequence proof of the first note
applies unchanged. The odd-prime J/H dimensions, grade-index formula,
Swan conductor and self-pair conductor depend on these same levels and
centralizer degrees, so all their formulas carry over unchanged.

\hypertarget{odd_tower-precisely-what-the-trace-supplies}{%
\subsection{7. Precisely what the trace
supplies}\label{odd_tower-precisely-what-the-trace-supplies}}

For any displayed correction delta and any scalar u in L\^{}times for
which its additive correction character is defined,

\begin{Verbatim}[fontsize=\small,breaklines=true,breakanywhere=true]
psi_delta(u I)=psi_L((u-1)tr_L(delta))=1.
\end{Verbatim}

This proves triviality of the actual additive correction on the scalar
subgroup of its domain. It does not, by itself, construct a simple
character with prescribed values on larger scalar groups. Compatible
extension, transfer, and intersection arguments for the full simple
character are separate obligations.

The note also does not identify the constructed candidates with any Weil
source, choose them by inverse LLC, realize arbitrary stabilizer
operators, or prove same-candidate H5. Its new output is an explicit
arithmetic defining sequence whose trace vanishes at every native and
amplified stage, suitable as input to that separate character argument.

\hypertarget{odd_tower-primary-dependency}{%
\subsection{Primary dependency}\label{odd_tower-primary-dependency}}

Bushnell--Kutzko, The admissible dual of GL(N) via compact open
subgroups, Remark 2.2.7 (printed p. 55), with the defining-sequence and
ring recursion already pinned in \path{RECURSIVE_TOWER166.md}. Retained source
SHA256:
\path{d69db0827256707a34cd2562af68b9de11fc889ab3f43535ec47b620d58ca9c8}. The
trace-dual and Artin--Schreier calculations above are given in full.

\clearpage
\section{Native-rank full characters and tame central completion}\label{app:odd_character}

Source record: Paper\_XI\_Base\_Realization165. Arithmetic
input: \path{attachment166/tower/TRACE_ZERO_AS_TOWER166.md},
together with its \path{RECURSIVE_TOWER166.md} defining-sequence
proof. Input hashes are recorded in
\path{TRACE_ZERO_CHARACTER_PINS166.json}.

Evidence: ordinary proof using the explicitly constructed arithmetic
parameters and the named Bushnell--Kutzko character and transfer
theorems. No arithmetic-induction/descent constructor or inverse
correspondence is used. No Lean verification is claimed.

\hypertarget{odd_character-statement}{%
\subsection{1. Statement}\label{odd_character-statement}}

Let p be odd, F=F\_p((t)), and

\begin{Verbatim}[fontsize=\small,breaklines=true,breakanywhere=true]
0<h_1<...<h_d,       p not dividing any h_i,
q_i=(p+1)h_i,        N=p^d.
\end{Verbatim}

Let beta\_d be the actual trace-zero Artin--Schreier matrix constructor
in the named tower note. There exists an actual full simple character
theta for its maximal simple stratum such that

\begin{Verbatim}[fontsize=\small,breaklines=true,breakanywhere=true]
theta restricted to U_F^1 is the trivial character.             (1)
\end{Verbatim}

Every step constructing theta is specified below: a finite abelian
character extension, the standard simple-character transfer, or the
explicit additive correction from the displayed defining sequence. The
field, order, and parameter are the actual tower ones, not an
unspecified arithmetic realization.

The actual canonical beta extension can then be completed to a maximal
extended type and supercuspidal pi with any prescribed smooth central
quasicharacter omega\_F trivial on U\_F\^{}1. Its rank, Swan conductor,
torsion number and self-pair conductor remain those proved for the same
tower. The residue and noncompact extension choices must be free or
already compatible with omega\_F.

For p=3, this gives full scalar-pushout attachment and same-candidate
sector recovery for any number d of the supplied cubic source planes,
with declared stabilizer C=\{1\} and the one-vertex contract. All
original full source character values are retained in the attachment. A
general nontrivial stabilizer, transport diagram, or H5 conclusion is
not part of this statement.

\hypertarget{odd_character-the-exact-arithmetic-inputs}{%
\subsection{2. The exact arithmetic
inputs}\label{odd_character-the-exact-arithmetic-inputs}}

For 1\textless=j\textless=d, put

\begin{Verbatim}[fontsize=\small,breaklines=true,breakanywhere=true]
b_j=B(F,t;h_(d-j+1),...,h_d) in M_(p^j)(F),
\end{Verbatim}

where B is the trace-zero constructor. The tower proof gives:

\begin{enumerate}
\def\labelenumi{\arabic{enumi}.}
\item
  K\_j=F{[}b\_j{]} is a totally ramified maximal field of degree p\^{}j,
  normalizing the native complete period-p\^{}j chain A\_j.
\item
  For \(1\le j<d\), in rank p\^{}(j+1), gamma\_j=b\_j tensor I\_p is a simple
  approximation to b\_(j+1), at the critical level
  r\_j=q\_(d-j)\textgreater0 of b\_(j+1). They are equivalent as strata
  at that level.
\item
  For \(1\le j<d\), delta\_j=b\_(j+1)-gamma\_j has grade -r\_j and F-matrix trace zero.
\item
  At native rank p, b\_1 is multiplication by alpha in the explicit
  Artin--Schreier field

\begin{Verbatim}[fontsize=\small,breaklines=true,breakanywhere=true]
E=F(alpha),       alpha^p-alpha=t^(-q_d).
\end{Verbatim}

  Its actual field trace is zero, since p is odd.
\end{enumerate}

For completeness, the tower's exact trace identity is

\begin{Verbatim}[fontsize=\small,breaklines=true,breakanywhere=true]
tr_L Lift_E(C)=Tr_(E/L)(kappa^(-1) tr_E(C)),
\end{Verbatim}

where s=kappa*pr is its normalized trace corestriction and
Lift\_E(C)=s(P0)\^{}(-1) C P0. The recursive inner matrix has trace
zero, so this proves zero trace at every native rank. It is not merely
the automatic zero obtained by amplifying any matrix p times in
characteristic p.~The literal compatibility of the truncations is proved
in Section~6 of the tower note.

\hypertarget{odd_character-a-finite-character-extension-fact}{%
\subsection{3. A finite character-extension
fact}\label{odd_character-a-finite-character-extension-fact}}

Let K/F be a finite totally ramified local extension. Suppose chi is a
smooth character of U\_K\^{}a, and chi is trivial on U\_K\^{}a intersect
U\_F\^{}1. Then chi extends to a smooth character of U\_K\^{}1 trivial
on U\_F\^{}1.

Indeed, on U\_K\^{}a U\_F\^{}1 define chi\_0(ku)=chi(k). Agreement on
the intersection makes this a character. Choose M\textgreater=a such
that chi is trivial on U\_K\^{}M. The character chi\_0 descends to a
subgroup of the finite abelian group \(U_K^1/U_K^M\). It
extends to that whole group by successively adjoining generators and
selecting complex roots. Pulling it back proves the assertion. This also
describes an actual finite procedure for making the extension choices.

This fact will be used only when the parameter field is maximal in the
current native matrix rank. Direct use in a nonmaximal centralizer would
require checking its determinant image and would not follow from the
scalar intersection alone.

\hypertarget{odd_character-base-character-at-native-rank-p}{%
\subsection{4. Base character at native rank
p}\label{odd_character-base-character-at-native-rank-p}}

Let E=K\_1, and let n\_1=q\_d be the native level of alpha. Since n\_1
is even, set a=n\_1/2+1. On U\_E\^{}a the forced field-unit character is

\begin{Verbatim}[fontsize=\small,breaklines=true,breakanywhere=true]
phi(v)=psi_F(Tr_(E/F)(alpha(v-1))).
\end{Verbatim}

It is trivial on U\_E\^{}a intersect U\_F\^{}1 because
Tr\_(E/F)(alpha)=0. Section~3 gives a full character chi\_E of U\_E\^{}1
extending phi and trivial on U\_F\^{}1. No additive-character
normalization is needed for this implication: the trace is exactly zero.

For the native minimal stratum, \(H^1=U_E^1\)
U\_(A\_1)\^{}(a). Define theta\_1 by chi\_E on U\_E\^{}1 and by
psi\_alpha on U\_(A\_1)\^{}a. They agree on the intersection, so BK 3.2.1
supplies the actual full simple character. It satisfies
theta\_1\textbar U\_F\^{}1=1.

\hypertarget{odd_character-transfer-the-full-field-character-before-adding-the-next-layer}{%
\subsection{5. Transfer the full field character before adding the next
layer}\label{odd_character-transfer-the-full-field-character-before-adding-the-next-layer}}

Suppose theta\_j has been constructed on the native rank-p\^{}j stratum
of b\_j, with theta\_j\textbar U\_F\^{}1=1. Its restriction to
U\_(K\_j)\^{}1 is the full beta-determinant character chi\_j, since its
native centralizer is K\_j itself.

Use the standard simple-character transfer to the multiplicity-p
realization gamma\_j=b\_j tensor I\_p in rank p\^{}(j+1). Denote the
transferred full character by Theta\_j. The exact primary result is
BK 3.6.13(ii), with existence in Theorem~3.6.14: transfer preserves the
full beta-determinant character on U\_(K\_j)\^{}1. Consequently its
restriction to the centralizer order is chi\_j composed with
det\_(K\_j).

For every literal scalar u in U\_F\^{}1 this gives

\begin{Verbatim}[fontsize=\small,breaklines=true,breakanywhere=true]
Theta_j(u I)=chi_j(u^p)=theta_j(u^p)=1.            (2)
\end{Verbatim}

The determinant exponent p is essential. We do not replace (2) by an
unsupported pullback along a scalar matrix embedding. This transfer also
deals with nonminimal b\_j: beta-determinant preservation is the general
theorem, not just the minimal-character formula.

\hypertarget{odd_character-add-the-actual-correction-and-extend-at-maximal-rank}{%
\subsection{6. Add the actual correction and extend at maximal
rank}\label{odd_character-add-the-actual-correction-and-extend-at-maximal-rank}}

Let r=r\_j=q\_(d-j), m=floor(r/2)=r/2, beta=b\_(j+1), gamma=gamma\_j,
and delta=beta-gamma. The actual defining sequence gives

\begin{Verbatim}[fontsize=\small,breaklines=true,breakanywhere=true]
D=H^(m+1)(beta)=H^(m+1)(gamma).
\end{Verbatim}

By BK 3.3.18, the character

\begin{Verbatim}[fontsize=\small,breaklines=true,breakanywhere=true]
vartheta=Theta_j restricted to D times psi_delta
\end{Verbatim}

is an actual simple character at level m for beta. For every scalar u in
D intersect U\_F\^{}1, equation (2) and the exact trace identity give

\begin{Verbatim}[fontsize=\small,breaklines=true,breakanywhere=true]
vartheta(u)=Theta_j(u)
                psi_F((u-1)tr_F(delta))=1.        (3)
\end{Verbatim}

The new parameter field K=K\_(j+1) is maximal and totally ramified, and
the native chain has K-period one. Thus

\begin{Verbatim}[fontsize=\small,breaklines=true,breakanywhere=true]
U_K^1 intersect D=U_K^(m+1),
U_F^1 intersect D=U_F^ceil((m+1)/p^(j+1)).
\end{Verbatim}

In particular (3) supplies the exact compatibility condition for
extending vartheta\textbar U\_K\^{}(m+1) while imposing triviality on
all U\_F\^{}1. Apply Section~3 to obtain such a character chi\_K of
U\_K\^{}1.

BK 3.2.5, including its explicit extension formula, now constructs

\begin{Verbatim}[fontsize=\small,breaklines=true,breakanywhere=true]
theta_(j+1)(kd)=chi_K(k) vartheta(d),
k in U_K^1, d in D.
\end{Verbatim}

The proposition supplies normalization of vartheta by the field units,
well-definedness, and the full simple-character condition. This
theta\_(j+1) has the claimed trivial wild center. Induction proves (1)
at the actual final rank p\^{}d.

The proof deliberately performs every arbitrary character extension at
maximal native rank. Extending directly along amplified intermediate
centralizers M\_(p\^{}(d-j))(K\_j) would introduce determinant-image
compatibility conditions. Merely knowing that their ambient matrix trace
is zero would not discharge those conditions. Transfer first, followed
by maximal-rank extension, removes that ambiguity rigorously.

\hypertarget{odd_character-complete-this-same-type-with-the-prescribed-tame-center}{%
\subsection{7. Complete this same type with the prescribed tame
center}\label{odd_character-complete-this-same-type-with-the-prescribed-tame-center}}

Keep the actual canonical beta extension kappa of the resulting theta.
The final K/F is totally ramified of degree N, with residue rank one.
The compact type group satisfies

\begin{Verbatim}[fontsize=\small,breaklines=true,breakanywhere=true]
J/J^1 = k_K^times = F_p^times.
\end{Verbatim}

On U\_F\^{}1 the beta extension acts by theta, hence trivially. For a
prescribed omega\_F trivial on U\_F\^{}1, choose the residue character

\begin{Verbatim}[fontsize=\small,breaklines=true,breakanywhere=true]
sigma=omega_F|F_p^times divided by omega_kappa|F_p^times.
\end{Verbatim}

Its inflation gives lambda=kappa tensor sigma with the right center on
O\_F\^{}times. Complete lambda to an actual extension Lambda\_0 on the
inducing group K\^{}times J using the retained maximal-type completion
theorem. If its center is omega\_0, choose s in C\^{}times with

\begin{Verbatim}[fontsize=\small,breaklines=true,breakanywhere=true]
s^N=omega_F(t)/omega_0(t).
\end{Verbatim}

Multiply Lambda\_0 by the character of K\^{}times J trivial on J and
taking a K-uniformizer to s. The quotient (K\^{}times J)/J is Z, and t
has image N; this proves exact agreement of the resulting central
character with omega\_F on all F\^{}times. Canonical kappa is unchanged.

The resulting pi=c-Ind Lambda is irreducible, admissible and
supercuspidal. All these choices are made before H5. If a residue
character or noncompact extension has been immutably supplied, it must
already satisfy this compatibility; an incompatible datum is not
silently replaced.

By the exact same parameter and defining-sequence calculations,

\begin{Verbatim}[fontsize=\small,breaklines=true,breakanywhere=true]
rank(pi)=N=p^d,
torsion(pi)=1,
Swan(pi)=p^(d-1)(p+1)h_d,
a(pi^vee x pi)=p^(2d)-1
                 + sum_i (p^(2i)-p^(2i-2))h_i.
\end{Verbatim}

No determinant twist changes the depth. These are the proved arithmetic
tower profile formulas, not a proof of H5.

\hypertarget{odd_character-full-packet-attachment-in-arbitrary-cubic-length}{%
\subsection{8. Full packet attachment in arbitrary cubic
length}\label{odd_character-full-packet-attachment-in-arbitrary-cubic-length}}

For p=3, retain d original supplied full characters theta\_i and their
actual character groups H\_i\^{}1 normal in J\_i\^{}1, with whole
symplectic quotients V\_i that are the d labeled planes. Assume the
declared stabilizer is C=\{1\} and the diagram is one vertex. The source
balanced scalar extension has quotient V\_src of dimension 2d and
Heisenberg carrier dimension 3\^{}d, the original rank.

The actual final host has total quotient dimension N-1. Its leading
nondegenerate symplectic layer has dimension

\begin{Verbatim}[fontsize=\small,breaklines=true,breakanywhere=true]
(p-1)p^(d-1)=2*3^(d-1) >= 2d.
\end{Verbatim}

This follows either from the explicit defining-sequence quotient
calculation or the independently proved leading-layer lemma. Choose a
labeled symplectic embedding of the supplied V\_src in that layer.
Write W=V\_src perpendicular-sum V\_perp. Choose a maximal isotropic
subspace L in V\_perp. Its associated eigenspace in the actual
Heisenberg carrier has dimension 3\^{}d, and L-perp/L is V\_src.

Form the actual reduced scalar extension using the full host character
and its extension to the inverse image of L. Form each original full
pushout

\begin{Verbatim}[fontsize=\small,breaklines=true,breakanywhere=true]
P_i=(C^times x J_i^1)/
           {(theta_i(h)^(-1),h): h in H_i^1}.
\end{Verbatim}

The center-fixing symplectic-extension isomorphism, constructed by
matched Darboux lifts and complex roots, maps their balanced product to
that actual reduced scalar extension. In particular each original h in
H\_i\^{}1 maps to the scalar theta\_i(h). Every prescribed shallow and
deep source character value is retained. Source character attachment
does not claim that the original field units are the same literal scalar
matrices in the new host.

Finite Stone--von Neumann supplies the full carrier intertwiner. The
retained sector induction and recovery idempotent then recover this same
Lambda and this same pi from the source-sized sector. The canonical beta
extension, residue factor, central character, and noncompact scalar are
those actually selected in Section~7. No source-independent alternate
candidate is substituted in recovery.

For the retained independent tensor family with d\textgreater=2, each
factor's wild determinant has order dividing p, because the
abelianization of its extraspecial order-p\^{}3 image is elementary
abelian. The tensor determinant raises each factor determinant to
p\^{}(d-1), so the full source determinant is trivial on wild inertia.
Thus the source's prescribed central quasicharacter lies in the class
covered by Section~7. At d=1, matching this central conclusion requires
the source determinant itself to be trivial on wild inertia; it is not
inferred merely by taking a tensor power of length one.

For p\textgreater3, the same arithmetic host and character theorem hold,
and the same attachment works for any actually supplied valid packet
whose whole quotient and markings meet the stated symplectic contract. A
proposed source plane is not automatically a whole rank-p type quotient:
that latter quotient has dimension p-1. No admission of all such plane
sources into original H2 is inferred.

\hypertarget{odd_character-boundary-and-primary-source-check}{%
\subsection{9. Boundary and primary-source
check}\label{odd_character-boundary-and-primary-source-check}}

This provides the previously missing full-character and actual
wild-center construction for the explicit odd-prime arithmetic tower in
every prime-power rank. The proof uses no classical
automorphic-induction/descent candidate. It retains the supplied finite
characters through the declared scalar-pushout attachments.

It does not preserve old summand matrix conjugacy classes, produce a
prescribed nontrivial stabilizer action, preserve arbitrary numerical
depth labels, handle arbitrary transport graphs, settle the dyadic case,
or prove original same-candidate H5. It therefore does not close
unrestricted claim~66.

The retained primary BK text, SHA256
\path{d69db0827256707a34cd2562af68b9de11fc889ab3f43535ec47b620d58ca9c8},
was directly checked at:

\begin{itemize}
\tightlist
\item
  Definition~3.2.1 and Proposition 3.2.5, printed pp. 98--100: minimal
  character condition and explicit full-character extension.
\item
  Corollary 3.3.18, printed pp. 112--113: the actual additive correction takes
  the approximation character to the next simple character.
\item
  Condition 3.6.13(ii) and Theorem~3.6.14, printed pp. 136--137: transfer
  preserves the full beta-determinant character, including at m0=0.
\end{itemize}

The actual field and defining-sequence assertions are the separately
proved and independently reviewed trace-zero tower inputs, with BK 2.2.7
and the retained simple-stratum theorems. No newly downloaded source or
new novelty claim is involved.

\clearpage
\section{The surviving leading symplectic layer}\label{app:odd_capacity}

Status: ordinary proof from the retained Bushnell--Kutzko theorems and,
for the source-defined candidate, the substantial published
Bushnell--Henniart arithmetic induction/descent constructor. No Lean
build or bounded computation is used. The accompanying audit files identify the scope of the internal review.

\hypertarget{odd_capacity-the-useful-simplification}{%
\subsection{The useful simplification}\label{odd_capacity-the-useful-simplification}}

Full critical-profile rigidity is unnecessary for the trivial-stabilizer
finite attachment contract. An even leading level produces a symplectic
subspace that survives every later simple approximation. For the
independent rank-p Heisenberg tensor sources, this single subspace
already has enough dimension to hold every source plane, however many
there are.

The argument below concerns the actual type in the source-defined
arithmetic candidate. It does not identify that type with a separately
guessed binder parameter.

\hypertarget{odd_capacity-leading-layer-theorem}{%
\subsection{1. Leading-layer theorem}\label{odd_capacity-leading-layer-theorem}}

Let F have residue field k of cardinality q and residue characteristic
p.~Let pi be an actual irreducible supercuspidal representation of
GL\_n(F) with torsion number one, where n=p\^{}r. Suppose

\begin{Verbatim}[fontsize=\small,breaklines=true,breakanywhere=true]
a(pi)=n+N,     N>0 even,     v_p(N)=r−1.
\end{Verbatim}

Then an actual maximal simple type in pi, with Heisenberg quotient
\(W=J^1/H^1\), contains a nondegenerate symplectic
F\_p-subspace S of dimension

\begin{Verbatim}[fontsize=\small,breaklines=true,breakanywhere=true]
[k:F_p] (p−1)p^(r−1).
\end{Verbatim}

For k=F\_p this is exactly (p−1)p\^{}(r−1). The subspace is the image of
the concrete congruence group U\^{}(N/2)(A) in \(J^1/H^1\).
It is present regardless of the later critical exponents, their
parities, or the number of intermediate approximation fields.

\hypertarget{odd_capacity-actual-parameter-field-and-leading-level}{%
\subsubsection{1.1. Actual parameter field and leading
level}\label{odd_capacity-actual-parameter-field-and-leading-level}}

By BK maximal-type existence, pi contains a maximal simple type. Its
positive level follows from a(pi)\textgreater n.~Write its stratum as
{[}A,N\_A,0,beta{]}, with parameter field K=F{[}beta{]}. BK 6.2.5 gives

\begin{Verbatim}[fontsize=\small,breaklines=true,breakanywhere=true]
t(pi)=n/e(K/F).
\end{Verbatim}

Thus e(K/F)=n.~Since {[}K:F{]} divides n, K/F is totally ramified of
degree n and has residue degree one. The principal-order assertion after
BK 6.2.5 gives e(A)=n.~Accordingly A is the complete principal
hereditary order. The established conductor formula

\begin{Verbatim}[fontsize=\small,breaklines=true,breakanywhere=true]
a(pi)=n+n N_A/e(A)
\end{Verbatim}

now implies N\_A=N. This is the same source-checked formula used in
Block165, originating in BH1999 Section 6.1 Lemma 2 and retained in
\path{PRINCIPAL_ORDER_CONDUCTOR161.md}.

Put P=rad(A), \(U^a(A)=1+P^a\) for a\textgreater=1, and
m=N/2.

\hypertarget{odd_capacity-the-terminal-minimal-approximation-has-degree-p}{%
\subsubsection{1.2. The terminal minimal approximation has degree
p}\label{odd_capacity-the-terminal-minimal-approximation-has-degree-p}}

Apply the defining-sequence construction BK 2.4.1--2 to beta. Write
gamma for its terminal minimal approximation. Every transition has a
critical level c\textless N, and gamma has the same A-level N as beta.
All approximation residue degrees divide f(K/F)=1, and all approximation
ramification degrees divide n.~Therefore E=F{[}gamma{]} is totally
ramified of degree d dividing n.

The scalar case is impossible: a scalar of F has A-valuation divisible
by n, whereas v\_p(N)=r−1. Thus gamma is noncentral and minimal in the
ordinary field sense. Its field valuation is

\begin{Verbatim}[fontsize=\small,breaklines=true,breakanywhere=true]
−v_E(gamma)=Nd/n.
\end{Verbatim}

Minimality, by BK 1.4.14--15, requires gcd(Nd/n,d)=1. Since d is a
p-power and v\_p(N)=r−1, integrality first gives p\textbar d, while
coprimality forces d=p.~Hence E/F is totally ramified of degree p.

Set B=End\_E(F\^{}n), B\_order=A∩B, and Q=rad(B\_order). The same
complete lattice chain has period n/p over E. Each P-grade has
k-dimension n and each Q-grade has k-dimension n/p.~In particular,

\begin{Verbatim}[fontsize=\small,breaklines=true,breakanywhere=true]
dim_k P^m/(Q^m+P^(m+1))=n−n/p.
\end{Verbatim}

No assertion that E is literally a subfield of the final field K is
needed. It is the actual minimal approximation field in the defining
sequence.

\hypertarget{odd_capacity-the-high-subgroup-survives-every-approximation}{%
\subsubsection{1.3. The high subgroup survives every
approximation}\label{odd_capacity-the-high-subgroup-survives-every-approximation}}

The minimal definitions BK 3.1.7--8 give

\begin{Verbatim}[fontsize=\small,breaklines=true,breakanywhere=true]
J_order(gamma)=B_order+P^m,
H_order(gamma)=B_order+P^(m+1).
\end{Verbatim}

At a nonminimal transition beta\_i→beta\_(i+1), with critical level
c\_i\textless N, the recursive J-order definition uses the threshold
ceil(c\_i/2). Since ceil(c\_i/2)\textless=m, the inclusion
P\^{}m⊂J\_order(beta\_(i+1)) implies P\^{}m⊂J\_order(beta\_i). Starting
at gamma, induction proves

\begin{Verbatim}[fontsize=\small,breaklines=true,breakanywhere=true]
U^m(A)⊂J^1(beta).
\end{Verbatim}

For H, BK 3.1.9(ii) states that
\(H^a(\beta_i)=H^a(\beta_{i+1})\) as soon as

\begin{Verbatim}[fontsize=\small,breaklines=true,breakanywhere=true]
a>=floor(c_i/2)+1.
\end{Verbatim}

Every c\_i is at most 2m−1, so this threshold is at most m. Consequently

\begin{Verbatim}[fontsize=\small,breaklines=true,breakanywhere=true]
H^m(beta)=H^m(gamma).
\end{Verbatim}

Here H\^{}m means \(H^1\cap U^m(A)\). Therefore the image S of
U\^{}m(A) in W is canonically

\begin{Verbatim}[fontsize=\small,breaklines=true,breakanywhere=true]
S = U^m(A)/H^m(beta)
  = U^m(A)/H^m(gamma)
  ≅ P^m/(Q^m+P^(m+1)).
\end{Verbatim}

The last group identification uses m\textgreater=1, so multiplication is
additive modulo P\^{}(m+1). Thus \textbar S\textbar=q\^{}(n−n/p).

\hypertarget{odd_capacity-its-actual-commutator-form-is-nondegenerate}{%
\subsubsection{1.4. Its actual commutator form is
nondegenerate}\label{odd_capacity-its-actual-commutator-form-is-nondegenerate}}

Choose the actual simple character theta contained in the type of pi. By
the standard deep-restriction property of simple characters, theta
restricts to psi\_beta on U\^{}(m+1)(A). This follows directly from BK
3.2.1--4 by induction along the defining sequence.

For u,v in U\^{}m(A), their group commutator is in
\(U^{2m}(A)=U^N(A)\), which lies in U\^{}(m+1)(A). The
terminal approximation satisfies

\begin{Verbatim}[fontsize=\small,breaklines=true,breakanywhere=true]
beta−gamma∈P^(1−N),
\end{Verbatim}

because every critical level in the sequence is less than N. Hence

\begin{Verbatim}[fontsize=\small,breaklines=true,breakanywhere=true]
(beta−gamma)([u,v]−1)∈P.
\end{Verbatim}

With the standard level-one additive character psi\_A=psi\_F∘tr\_A/F,
psi\_A is trivial on P. We therefore have the exact equality

\begin{Verbatim}[fontsize=\small,breaklines=true,breakanywhere=true]
theta([u,v]) = psi_beta([u,v]) = psi_gamma([u,v]).
\end{Verbatim}

The right side is the Heisenberg commutator pairing for the minimal
simple stratum of gamma. Its quotient
\(J^1(\gamma)/H^1(\gamma)\) is precisely
\(U^m(A)/H^m(\gamma)\), because
\(J^1(\gamma)=U^1(B_{\mathrm{order}})U^m(A)\), with
\(U^1(B_{\mathrm{order}})\subset H^1(\gamma)\). BK's nondegeneracy theorem for
simple-character commutators therefore says that the form on S is
nondegenerate.

For general k, this is a nondegenerate alternating F\_p-form of
dimension \([k:\FF_p](n-n/p)\). For \(k=\FF_p\), it has dimension \(n-n/p\). This
proves the theorem.

The argument uses no self-pair conductor and no determination of the
later critical levels. Those invariants remain useful for finer
arithmetic identification, but they are not needed for this
finite-sector capacity claim.

\hypertarget{odd_capacity-arbitrary-length-odd-prime-source-family}{%
\subsection{2. Arbitrary-length odd-prime source
family}\label{odd_capacity-arbitrary-length-odd-prime-source-family}}

Fix an odd prime p and F=F\_p((t)). Let the complete marked Weil source
rho be an independent tensor product of r rank-p wild Heisenberg factors
at distinct positive levels

\begin{Verbatim}[fontsize=\small,breaklines=true,breakanywhere=true]
0<h_1<...<h_r,     p does not divide h_i.
\end{Verbatim}

Here independence means the stipulated balanced-central-product
description of the wild image is valid, and rho\textbar P\_F is
irreducible. We do not infer independence from the inequalities alone.
Suppose the retained Tsushima source calculation supplies the usual
highest slope

\begin{Verbatim}[fontsize=\small,breaklines=true,breakanywhere=true]
slope(rho)=(p+1)h_r/p.
\end{Verbatim}

Then rank(rho)=n=p\^{}r and

\begin{Verbatim}[fontsize=\small,breaklines=true,breakanywhere=true]
Swan(rho)=N=(p+1)p^(r−1)h_r.
\end{Verbatim}

Let Pi\_BH(rho) be the published source-defined arithmetic
induction/descent candidate, with its full prescribed central character,
as audited in Block165. The construction is applied to the complete
actual rho, including its scalar normalization. Its codomain is the
torsion-one supercuspidal class, and its standard epsilon-function
property gives a(Pi\_BH(rho))=n+N. These are uses of the substantial
classical constructor, not consequences of finite Heisenberg algebra
alone and not a selection of an inverse LLC image.

For odd p, N is even and v\_p(N)=r−1. The leading-layer theorem
therefore puts a symplectic space of dimension

\begin{Verbatim}[fontsize=\small,breaklines=true,breakanywhere=true]
D=(p−1)p^(r−1)
\end{Verbatim}

inside the actual Heisenberg quotient of a type in this same
Pi\_BH(rho).

The source packet is the orthogonal sum of r planes, so its symplectic
dimension is 2r. The elementary inequality

\begin{Verbatim}[fontsize=\small,breaklines=true,breakanywhere=true]
(p−1)p^(r−1)>=2r
\end{Verbatim}

holds for every odd p and r\textgreater=1: it suffices that
3\^{}(r−1)\textgreater=r, which follows by induction. Thus the actual
candidate's leading symplectic layer alone can receive the entire source
packet.

\hypertarget{odd_capacity-full-characters-and-recovery-for-declared-c1}{%
\subsubsection{2.1. Full characters and recovery for declared
C=\{1\}}\label{odd_capacity-full-characters-and-recovery-for-declared-c1}}

Work under the same finite scalar-pushout contract as
\path{SAME_CANDIDATE_ATTACHMENT165.md}: declared stabilizer C=\{1\}, one
vertex, and full source character extensions retained, with no
requirement of literal old matrix-group inclusion or equality of old
numerical depths with host depths.

Let P\_src be the balanced C×-central product of the r supplied full
source character extensions. Its quotient V\_src has dimension 2r over
F\_p and its central-character-one Heisenberg representation has
dimension p\^{}r=n.~Choose a symplectic embedding V\_src→S⊂W.

The C×-pushout of the actual type character, restricted to this
subspace, has the same alternating commutator as P\_src. The retained
extension theorem supplies a center-fixing isomorphism from P\_src to
that restriction. Consequently every source h∈H\_i\^{}1 is carried to
the scalar theta\_i(h), preserving the full character value. No shallow
source values are discarded.

Let V be the image and write W=V⊥V\_perp. Stone--von Neumann gives

\begin{Verbatim}[fontsize=\small,breaklines=true,breakanywhere=true]
eta_W ≅ eta_src⊗eta_perp.
\end{Verbatim}

Choose a line in eta\_perp and the corresponding rank-n projector e in
the actual finite Heisenberg image algebra. Then e eta\_W≅eta\_src.
Applying the retained Block159 induction/idempotent theorem to the
unchanged actual extended type Lambda in Pi\_BH(rho) recovers Lambda and
then this exact Pi\_BH(rho). All source normalization and the actual
noncompact scalar remain fixed.

This proves same-candidate finite attachment for arbitrary r in this
C=\{1\} odd-prime family. Original H5, when supplied for this same
candidate, remains exactly the original conditional comparison premise.
We do not prove it here.

\hypertarget{odd_capacity-first-new-rank-gl_27}{%
\subsubsection{2.2. First new rank: GL\_27}\label{odd_capacity-first-new-rank-gl_27}}

For p=3,r=3, the source has rank~27 and V\_src has dimension 6. Put h=h\_3
with 3∤h. The actual candidate has

\begin{Verbatim}[fontsize=\small,breaklines=true,breakanywhere=true]
N=36h,       m=18h,
dim_F3 S=27−9=18.
\end{Verbatim}

Thus its high layer contains the entire three-plane source packet. At
least a 3\^{}9-dimensional Heisenberg contribution is present, while the
required source sector has dimension 27=3\^{}3. The six surplus
symplectic pairs can be projected onto a line in their representation,
together with any additional lower layers. This proves the required
capacity and sector realization for the new rank~27 family without
solving its two possible lower critical levels.

In particular, for (h\_1,h\_2,h\_3)=(1,2,4), N=144 and the same high
subspace has dimension 18. These numbers are a specialization of the
universal proof, not a numerical experiment.

\hypertarget{odd_capacity-why-trying-to-force-every-critical-level-was-unnecessary}{%
\subsection{3. Why trying to force every critical level was
unnecessary}\label{odd_capacity-why-trying-to-force-every-critical-level-was-unnecessary}}

At rank~27, a hypothetical complete degree pattern 27→9→3 has leading
level N=36h\_3 and lower critical levels a\textless b\textless N. The
standard full-profile conductor count would give

\begin{Verbatim}[fontsize=\small,breaklines=true,breakanywhere=true]
Swan(End rho)=18N+6b+2a.
\end{Verbatim}

Comparing with the stipulated source formula 8(h\_1+9h\_2+81h\_3) only
gives

\begin{Verbatim}[fontsize=\small,breaklines=true,breakanywhere=true]
a+3b=4h_1+36h_2.
\end{Verbatim}

That single linear relation does not determine both a and b, and by
itself does not force each parity. Nor does it establish that all three
degrees occur rather than skipping degree~9. This observation is only a
warning about the insufficient equations; it is not presented as an
arithmetic counterexample or as existence of every numerical profile
satisfying the equation.

The leading-layer theorem bypasses these unresolved profile choices. It
supplies a fixed nondegenerate high subspace in every actual candidate
with the two required invariants.

\hypertarget{odd_capacity-dyadic-capacity-boundary}{%
\subsection{4. Dyadic capacity
boundary}\label{odd_capacity-dyadic-capacity-boundary}}

The leading-layer theorem itself is valid at p=2. If an actual
torsion-one candidate of rank~2\^{}r has

\begin{Verbatim}[fontsize=\small,breaklines=true,breakanywhere=true]
N=3·2^(r−1)h,     h odd,     r>=2,
\end{Verbatim}

then its high subspace has dimension 2\^{}(r−1) over F\_2. This alone is
at least 2r for r\textgreater=4, with equality at r=4. It does not
supply enough capacity by itself for r=2 or 3. For r=1, N is odd and this
theorem does not apply.

This is a capacity statement only. Applying it to a dyadic source
requires the actual source family, its arithmetic candidate and whatever
quadratic/phase markings the original contract requires. No universal
dyadic attachment theorem is claimed from dimension alone.

\hypertarget{odd_capacity-boundaries-that-remain}{%
\subsection{5. Boundaries that remain}\label{odd_capacity-boundaries-that-remain}}

\begin{enumerate}
\def\labelenumi{\arabic{enumi}.}
\tightlist
\item
  The full character attachment here is through center-fixing
  scalar-pushout maps. Literal embeddings of prescribed old arithmetic
  matrix groups are not proved.
\item
  No nontrivial prescribed stabilizer action is attached. Capacity does
  not imply equivariant placement.
\item
  All planes may be placed in the leading host layer. Preservation of
  supplied numerical ramification depths, or a specified map into
  individual host depth layers, would be an additional requirement and
  is not claimed.
\item
  General packet transport graphs are not handled here. The statement is
  for C=\{1\} and one vertex.
\item
  The arithmetic induction/descent theorem is an explicit major external
  input. This is a new capacity/attachment consequence for its actual
  output, not a new elementary reconstruction of that classical
  constructor.
\item
  H5 remains conditional and refers to that unchanged actual candidate.
  No additional pair-factor premise is introduced.
\item
  No Lean verification is claimed.
\end{enumerate}

\hypertarget{odd_capacity-primary-locators-and-retained-dependencies}{%
\subsection{6. Primary locators and retained
dependencies}\label{odd_capacity-primary-locators-and-retained-dependencies}}

The authorized primary BK text was read at these exact results:
1.4.14--15 (minimality), 2.4.1--2 (defining sequence and degree
divisibility), 3.1.7--9 (recursive orders and high-depth equality),
3.2.1--4 (deep restriction of simple characters), 3.4.1 (nondegenerate
Heisenberg commutator), and 6.2.5 with the following paragraph (torsion
number and principal period).

The standard conductor and published arithmetic-constructor properties
are reused from the source-pinned Block165 audit and Block161 conductor
proof. The finite scalar-pushout and sector-recovery arguments are
exactly those retained and restated in
\path{SAME_CANDIDATE_ATTACHMENT165.md}. No entire copyrighted source is
copied into this continuation.

\clearpage
\section{Actual scalar inertia phases and finite-cover descent}\label{app:odd_scalar}

Evidence: ordinary proof using rank-one reciprocity, the given central
character, and elementary finite group arguments. The result constructs
actual central matrix operators through a finite cover and proves their
descent. It does not construct a new arithmetic candidate or alter a
canonical beta extension.

\hypertarget{odd_scalar-exact-statement}{%
\subsection{1. Exact statement}\label{odd_scalar-exact-statement}}

Let F be a nonarchimedean local field with residue field k of
cardinality q and residue characteristic p.~Put Q=q-1 and n=p\^{}r,
allowing r=0. Let rho:W\_F -\textgreater{} GL\_n(C) be a smooth Weil
representation with finite inertia image. Use the fixed reciprocity
convention identifying its determinant with a character omega:F×
-\textgreater{} C×.

Let C be a finite group of order prime to p, and suppose its prescribed
operators on the \textbf{full source carrier} are

\begin{Verbatim}[fontsize=\small,breaklines=true,breakanywhere=true]
Omega_src(c)=chi(c)I_n,
\end{Verbatim}

where chi is a character and each Omega\_src(c) belongs to the actual
matrix image rho(I\_F). A finite cyclic scalar tame image is the
intended example. This hypothesis refers to actual source operators, not
freely chosen phases on an unrelated multiplicity-space factor.

Let Pi be any already constructed irreducible smooth GL\_n(F)
representation with central character omega. Suppose it has its supplied
actual inducing representation Lambda on N, with F× contained in N, and
an attached finite sector E=eLambda realizing the prescribed source
Heisenberg packet. Assume the prescribed C action on the source scalar
extension is trivial, and that the honest scalar phase on the packet is
the same chi above.

\textbf{Theorem.} There is a finite group C\_tilde surjecting onto C and
a genuine homomorphism

\begin{Verbatim}[fontsize=\small,breaklines=true,breakanywhere=true]
g:C_tilde -> F× I_n contained in N
\end{Verbatim}

whose action on Lambda and on Pi descends to the prescribed honest C
action chi(c)I. Every finite sector is stable under this action. Any
center-preserving Heisenberg attachment and its intertwiner are
C-equivariant. The retained sector-induction recovery maps recover the
same Lambda and Pi equivariantly. All actual inducing data, including
the fixed canonical beta extension, remain unchanged.

The arithmetic action is supplied by actual central matrices, with an
explicit finite-cover descent. The theorem does not assert that C itself
embeds as a matrix subgroup. Such an embedding can fail even for a
genuine scalar tame source, as Section~4 shows.

\hypertarget{odd_scalar-construction-and-proof}{%
\subsection{2. Construction and proof}\label{odd_scalar-construction-and-proof}}

Let D=k×, and let Teich:D -\textgreater{} O\_F× be the canonical
multiplicative lift. For p=2 this is still the prime-to-p root-of-unity
subgroup. When q=2, D is trivial and the argument below forces chi
trivial as well.

Local reciprocity identifies the determinant image of inertia with

\begin{Verbatim}[fontsize=\small,breaklines=true,breakanywhere=true]
det rho(I_F)=omega(O_F×).
\end{Verbatim}

The unit group decomposes as

\begin{Verbatim}[fontsize=\small,breaklines=true,breakanywhere=true]
O_F× = Teich(D) times U_F¹.
\end{Verbatim}

The smooth character image omega(U\_F¹) is a finite p-group, while
omega(Teich(D)) has order prime to p.~Therefore the prime-to-p elements
of omega(O\_F×) are exactly omega(Teich(D)).

For every c in C, the actual-source hypothesis implies

\begin{Verbatim}[fontsize=\small,breaklines=true,breakanywhere=true]
chi(c)^n=det(Omega_src(c)) in det rho(I_F).
\end{Verbatim}

This value has order prime to p, and hence lies in omega(Teich(D)).
Since n is a p-power and chi(C) has order prime to p, the map z
-\textgreater{} z\^{}n permutes chi(C). Consequently

\begin{Verbatim}[fontsize=\small,breaklines=true,breakanywhere=true]
chi(C) is contained in omega(Teich(D)).                 (1)
\end{Verbatim}

In particular every source phase has order dividing Q. This conclusion
does not require selecting a splitting of the tame quotient or choosing
Frobenius representatives.

Choose an integer k\_0 with n k\_0 congruent to 1 modulo Q. When Q=1 take
k\_0=1. Define

\begin{Verbatim}[fontsize=\small,breaklines=true,breakanywhere=true]
delta(u)=omega(Teich(u))^(k_0),       u in D.
\end{Verbatim}

Because exponentiation by k\_0 is an automorphism on every group of
roots of unity of order dividing Q, delta(D)=omega(Teich(D)). Define the
finite fiber product

\begin{Verbatim}[fontsize=\small,breaklines=true,breakanywhere=true]
C_tilde={(c,u) in C times D : chi(c)=delta(u)}.          (2)
\end{Verbatim}

It is a subgroup, and projection pr:C\_tilde -\textgreater{} C is
surjective by (1). Define the actual central matrix homomorphism

\begin{Verbatim}[fontsize=\small,breaklines=true,breakanywhere=true]
g(c,u)=Teich(u^(k_0)) I_n.                             (3)
\end{Verbatim}

Its action on Lambda and Pi is the central-character scalar

\begin{Verbatim}[fontsize=\small,breaklines=true,breakanywhere=true]
Lambda(g(c,u))=omega(Teich(u^(k_0)))I
              =delta(u)I=chi(c)I.                     (4)
\end{Verbatim}

If (1,u) belongs to ker pr, (4) is the identity. Hence the represented
action factors uniquely through C. It is exactly the prescribed source
action. Although selecting a particular matrix representative for c may
require a choice of u, the represented operator is independent of that
choice. Changing k\_0 by a multiple of Q leaves (3) unchanged.

All matrices (3) are central. Their conjugation on actual groups,
characters, and scalar pushouts is therefore the identity, and their
Lambda operators are scalar. Thus they preserve every e and every sector
E. The source extension action is also the identity by hypothesis, so
any attachment of those scalar extensions is equivariant. Both honest
carrier actions are chi(c)I, so every Heisenberg intertwiner intertwines
C without a new phase choice.

For the recovery step, put P=\{a in N:Lambda(a)eLambda(a)\^{}(-1)=e\}.
Then F× is contained in P. The retained maps

\begin{Verbatim}[fontsize=\small,breaklines=true,breakanywhere=true]
Phi:Ind_P^N E -> Lambda,
j:Lambda -> Ind_P^N E
\end{Verbatim}

are N-equivariant. They therefore commute with all actual matrices (3),
and ker pr also acts trivially on Ind\_P\^{}N E: its matrices are
central elements of P acting trivially on E, hence trivially on every
induced translate. Thus the equality descends through pr to C. The
recovering idempotent, and then compact induction to GL\_n(F), preserve
this equivariance. They recover the same fixed Lambda and Pi. QED.

\hypertarget{odd_scalar-what-this-closes-and-what-it-does-not}{%
\subsection{3. What this closes and what it does
not}\label{odd_scalar-what-this-closes-and-what-it-does-not}}

For a source packet whose marked tensor carrier is the entire rank-n
source, this proves the stabilizer step whenever its genuine tame
operators are scalar and an actual host with the correct central
character and full-character attachment is already present. There is no
remaining scalar-phase equation to solve in that situation: the
determinant and central character supply it.

If only the source symplectic C action is known to be trivial, the
preceding \path{STABILIZER_CONTRACT166.md} proves triviality of its
full scalar-extension action for prime-to-p C. Schur then gives a scalar
packet character. To apply the present theorem one must still identify
that character with the action of actual inertia on the full source.
That identification holds when the supplied honest packet carrier is the
whole source carrier. It is not automatic for a proper
multiplicity-space factor, whose omitted factor can contribute to the
determinant.

For several factors the total packet character is the product of the
supplied factor characters. The theorem matches that product on the
whole attached packet. It does not represent independent factor actions
by distinct central matrices unless those additional marked actions are
separately specified and verified.

The result uses the fixed full inducing representation Lambda. If
Lambda\textbar J=kappa\_can tensor sigma, then kappa\_can and sigma are
both unchanged. Equation (4) concerns Lambda, and does not assert that
kappa\_can alone has the same scalar character. A contract prescribing
an action specifically on the canonical beta extension requires that
distinction to remain explicit.

No nontrivial source symplectic action is realized here. No new base
candidate, new H5 comparison, dyadic quadratic refinement, or general
rank theorem is claimed. Existing quadratic refinements are preserved
because the actual central conjugation action is the identity.

\hypertarget{odd_scalar-why-a-direct-map-c-into-residue-units-is-too-strong}{%
\subsection{4. Why a direct map C into residue units is too
strong}\label{odd_scalar-why-a-direct-map-c-into-residue-units-is-too-strong}}

Take q=5, n=5\^{}r, and let omega on k×=F5× be the quadratic character.
Let C=C2 and let chi be its nontrivial character. Such scalar tame
behavior is obtained by twisting a tame-trivial rank-5\^{}r source by
the quadratic tame character, whenever such a source is supplied; the
retained odd-prime two-factor family supplies rank~25 examples.

A group homomorphism a:C2 -\textgreater{} F5× must send its generator
to 1 or -1. Both have quadratic character 1. Consequently

\begin{Verbatim}[fontsize=\small,breaklines=true,breakanywhere=true]
omega(a(c)) = chi(c)^n
\end{Verbatim}

is impossible at the nontrivial c, since n is odd. The proposed direct
section is not forced by genuine scalar tame source data. Reciprocity
supplies a quotient of residue units; that quotient need not split.

The construction above still works. Here k\_0=1 modulo 4, C\_tilde is the
cyclic group of order 4, and its central matrices act by the quadratic
character. The kernel of C\_tilde -\textgreater{} C2 acts trivially on
the fixed candidate, so the actual operator action descends correctly.
This example explains why finite-cover descent is necessary rather than
a cosmetic reformulation.

\hypertarget{odd_scalar-dependencies-and-review-status}{%
\subsection{Dependencies and review
status}\label{odd_scalar-dependencies-and-review-status}}

External input: rank-one local class field theory, with the same fixed
normalization used for omega=det rho; the existing actual candidate's
central character identity; the retained sector recovery theorem. The
finite-cover construction and its descent are proved above and use no
local Langlands candidate selection. No computation is needed. This
version has undergone author self-review; a parent audit is requested.
No Lean run is claimed.

\clearpage
\section{The precise stabilizer-action contract and normalizer obstruction}\label{app:stabilizer_contract}

Evidence: direct audit of the original statement, a new elementary
reduction, and an explicit actual-source/actual-type H1--H2 witness.
No Lean closure, class-wide arithmetic
realization, or H5 witness is claimed.

\textbf{Primary audit verdict:} incomplete, with the smallest missing
implication being the identification of the declared finite action with
the action intended on a newly constructed common arithmetic candidate.
The original text does not define that identification by a common matrix
normalizer.

\hypertarget{stabilizer_contract-exact-contract-and-what-is-already-an-input}{%
\subsection{1. Exact contract, and what is already an
input}\label{stabilizer_contract-exact-contract-and-what-is-already-an-input}}

The controlling text is
\path{closure_attempt162/inputs/original_reading.txt}:

\begin{itemize}
\tightlist
\item
  Proposition 9.8, starting line 4834: a finite group acting by
  automorphisms of the one-break Clifford datum acts symplectically on
  its intrinsic envelope.
\item
  Definition~11.1, lines 4911--4923: a bridge supplies an equivariant
  isometry to an actual type quotient for a finite cyclic stabilizer C,
  with common central character and additional dyadic lift data.
\item
  Proposition 11.3 and Theorem~11.5: once the two honest cyclic actions
  are chosen, their discrepancy is a character, killed by
  generator-phase equality.
\item
  Theorem~15.4 and Theorem~16.2: breakwise comparisons and transports are
  supplied and tensored. They do not construct a common arithmetic host
  from separate actual type targets.
\end{itemize}

Thus H2 and H3 already settle the finite action comparison \textbf{on
each supplied break target}. They do not leave a second free stabilizer
action to invent on that same target. Their transport into a new common
arithmetic type is another assertion. Neither a map C into its matrix
normalizer nor arbitrary transported operators are identified with that
assertion by the original text.

The cyclic group in Definition~11.1 is not explicitly required there to
have order prime to p.~The reduction below assumes that order condition;
it applies to the declared tame prime-to-p case, and cannot silently be
applied to an arbitrary finite-order Frobenius operator.

\hypertarget{stabilizer_contract-full-scalar-extension-rigidity-for-a-native-tame-action}{%
\subsection{2. Full scalar-extension rigidity for a native tame
action}\label{stabilizer_contract-full-scalar-extension-rigidity-for-a-native-tame-action}}

Let an actual maximal simple type in GL\_n(F) have a totally ramified
parameter field K/F of degree n=p\^{}r and residue matrix rank one. Let

\begin{Verbatim}[fontsize=\small,breaklines=true,breakanywhere=true]
W=J¹/H¹,  P=(C××J¹)/{(theta(h)^(-1),h):h in H¹},
\end{Verbatim}

and let eta be its irreducible Heisenberg representation. Let C be
finite of order prime to p.~Suppose the C action on P fixes the scalar
center and is genuinely induced by character-preserving matrix
normalizers of this type. It is enough that it is an honest action whose
action on W comes from those normalizers. Then:

\begin{enumerate}
\def\labelenumi{\arabic{enumi}.}
\tightlist
\item
  C acts trivially on W.
\item
  C acts trivially on the full scalar extension P.
\item
  Every honest C action Omega on eta implementing this action is
  Omega(c)=chi(c)I for one character chi of C.
\end{enumerate}

\textbf{Proof.} The retained normalizer theorem in Block162 gives

\begin{Verbatim}[fontsize=\small,breaklines=true,breakanywhere=true]
N=K×J¹,          N/(F×J¹) = Z/nZ.
\end{Verbatim}

The action on W factors through this cyclic p-group, so a homomorphic
image of the prime-to-p group C is trivial.

For any central extension P of the elementary abelian p-group W by C×,
the kernel of

\begin{Verbatim}[fontsize=\small,breaklines=true,breakanywhere=true]
Aut_(C×)(P) -> Aut(W)
\end{Verbatim}

is naturally Hom(W,C×). Indeed, an automorphism phi fixing the quotient
and the center has phi(x)=a(x\_bar)x. The scalar a is independent of the
lift x and multiplication shows that a is a character. Conversely every
such character defines an automorphism. Since W has exponent p, this
kernel is a finite elementary abelian p-group. The action of C lies in
this kernel and therefore is trivial. Finally an implementing operator
commutes with the irreducible P action, so is scalar by Schur. Honesty
makes its scalar a character. QED.

This strengthens the earlier W-level obstruction: passing to the full
scalar pushout cannot hide a nontrivial native prime-to-p action above
an identity quotient action.

\hypertarget{stabilizer_contract-exact-attachment-criterion-in-this-native-setting}{%
\subsection{3. Exact attachment criterion in this native
setting}\label{stabilizer_contract-exact-attachment-criterion-in-this-native-setting}}

Suppose a marked source scalar extension P\_src has a symplectic
embedding into W, and an actual finite sector of eta realizes that
extension. Fix its center-preserving attachment and a Heisenberg
intertwiner T. If the host has the native tame action of Section~2, the
attachment intertwines the full honest C actions if and only if

\begin{Verbatim}[fontsize=\small,breaklines=true,breakanywhere=true]
the source action on P_src is trivial, and
chi_src = chi_host.
\end{Verbatim}

Necessity follows from equivariance and Section~2. For sufficiency,
trivial extension actions make both honest carrier actions scalar by
Schur, and equal characters commute through every Heisenberg
intertwiner. Any sector is stable under the scalar host action.

If the source C action on its symplectic quotient is trivial, its action
on the full scalar extension is already trivial by the Hom-kernel
argument. For several break factors, the source character is the product
of their characters. Thus the criterion is a single scalar-character
equality, rather than a further system of symplectic placement
equations.

This applies to the original C=\{1\} fixture immediately. For nontrivial
C, H3 equates source and \textbf{supplied break-target} phases. It does
not equate their product with a separately built common host's phase
without a dictionary. We do not change a fixed canonical beta extension
to enforce this equality.

\hypertarget{stabilizer_contract-h2-and-matching-carrier-dimensions-do-not-force-trivial-c}{%
\subsection{4. H2 and matching carrier dimensions do not force trivial
C}\label{stabilizer_contract-h2-and-matching-carrier-dimensions-do-not-force-trivial-c}}

Here is an actual H1--H2 witness under the literal original allowance of
separate type targets. It is not a claimed H5-admitted input, and does
not assert equal ambient matrix ranks for source and break target.

Take F=Fp((t)), p odd. The retained primary source by Tsushima,
Definition~3.1, Lemmas 2.6, 2.8, 3.9, 3.12, 3.14 and Corollary 3.15, with
R(X)=X\^{}p and m=1, constructs a genuine rank-p Weil representation
irreducible on wild inertia. Its wild image is extraspecial of order p³.
The source envelope is a symplectic plane, and its actual tame image
C=mu\_(p+1) acts by multiplication on the roots of X\^{}(p²)+X. This
action is nontrivial and has order p+1. The source's finite Heisenberg
carrier has dimension p.

For the actual type target, take the unramified quadratic extension E/F
and the maximal order A=End\_OF(OE). Let beta=t\^{}(-2)a, where a is a
unit generating the residue extension. This is a minimal simple stratum
of even level two. The standard minimal simple-character formulas give

\begin{Verbatim}[fontsize=\small,breaklines=true,breakanywhere=true]
J¹=U_A¹,  H¹=U_E¹ U_A²,
W_theta=End_Fp(Fp²)/Fp²,     dim_Fp W_theta=2,
dim eta_theta=p.
\end{Verbatim}

These are actual groups of a simple character, not an abstract
replacement. The type residue units act by conjugation. Write
sigma(x)=x\^{}p and decompose

\begin{Verbatim}[fontsize=\small,breaklines=true,breakanywhere=true]
End_Fp(Fp²)=Fp² direct-sum Fp² sigma.
\end{Verbatim}

On the quotient represented by b sigma, conjugation by u in Fp²× is

\begin{Verbatim}[fontsize=\small,breaklines=true,breakanywhere=true]
b -> u^(1-p)b.
\end{Verbatim}

Its image is the norm-one torus, of order p+1. Consequently the actual
projective residue action of Fp²×/Fp× on the two-dimensional quotient is
nontrivial. It has the same symplectic C-module as the source: identify
the cyclic groups using the displayed ratio map. The two invariant
alternating forms differ by a nonzero Fp scalar, and multiplication by
lambda in Fp²× rescales an alternating form by Norm(lambda).
Surjectivity of the norm gives an equivariant symplectic isometry.

This supplies the actual source, actual type quotient, matching degree-p
Heisenberg carriers, and nontrivial equivariant H2 action. For p=3, C
has order 4 and the action is an irreducible rotation of the F3 plane.

The target here is in GL\_2(F), while the source has rank p.~This is
permitted by the literal breakwise definitions, which contain no
relation between the individual type targets' ambient ranks. Block161
already recorded that exact omission. If a stricter intended contract
requires every break target in a specified GL\_(p\^{}r)(F), this example
does not settle that stricter question. We have not supplied H5 for a
common candidate or certified fixed canonical beta-extension phases for
this example. It refutes only the proposed inference that actual H2 plus
equality of finite carrier dimensions forces C trivial.

\hypertarget{stabilizer_contract-result-for-the-next-construction-attempt}{%
\subsection{5. Result for the next construction
attempt}\label{stabilizer_contract-result-for-the-next-construction-attempt}}

There are three distinct cases:

\begin{longtable}[]{@{}
  >{\raggedright\arraybackslash}p{(\columnwidth - 2\tabcolsep) * \real{0.5000}}
  >{\raggedright\arraybackslash}p{(\columnwidth - 2\tabcolsep) * \real{0.5000}}@{}}
\toprule\noalign{}
\begin{minipage}[b]{\linewidth}\raggedright
Given contract
\end{minipage} & \begin{minipage}[b]{\linewidth}\raggedright
Exact remaining issue
\end{minipage} \\
\midrule\noalign{}
\endhead
\bottomrule\noalign{}
\endlastfoot
Original C=\{1\} fixture & No stabilizer issue remains. \\
Native prime-to-p C action in a degree-p\^{}r totally ramified host &
Source extension action must be trivial; then compare one scalar
character. \\
Nontrivial source symplectic C action under the original finite contract
& A common-host arithmetic action dictionary is needed. Requiring an
inner normalizer makes this impossible for the preceding host;
transporting external operators alone was already done in Block 161. \\
\end{longtable}

Thus ``solve arbitrary nontrivial stabilizer attachment by adjusting the
current totally ramified binder'' is not a viable native-inner target.
That is a precise failure of that stronger construction route, not a
refutation of 16.2.

\hypertarget{stabilizer_contract-evidence-and-self-review}{%
\subsection{Evidence and self-review}\label{stabilizer_contract-evidence-and-self-review}}

Dependencies: original Sections~9--11,15--16; Block162's source-backed
normalizer theorem; Block165's base realization contract and
same-candidate attachment; the retained published Tsushima text; the
retained minimal-stratum group formulas. The new Hom-kernel lemma and
attachment criterion are ordinary elementary proofs. The witness uses
the displayed arithmetic constructor inputs, not a numerical experiment.
No independent review of this new note is claimed in this file, and no
new Lean compilation was attempted.

Self-review explicitly checked: the prime-to-p qualification, honest
action versus projective implementers, center-fixed scalar extension
automorphisms, separate target rank versus original rank, and the
absence of an H5 witness.

\clearpage
\section{Dyadic synthesis and all-length assembly scope}\label{app:dyadic_synthesis}

The odd-prime construction used here is pinned in
\path{TARGET167.md}. Evidence: ordinary mathematical proof and
independent component reviews. No Lean verification or universal claim~66
closure is asserted.

\hypertarget{dyadic_synthesis-combined-arithmetic-theorem}{%
\subsection{1. Combined arithmetic
theorem}\label{dyadic_synthesis-combined-arithmetic-theorem}}

Let \(F=\mathbb F_2((t))\), fix an additive character \(\psi_F\) of
conductor \(t\mathcal O_F\), and take increasing odd positive integers
\(h_1<\cdots<h_d\). Put \[
q_i=3h_i,\qquad N=2^d,\qquad n_j=2^{d-j}q_{d-j+1}.
\] There is an explicit \(\beta\in M_N(F)\) generating a separable
totally ramified degree-\(N\) field that normalizes the complete lattice
chain. It gives a maximal simple stratum with leading level \(n_1\),
final critical exponent \(-q_1\), and successive approximation-field
degrees \(2,4,\ldots,2^d\), with critical levels \(n_2,\ldots,n_d\).

The actual stratum admits a full simple character \(\theta\) trivial on
\(U_F^1\). Its canonical beta extension can be completed to an actual
irreducible admissible supercuspidal \(\pi\) of \(\mathrm{GL}_N(F)\)
with any prescribed smooth central quasicharacter \(\omega\) trivial on
\(U_F^1\), provided the noncompact completion is free or already
compatible. Over this residue field, such \(\omega\) is unramified.

For this same constructed candidate, \[
\begin{aligned}
t(\pi)&=1,& \operatorname{Sw}(\pi)&=3\,2^{d-1}h_d,\\
\dim_{\mathbb F_2}W&=N-2,&
\dim\eta&=2^{N/2-1},\qquad W=J^1/H^1,\\
a(\pi^\vee\times\pi)
 &=N^2-1+\sum_{i=1}^d3\,4^{i-1}h_i.
\end{aligned}
\] Matching these profiles does not identify the source, its endo-class,
or the pair factors in original H5.

\hypertarget{dyadic_synthesis-the-explicit-construction-and-its-depth-invariant}{%
\subsection{2. The explicit construction and its depth
invariant}\label{dyadic_synthesis-the-explicit-construction-and-its-depth-invariant}}

At a recursive coefficient field \(L\) of characteristic two, with
residue field \(\mathbb F_2\) and uniformizer \(t_L\), put \[
E=L(\alpha),\qquad \alpha^2-\alpha=t_L^{-q},\qquad
a=(q+1)/2,\qquad z=t_L^a\alpha.
\] The extension is separable and totally ramified of degree two; \(z\)
is a uniformizer and \[
z^2+t_L^az+t_L=0,\qquad \operatorname{Tr}_{E/L}(\alpha)=1.
\] For constant-coefficient projection \(P_0\) in the basis \(1,z\), the
normalized trace corestriction and exact filtered right inverse are \[
\kappa=\alpha^{-1},\qquad s=\kappa\operatorname{pr},\qquad
s(P_0)=t_L/z^2,\qquad
\operatorname{Lift}_E(C)=(z^2/t_L)CP_0.
\] They satisfy \(s(\operatorname{Lift}_E(C))=C\), preserve the exact
filtration valuation, and give the identity \[
\operatorname{tr}_L\operatorname{Lift}_E(C)
 =\operatorname{Tr}_{E/L}
   \bigl(\alpha\operatorname{tr}_E C\bigr).
\] The actual recursive matrix is \[
B(L,t_L;h_1,\ldots,h_d)
 =\alpha I_{2^{d-1}}+
   \operatorname{Lift}_E\bigl(B(E,z;h_1,\ldots,h_{d-1})\bigr),
\] with multiplication by \(\alpha\) as the rank-two base case.

Every native constructor has trace one. At every next native rank, the
preceding parameter amplified by two has trace zero, so the native
correction also has trace one. At full ambient rank, the earlier
amplified approximations instead have trace zero and the final one has
trace one. These are different statements; no amplified trace is used as
a substitute for the native trace needed in character construction.

At an inductive field step \(d\ge2\), set \(n=2^{d-1}q_d\),
\(R=2^{d-2}q_{d-1}\). Then \(n>2R\) and \(-2R>-n-q_1\). BK Remark 2.2.7
gives the maximal field generated by the actual displayed matrix.
Literal truncation and the filtered lift give the defining sequence.
Nonzero native trace also proves separability: the irreducible
characteristic polynomial has a nonzero coefficient of \(X^{N-1}\), so
its derivative is nonzero. See \path{tower/DYADIC_AS_TOWER167.md}
for all field and filtration details.

The replacement for trace zero is the depth inequality \[
\boxed{v_F(\operatorname{tr}\delta)+s\ge1.}
\] It makes the additive correction trivial on scalar \(U_F^s\), because
\((u-1)\operatorname{tr}\delta\in t\mathcal O_F\). Our native
corrections have trace one, so the inequality holds at \(s=1\). This is
the exact source-to-target transfer of the weighted-visibility
principle; the Euler coefficient theorem is motivation, not an input to
the arithmetic proof. In characteristic two, the weights are uniformizer
depths, rather than integer powers of two in the coefficient field.

\hypertarget{dyadic_synthesis-full-characters-including-the-odd-critical-endpoint}{%
\subsection{3. Full characters, including the odd critical
endpoint}\label{dyadic_synthesis-full-characters-including-the-odd-critical-endpoint}}

At the minimal rank-two base, extend the forced character on
\(U_E^{(q+1)/2}\) to \(U_E^1\), imposing triviality on \(U_F^1\). The
prescriptions agree on their intersection because the parameter trace is
one and \(\psi_F(u-1)=1\). This is an actual finite abelian character
extension and gives the full minimal simple character.

For a later native step, first transfer the full beta-determinant unit
character to multiplicity two. BK 3.6.13(ii) and 3.6.14 give, on scalar
units, \[
\Theta_j(uI)=\theta_j(u^2)=1.
\] For odd critical level \(r\), use \(m=\lfloor r/2\rfloor\) and
\(D=H^{m+1}(\beta)=H^{m+1}(\gamma)\). The endpoint is valid because
\(2(m+1)-r=1\): the additive-character cross term lies in the radical on
which the matrix additive character is trivial. BK 3.3.18 supplies
\(\vartheta=\Theta_j|_D\,\psi_{\beta-\gamma}\). Integral trace makes it
trivial on \(D\cap U_F^1\).

The next parameter field is maximal. Extend the forced character on its
deep units together with the trivial character of \(U_F^1\), then apply
BK 3.2.5 to obtain the next full character. All new extensions occur at
maximal native rank. No unverified determinant-image extension in a
nonmaximal centralizer is used.

For the final totally ramified field, \(J/J^1=\mathbb F_2^\times=1\), so
\(J=J^1\) and the canonical beta extension is \(\eta\). Extend to
\(K^\times J\), and choose the noncompact scalar to match \(\omega(t)\).
The needed \(N\)-th root exists in \(\mathbb C^\times\). The compact
type stays fixed. All choices precede H5. See
\path{central/DYADIC_FULL_CHARACTER167.md} for the complete proof.

\hypertarget{dyadic_synthesis-exact-packet-capacity-and-marked-attachment}{%
\subsection{4. Exact packet capacity and marked
attachment}\label{dyadic_synthesis-exact-packet-capacity-and-marked-attachment}}

The dyadic floors cannot be replaced by the odd-prime dimension formula.
At stage \(j\), the centralizer grade has dimension \(2^{d-j}\). Its
contribution to \(J^1/H^1\) is \[
\bigl(2^{d-j+1}-2^{d-j}\bigr)\,\mathbf1_{n_j\text{ even}}.
\] Every \(n_j\) for \(j<d\) is even; \(n_d=q_1\) is odd and contributes
zero. The sum is exactly \(N-2\). BK nondegeneracy makes this a
symplectic space. A whole source packet of \(d\) planes fits if and only
if \(N-2\ge2d\), equivalently \(d\ge3\).

\begin{longtable}[]{@{}
  >{\raggedright\arraybackslash}p{(\columnwidth - 8\tabcolsep) * \real{0.2000}}
  >{\raggedright\arraybackslash}p{(\columnwidth - 8\tabcolsep) * \real{0.2000}}
  >{\raggedright\arraybackslash}p{(\columnwidth - 8\tabcolsep) * \real{0.2000}}
  >{\raggedright\arraybackslash}p{(\columnwidth - 8\tabcolsep) * \real{0.2000}}
  >{\raggedright\arraybackslash}p{(\columnwidth - 8\tabcolsep) * \real{0.2000}}@{}}
\toprule\noalign{}
\begin{minipage}[b]{\linewidth}\raggedright
Ambient rank
\end{minipage} & \begin{minipage}[b]{\linewidth}\raggedright
Host quotient dimension
\end{minipage} & \begin{minipage}[b]{\linewidth}\raggedright
Host Heisenberg carrier
\end{minipage} & \begin{minipage}[b]{\linewidth}\raggedright
Required source carrier
\end{minipage} & \begin{minipage}[b]{\linewidth}\raggedright
Outcome
\end{minipage} \\
\midrule\noalign{}
\endhead
\bottomrule\noalign{}
\endlastfoot
2 & 0 & 1 & 2 & Insufficient \\
4 & 2 & 2 & 4 & Insufficient \\
8 & 6 & 8 & 8 & Exact capacity \\
16 & 14 & 128 & 16 & Spare capacity \\
\end{longtable}

Now assume \(d\ge3\), the one-vertex contract, declared \(C=\{1\}\), and
actually supplied full character extensions with labeled whole-plane
quotients. Retain their complete dyadic finite central-extension
markings, including their specified central embeddings, lifts, square
values, refinements and phases. Any central-character kernel already
factored out must be recorded; a scalar pushout cannot recover invisible
kernel elements.

Under the scalar-pushout realization contract, with no independently
fixed conflicting host marking, choose a labeled symplectic embedding
into \(W\). Matching commutators gives a center-fixing isomorphism
between the source \(\mathbb C^\times\) extension and the corresponding
host restriction. \textbf{Retain the image of the whole marked finite
subgroup and every lift.} For this isomorphism \(F_0\), \[
F_0(a)^2=F_0(a^2).
\] Thus its square bits, refinement and Gauss sign survive, along with
all supplied full character values. They are not inferred from the
unmarked pushout. Stone--von Neumann gives the source factor inside
\(\eta\); a complementary line gives an actual rank-\(N\) image-algebra
projector. The retained induction/recovery idempotent recovers the same
\(\Lambda\) and the same \(\pi\).

This gives an actual represented marked finite subgroup in the scalar
extension of the actual type. It permits complex scalars times type
operators, as the declared contract does. It does not assert literal
inclusion of old matrix groups, numerical depth placement, or
identification with a separately fixed canonical native finite subgroup.

If a genuine host-native quadratic refinement and its finite
\(\mu_2\)-central extension are independently supplied, quadratic
embedding has a sharper condition. With otherwise free embeddings, at
rank eight the equal-dimensional source and host must have equal Arf
invariants. For \(d\ge4\), the host has at least one spare symplectic
plane, so either source Arf type embeds. If the original embedding is
preassigned, it must instead satisfy the pointwise equation
\(q_W(f(v))=q_{\rm src}(v)\). These conclusions require the actual host
marking as input; dimension does not manufacture it.

The full finite proofs are in \path{markings/DYADIC_MARKINGS167.md}.
In particular, the D8 and Q8 unmarked complex pushouts coincide while
their whole marked finite \(\mu_2\)-central extensions differ. Merely
recording the central \(\{\pm1\}\) subgroup would lose that distinction.

\hypertarget{dyadic_synthesis-failed-repairs-and-boundaries}{%
\subsection{5. Failed repairs and
boundaries}\label{dyadic_synthesis-failed-repairs-and-boundaries}}

The rank-two and rank-four carrier failures cannot be repaired by
changing the maximal simple parameter while keeping the specified Swan
exponent and torsion number. At rank two, odd Swan forces a minimal
odd-level stratum and \(J^1=H^1\). At rank four, Swan \(6h_2\) with
\(h_2\) odd forces the degree pattern \(4\to2\); the maximal quadratic
derived field forces an odd lower critical level. Every such
maximal-type carrier has dimension two, below the required four. Scalar
twists preserve this capacity.

This is an obstruction to the whole-source sector inside that
maximal-type carrier. It is not a counterexample to the original theorem
or a proof that a different arithmetic carrier in the same
representation cannot work. Section 7 constructs such carriers and
proves their exact recovery.

At rank eight, numerical profiles alone also fail to determine capacity.
\path{capacity/DYADIC_CAPACITY167.md}, Section 6, constructs an
actual skipped-degree \(8\to2\) host with the same Swan, torsion and
self-pair conductor but quotient dimension four rather than six. Our
successful result uses the explicit full tower, rather than inferring
its profile from invariants.

Thus the completed dyadic scope is:

\begin{itemize}
\tightlist
\item
  Actual fields, full characters and central completion in every rank
  \(2^d\).
\item
  Full marked scalar-pushout packet realization for the stated family
  when \(d\ge3\), with \(C=1\), one vertex, and transported markings
  retained.
\item
  A conditional quadratic embedding into an independently supplied
  native marked host, with the exact Arf or pointwise condition above.
\end{itemize}

Section 7 constructs different actual arithmetic carriers at the two low
ranks and recovers the same candidate from them. The maximal-type
capacity failure above remains true. Original H5 stays the original
same-candidate conditional premise; these conductor identities and
noncompact root choices do not prove it.

\hypertarget{dyadic_synthesis-evidence}{%
\subsection{6. Evidence}\label{dyadic_synthesis-evidence}}

The arbitrary-length field and full-character arguments are ordinary
proofs. The exact rank-eight run for \(h=(1,3,5)\) checks actual matrix
identities, normalized lifts, leading and adjacent-correction grades
\(-60,-18,-3\), native traces one and ambient traces \(0,0,1\). It
outputs a matrix with 25 nonzero Laurent-polynomial entries. The
corresponding critical exponents are proved symbolically. The run does
not prove simplicity, full-character existence, or universal marked
realization. Those conclusions use the separate proofs and their stated
scope. No Lean build was run.

The base synthesis before the low-rank continuation is preserved in
\path{history/PROOF167_BEFORE_ENLARGEMENT.md}, together with its
review. Its previously open enlarged-carrier route is executed next.

\hypertarget{dyadic_synthesis-the-low-rank-repair-enlarge-the-actual-inducing-carrier}{%
\subsection{7. The low-rank repair: enlarge the actual inducing
carrier}\label{dyadic_synthesis-the-low-rank-repair-enlarge-the-actual-inducing-carrier}}

The maximal-type carrier is not the only arithmetic carrier available in
the constructed representation. The following constructions keep the
original \(\beta,\theta,\Lambda\) and \(\pi\) fixed.

\hypertarget{dyadic_synthesis-rank-two}{%
\subsubsection{Rank two}\label{dyadic_synthesis-rank-two}}

Write the odd minimal level as \(q=2m+1\), with \(m\ge1\), and let
\(J=H^1=J^1=U_K^1U_A^{m+1}\). Define \[
P=U_K^1U_A^m.
\] This is an actual compact open group, with \(J\triangleleft P\),
\([P:J]=2\) and \(P\cap(K^\times J)=J\). For \(g\in P\setminus J\), the
actual intertwining theorem makes \(\lambda=\theta^g/\theta\)
nontrivial. The commutator \([P,J]\subset U_A^{m+1}\) lies in the forced
additive-character domain, so \(\lambda\) has values \(\{\pm1\}\). Put
\(H_0=\ker\lambda\). Then \[
P/H_0\simeq\mathbb F_2^2,\qquad
E=\operatorname{Ind}_J^P\theta\text{ is irreducible of dimension }2.
\] The actual represented generators are a diagonal sign and an
off-diagonal matrix that anticommute. Thus their commutator form is
nondegenerate, and \(\theta|_{H_0}\) is the full character of an actual
new arithmetic Heisenberg pair. No abstract operators create that group
action.

The identity compact-induction double coset realizes \(E\) inside the
same \(\pi\), as functions supported on \((K^\times J)P\). The complete
group and recovery proof is
\path{lowrank/RANK2_ENLARGED_CARRIER167.md}.

\hypertarget{dyadic_synthesis-rank-four}{%
\subsubsection{Rank four}\label{dyadic_synthesis-rank-four}}

Use the actual rank-four outer approximation field \(E_0/F\) of degree
two and its centralizer \(B=\operatorname{End}_{E_0}(F^4)\). The odd
lower critical level is \(r=3h_1=2m+1\); the half leading level is
\(a=3h_2>r\). Define \[
P=J\,U_B^m,
\] where \(J=J^1\) is the actual compact type group. The filtered
centralizer estimate for the actual final field, without assuming
\(E_0\subset K\), proves that \(P\) is a compact group with \(J\) normal
of index two. For a concrete new coset representative \(u\),
\(\lambda=\theta^u/\theta\) on the actual \(H=H^1\) is nontrivial with
values \(\{\pm1\}\). Put \(H_0=\ker\lambda\subset H\). The calculation
proves \[
P/H_0\simeq\mathbb F_2^4,\qquad
E=\operatorname{Ind}_J^P\eta\text{ is irreducible of dimension }4.
\] The new represented commutators are scalar, its quotient has exponent
two, and nondegeneracy follows from its exact scalar preimage and
irreducibility. The finite scalar center may be larger than \(\mu_2\); a
literal extraspecial quotient is not asserted. The full character values
are retained by the actual finite Heisenberg extension. Again \(E\)
occurs inside the same \(\pi\) through the identity double coset. See
\path{lowrank/RANK4_ENLARGED_CARRIER167.md}.

\hypertarget{dyadic_synthesis-exact-recovery-with-the-noncompact-normalization-retained}{%
\subsubsection{Exact recovery, with the noncompact normalization
retained}\label{dyadic_synthesis-exact-recovery-with-the-noncompact-normalization-retained}}

At both low ranks, \(K^\times\) normalizes the new \(P\). In rank two
this follows directly from normalization of \(A\) and its radical. In
rank four, a \(K\)-uniformizer is approximated by the derived chain
shift \(R\) times a unit of \(B\), with relative error in \(P_A^{n-r}\).
Its conjugation error on \(U_B^m\) lies in \(U_A^{n-r+m}\subset J\),
which proves normalization. The exact argument is
\path{lowrank/RANK4_INDUCTION_BY_STAGES167.md}.

Consequently \(\mathcal M=K^\times P\) is an actual open group compact
modulo the center, \(\mathcal N=K^\times J\) is normal of index two in
\(\mathcal M\), and \[
\Sigma=\operatorname{Ind}_{\mathcal N}^{\mathcal M}\Lambda,
\qquad \Sigma|_P\simeq E,
\qquad
\boxed{\pi\simeq\operatorname{c\!-\!Ind}_{\mathcal M}^G\Sigma.}
\] Restriction to \(P\) is irreducible, so \(\Sigma\) is irreducible.
Induction in stages proves the last equality for the original
\(\Lambda\), including its fixed central character and noncompact root.
This is a genuine enlarged inducing representation; it changes neither
the candidate nor its original canonical beta-extension data.

There is also a recovery from compact data explicitly recording that
root. With \(A_0=\Lambda(\varpi_K)\) and
\(X=\operatorname{c\!-\!Ind}_J^G\eta\), define
\((Tf)(g)=A_0^{-1}f(\varpi_Kg)\). Then \(X/(T-1)X\simeq\pi\). After
imposing the fixed central character, \(\mathcal N/(F^\times J)\) is
cyclic of order \(N\), and the finite average
\(N^{-1}\sum_{i=0}^{N-1}T^i\) selects this same extension. These
operators act on complex representation spaces, so their scalar division
by \(N\) is valid. This is a different recovery from the Block159
projector onto a sector of \(\eta\). The exact kernel and finite
projector are proved in
\path{lowrank/ENLARGED_CARRIER_CONTRACT167.md}.

\hypertarget{dyadic_synthesis-what-the-low-rank-repair-admits}{%
\subsubsection{What the low-rank repair
admits}\label{dyadic_synthesis-what-the-low-rank-repair-admits}}

Apply the marked scalar-pushout theorem of Section 4 to the
\textbf{actual new} pair \((H_0,P)\) and its carrier \(E\). It has
exactly the capacity of the whole source at each low rank. All full
source character values and the transported whole finite subgroup,
lifts, squares and phase markings can be attached and retained. The
displayed \(\Sigma\) and induction in stages recover the same candidate.

The old H2 break targets remain their supplied actual simple-character
quotients. The enlarged pair is their common arithmetic assembly
carrier; it is not renamed as the original maximal-simple \(H^1,J^1\)
pair. If a stronger reading requires the whole packet itself to be the
quotient of a maximal simple character in the same candidate, the
earlier dimension obstruction still applies. Likewise, agreement with an
independently fixed native quadratic marking must be checked separately.
The contract note proves this distinction explicitly.

\hypertarget{dyadic_synthesis-final-scope-and-next-construction-question}{%
\subsection{8. Final scope and next construction
question}\label{dyadic_synthesis-final-scope-and-next-construction-question}}

The weighted dyadic mechanism now gives actual arithmetic fields, full
characters and tame central completion in every rank \(2^d\) of the
stated increasing odd-break family. Under the one-vertex, \(C=1\), full
marked scalar-pushout \textbf{assembly} contract, the packet carrier can
be realized at every length: inside the original type for \(d\ge3\), and
in the explicit enlarged inducing carrier at \(d=1,2\). The source's
prescribed central character must belong to the covered tame class. In
particular no arbitrary single-factor wild determinant is assumed to be
trivial.

Remaining universal obligations are construction/comparison of
independently prescribed canonical native markings, nontrivial
symplectic stabilizer actions, general source/transport diagrams and any
required numerical layer placement. Original H5 remains conditional for
this same fixed candidate. No equality of profiles or choice among
noncompact roots is promoted to H5.

The next concrete target is equivariance of the arithmetic attachment
for a nontrivial prescribed symplectic action. Another matrix search
inside the old low-rank maximal-type carrier would repeat its already
proved obstruction. \path{STATUS167.json} and \path{ROUTES167.md}
record the achieved scopes and open routes.

\clearpage
\section{The explicit dyadic arithmetic tower}\label{app:dyadic_tower}

No character, quotient-dimension,
source attachment, pair-factor, or universal Chapter XI closure is
asserted here.

\hypertarget{dyadic_tower-statement}{%
\subsection{1. Statement}\label{dyadic_tower-statement}}

Let F=F\_2((t)), let 0\textless h\_1\textless\ldots\textless h\_d be odd
positive integers, and set

\begin{Verbatim}[fontsize=\small,breaklines=true,breakanywhere=true]
q_i=3h_i,  N=2^d.
\end{Verbatim}

The explicit recursion below constructs B\_d in M\_N(F) such that

\begin{Verbatim}[fontsize=\small,breaklines=true,breakanywhere=true]
F[B_d] is a totally ramified maximal field normalizing
   the complete period-N lattice chain,
v_A(B_d)=-2^(d-1)q_d,
k_0(B_d,A)=-q_1,
tr_F(B_d)=1.
\end{Verbatim}

There are explicit native constructors of ranks 2\^{}j for the largest j
exponents. Their identity-amplifications give the actual defining
sequence with field degrees 2,4,\ldots,2\^{}d and critical levels

\begin{Verbatim}[fontsize=\small,breaklines=true,breakanywhere=true]
n_j=2^(d-j)q_(d-j+1),  1<=j<=d.
\end{Verbatim}

Every parameter in this ambient defining sequence has the same leading
level n\_1. The quantities n\_2,\ldots,n\_d are its successive critical
approximation levels, not new leading levels of the ambient matrices.

Every native constructor has trace 1. Every correction from an
identity-2 amplification of the preceding native constructor to the next
native constructor also has trace 1. In full ambient rank 2\^{}d, the
earlier amplified approximations instead have trace 0, because their
multiplicities are even. This distinction is essential.

\hypertarget{dyadic_tower-the-minimal-degree-two-building-block}{%
\subsection{2. The minimal degree-two building
block}\label{dyadic_tower-the-minimal-degree-two-building-block}}

Let L be a complete discretely valued field of characteristic 2 with
residue field F\_2, and choose a uniformizer t\_L. For odd
q\textgreater0 put a=(q+1)/2 and define

\begin{Verbatim}[fontsize=\small,breaklines=true,breakanywhere=true]
E=L(alpha),  alpha^2-alpha=t_L^(-q),
z=t_L^a alpha.
\end{Verbatim}

The Artin--Schreier polynomial is irreducible: an element x\^{}2-x of
negative L-valuation has even negative valuation. Its derivative is 1,
and the valuation equation forces ramification index 2. Thus E/L is
separable and totally ramified, v\_E(alpha)=-q, and alpha is minimal.

The element z has E-valuation 2a-q=1 and satisfies

\begin{Verbatim}[fontsize=\small,breaklines=true,breakanywhere=true]
f(z)=z^2+t_L^a z+t_L=0.
\end{Verbatim}

Thus 1,z is a uniformizer power basis, and

\begin{Verbatim}[fontsize=\small,breaklines=true,breakanywhere=true]
alpha=z/t_L^a,  Tr_(E/L)(alpha)=1.
\end{Verbatim}

\hypertarget{dyadic_tower-an-explicit-normalized-corestriction-and-lift}{%
\subsection{3. An explicit normalized corestriction and
lift}\label{dyadic_tower-an-explicit-normalized-corestriction-and-lift}}

On E use the coefficient basis 1,z. Its trace-dual elements are

\begin{Verbatim}[fontsize=\small,breaklines=true,breakanywhere=true]
z_0^*=(z+t_L^a)/t_L^a,
z_1^*=1/t_L^a.
\end{Verbatim}

Let pr(T\_ij)=z\^{}i z\_j\^{}* be the trace corestriction. Put

\begin{Verbatim}[fontsize=\small,breaklines=true,breakanywhere=true]
kappa=t_L^a/z=alpha^(-1),
s=kappa pr,
u_0=s(P0)=t_L/z^2.
\end{Verbatim}

Here P0 is constant-coefficient projection, applied separately to each
coordinate when acting on E\^{}m. The identity for u\_0 follows from
f(z)=0. It is a unit: v\_E(t\_L)=v\_E(z\^{}2)=2. More explicitly,

\begin{Verbatim}[fontsize=\small,breaklines=true,breakanywhere=true]
s(T_i0)=z^i(1+t_L^a/z),
s(T_i1)=z^(i-1).
\end{Verbatim}

Since v\_E(t\_L\^{}a/z)=2a-1=q\textgreater0, these preserve the expected
leading grades. On End\_L(E\^{}m), take the complete interleaved chain
in the basis w\_(mi+b)=z\^{}i v\_b. Its centralizer is M\_m(E) with its
complete E-chain. The formulas give s(P\^{}r) contained in Q\^{}r for
all integers r.

For C in M\_m(E), define the explicit L-linear operator

\begin{Verbatim}[fontsize=\small,breaklines=true,breakanywhere=true]
Lift_E(C)=u_0^(-1) C P0=(z^2/t_L) C P0.
\end{Verbatim}

B-bimodularity gives s(Lift\_E(C))=C. The lift preserves exact
valuation: u\_0\^{}(-1) is an E-unit, and expansion in 1,z gives
coefficient valuations distinct modulo 2. Each coefficient entry has the
same grade in the interleaved L-chain and the E-chain. In particular
\(s(P^r)=Q^r\).

\hypertarget{dyadic_tower-exact-trace-transport}{%
\subsection{4. Exact trace transport}\label{dyadic_tower-exact-trace-transport}}

For every C in M\_m(E), ordinary matrix traces satisfy

\begin{Verbatim}[fontsize=\small,breaklines=true,breakanywhere=true]
tr_L(Lift_E(C))
   =Tr_(E/L)(z_0^* u_0^(-1) tr_E(C))
   =Tr_(E/L)(kappa^(-1) tr_E(C))
   =Tr_(E/L)(alpha tr_E(C)).                 (T2)
\end{Verbatim}

The first equality follows by taking the L-trace of each rank-one
operator yP0, which is the constant coefficient of y. This coefficient
is Tr\_(E/L)(z\_0\^{}* y).

Thus the trace of a lift is 1 whenever the inner ordinary E-matrix trace
is 1. It is not asserted that arbitrary trace values are unchanged.

\hypertarget{dyadic_tower-recursive-actual-matrices-and-fields}{%
\subsection{5. Recursive actual matrices and
fields}\label{dyadic_tower-recursive-actual-matrices-and-fields}}

Define B(L,t\_L;h\_1,\ldots,h\_d) as follows. For d=1 use the
multiplication matrix of alpha from Section~2 with q=q\_1. For
d\textgreater=2 apply Section~2 with q=q\_d, set m=2\^{}(d-1),
recursively construct

\begin{Verbatim}[fontsize=\small,breaklines=true,breakanywhere=true]
C=B(E,z;h_1,...,h_(d-1)),
\end{Verbatim}

and put

\begin{Verbatim}[fontsize=\small,breaklines=true,breakanywhere=true]
B(L,t_L;h_1,...,h_d)=alpha I_m+(z^2/t_L) C P0.       (D2)
\end{Verbatim}

This is a finite algebraic recipe. It uses specified separable quadratic
extensions and no unknown support, genericity choice, or correspondence
inverse.

The degree and simplicity proof is by induction. The base alpha is
minimal and generates E. At a later step let gamma=alpha I\_m and
b=Lift\_E(C). The outer minimal stratum has

\begin{Verbatim}[fontsize=\small,breaklines=true,breakanywhere=true]
n=2^(d-1)q_d,  k_0(gamma,A)=-n.
\end{Verbatim}

The inner induction hypothesis gives an actual maximal field E{[}C{]}
normalizing the complete E-chain with

\begin{Verbatim}[fontsize=\small,breaklines=true,breakanywhere=true]
R=2^(d-2)q_(d-1),  v(C)=-R,  k_0(C)=-q_1.
\end{Verbatim}

Because h\_d\textgreater h\_(d-1), we have n\textgreater2R. Therefore

\begin{Verbatim}[fontsize=\small,breaklines=true,breakanywhere=true]
-2R > -n-q_1 = k_0(gamma,A)+k_0(C).
\end{Verbatim}

The derived stratum with parameter s(b)=C is pure, and its field is
maximal. Bushnell--Kutzko Remark 2.2.7 consequently gives the field
generated by the actual gamma+b, not merely an equivalent replacement.
It has ramification and degree 2\^{}d, residue degree one, normalizes
the complete chain, and has k\_0=max(-n,-q\_1)=-q\_1. Since
v(b)\textgreater v(gamma), its level is n and {[}A,n,0,gamma+b{]} is
simple.

The derived model E{[}C{]} is not identified with L{[}gamma+b{]}. In
particular no literal embedding of E into the latter field is claimed.

\hypertarget{dyadic_tower-trace-1-at-every-native-stage}{%
\subsection{6. Trace 1 at every native
stage}\label{dyadic_tower-trace-1-at-every-native-stage}}

The base matrix has trace Tr\_(E/L)(alpha)=1. Inductively tr\_E(C)=1, so
(T2) gives tr\_L(Lift\_E(C))=1. For d\textgreater=2 the L-trace of alpha
I\_m is

\begin{Verbatim}[fontsize=\small,breaklines=true,breakanywhere=true]
m Tr_(E/L)(alpha)=0,
\end{Verbatim}

because m is even. Hence tr\_L(B)=1 at every native rank.

For any j\textgreater=1, let A\_j denote the native rank-2\^{}j
constructor using a fixed list of the largest j exponents. The next
native rank-2\^{}(j+1) constructor A\_(j+1) uses the same largest
exponents and adds the next smaller one. With their compatible
interleaved bases, put

\begin{Verbatim}[fontsize=\small,breaklines=true,breakanywhere=true]
Delta_(j+1)=A_(j+1)-A_j tensor I_2.
\end{Verbatim}

Then

\begin{Verbatim}[fontsize=\small,breaklines=true,breakanywhere=true]
tr(Delta_(j+1))=1-2*1=1                    in characteristic 2.
\end{Verbatim}

This native trace identity is stronger and more relevant than a claim
obtained merely by amplifying to a rank divisible by two.

\hypertarget{dyadic_tower-full-arithmetic-defining-sequence}{%
\subsection{7. Full arithmetic defining
sequence}\label{dyadic_tower-full-arithmetic-defining-sequence}}

For the final rank 2\^{}d define

\begin{Verbatim}[fontsize=\small,breaklines=true,breakanywhere=true]
B^[j]=B(F,t;h_(d-j+1),...,h_d) tensor I_(2^(d-j)).
\end{Verbatim}

All truncations use the same top q\_d to select the first extension and
uniformizer, the same next q\_(d-1) to select the next extension, and so
on. The coefficient projection and lift commute with amplification. Thus
these are literal truncations of (D2).

Their parameter fields have degree 2\^{}j, and their critical levels in
the ambient complete chain are n\_j=2\^{}(d-j)q\_(d-j+1).
Valuation-preserving lifts give

\begin{Verbatim}[fontsize=\small,breaklines=true,breakanywhere=true]
v_A(B^[j+1]-B^[j])=-n_(j+1).
\end{Verbatim}

Later differences have strictly greater valuation, since
n\_j\textgreater2n\_(j+1). This proves that the displayed parameters
form an actual defining sequence of simple approximations.

Their ambient traces are

\begin{Verbatim}[fontsize=\small,breaklines=true,breakanywhere=true]
tr_F(B^[j])=0 for j<d,  tr_F(B^[d])=1.
\end{Verbatim}

Consequently ambient corrections before the last step have trace 0, and
the last correction has trace 1. Every one of these trace values is
integral. On a scalar principal unit u, the additive correction
character has value psi\_F((u-1)tr(delta)); this is 1 for an additive
character trivial on the maximal ideal. No assertion about full simple
character extension follows here without its separate intersection and
transfer argument.

\hypertarget{dyadic_tower-concrete-rank-eight-instance-and-limits}{%
\subsection{8. Concrete rank-eight instance and
limits}\label{dyadic_tower-concrete-rank-eight-instance-and-limits}}

For h=(1,3,5), the construction gives actual approximation-field degrees
2,4,8 and critical levels 60,18,3. The supplementary exact program
constructs its Laurent-polynomial matrices, verifies the normalized lift
identities, and checks the native/ambient trace distinctions and matrix
grades. It does not prove field existence or simplicity by bounded
computation: those are consequences of Section~5.

The odd-prime Heisenberg quotient and conductor counting from Block166
cannot simply be copied: q\_1 is odd here. This note makes no dyadic
quotient, full-character, stabilizer, phase, source-attachment, or H5
claim. It also changes lower parameter conjugacy classes, as in the
corresponding odd-prime Artin--Schreier construction.

\hypertarget{dyadic_tower-dependencies-and-provenance}{%
\subsection{Dependencies and
provenance}\label{dyadic_tower-dependencies-and-provenance}}

Primary theorem: Bushnell--Kutzko, The admissible dual of GL(N) via
compact open subgroups, Remark 2.2.7 (printed p. 55), together with the
amplification normalization and defining-sequence conventions already
checked in Block166.

Retained primary source text SHA256:
\path{d69db0827256707a34cd2562af68b9de11fc889ab3f43535ec47b620d58ca9c8}.

Block166 inputs are pinned separately in \path{INPUT_PINS167.json} and are not
modified by this continuation. The Artin--Schreier, trace-dual, and
trace calculations needed specifically for p=2 are proved above.

\clearpage
\section{Dyadic full characters and the odd critical endpoint}\label{app:dyadic_character}

Source record: Paper\_XI\_Dyadic167. Dependency:
Paper\_XI\_Attachment166.

Evidence: ordinary mathematical proof using the actual Artin--Schreier
matrix tower and the retained Bushnell--Kutzko character and transfer
theorems. No Lean verification is claimed. This note proves no dyadic
quotient-capacity, quadratic-phase, packet-attachment, or H5 result.

\hypertarget{dyadic_character-exact-theorem}{%
\subsection{1. Exact theorem}\label{dyadic_character-exact-theorem}}

Let F=F\_2((t)), let psi\_F be an additive character of conductor t O\_F
(trivial there and nontrivial on O\_F), and let

\begin{Verbatim}[fontsize=\small,breaklines=true,breakanywhere=true]
0<h_1<...<h_d,       each h_i odd,
q_i=3h_i,           N=2^d.
\end{Verbatim}

Use the dyadic Artin--Schreier recursion specified in Section~2. It gives
the actual maximal simple parameter beta\_d with its native arithmetic
defining sequence. There is an actual full simple character theta of its
maximal simple stratum such that

\begin{Verbatim}[fontsize=\small,breaklines=true,breakanywhere=true]
theta restricted to U_F^1 is trivial.                         (1)
\end{Verbatim}

The construction retains the actual canonical beta extension and can be
completed with any prescribed central quasicharacter omega\_F trivial on
U\_F\^{}1. Because \(k_F^\times=\FF_2^\times\) is trivial,
such an omega\_F is unramified. Thus this covers every tamely ramified
central quasicharacter in this particular residue-field setting.

The corresponding compact induction is an actual irreducible admissible
supercuspidal of rank N, by the retained maximal-simple-type theorem. No
claim is made here about its Heisenberg quotient size or about
realization of a supplied dyadic finite packet.

The trace-zero argument from Block 166 does not apply verbatim: at p=2 the
primitive Artin--Schreier parameter has trace one. The correct invariant
is \textbf{integral trace}, not necessarily zero trace. Since psi\_F
kills t O\_F, integral correction trace is sufficient on scalar
principal units. The recursion below preserves the stronger exact value
one.

\hypertarget{dyadic_character-actual-dyadic-recursion-and-exact-native-traces}{%
\subsection{2. Actual dyadic recursion and exact native
traces}\label{dyadic_character-actual-dyadic-recursion-and-exact-native-traces}}

At a recursive step let L be a characteristic-two local field with
residue field F\_2, uniformizer t\_L, and a positive odd q. Define

\begin{Verbatim}[fontsize=\small,breaklines=true,breakanywhere=true]
E=L(alpha),       alpha^2-alpha=t_L^(-q).
\end{Verbatim}

The right side is not an Artin--Schreier coboundary: a negative
valuation of x\^{}2-x is even, whereas -q is odd. The degree-two
polynomial is therefore irreducible and separable. Its valuation
equation gives total ramification and v\_E(alpha)=-q. Put

\begin{Verbatim}[fontsize=\small,breaklines=true,breakanywhere=true]
a=(q+1)/2,       z=t_L^a alpha.
\end{Verbatim}

Then z is an actual uniformizer and its Eisenstein polynomial is

\begin{Verbatim}[fontsize=\small,breaklines=true,breakanywhere=true]
f(Z)=Z^2+t_L^a Z+t_L.                                      (2)
\end{Verbatim}

Its two conjugate alpha values are alpha and alpha+1, so

\begin{Verbatim}[fontsize=\small,breaklines=true,breakanywhere=true]
Tr_(E/L)(alpha)=1.                                         (3)
\end{Verbatim}

Use the coefficient basis 1,z, the trace corestriction pr, and the same
normalized corestriction convention as in Block 166:

\begin{Verbatim}[fontsize=\small,breaklines=true,breakanywhere=true]
kappa=-f'(z)/z = t_L^a/z = alpha^(-1),
s=kappa pr.
\end{Verbatim}

Minus and plus agree in characteristic two. If P0 extracts the constant
coefficient on every E-coordinate, then

\begin{Verbatim}[fontsize=\small,breaklines=true,breakanywhere=true]
u_0=s(P0)=t_L/z^2 is a unit,
Lift_E(C)=u_0^(-1) C P0,
s(Lift_E(C))=C.
\end{Verbatim}

The trace-dual calculation in Block 166 gives, as an identity of L-matrix
traces for every matrix C over E,

\begin{Verbatim}[fontsize=\small,breaklines=true,breakanywhere=true]
tr_L Lift_E(C)
   =Tr_(E/L)(kappa^(-1) tr_E(C))
   =Tr_(E/L)(alpha tr_E(C)).                               (4)
\end{Verbatim}

The lift preserves the complete-chain filtration valuation, by expansion
in 1,z. Thus it is the actual filtered right inverse used in the field
constructor, not an arbitrary inverse on vector spaces.

Define B(L,t\_L;h\_1,\ldots,h\_s) recursively. At s=1 use the L-matrix
of multiplication by alpha with q=q\_1. At s\textgreater=2 use (2) with
q=q\_s, m=2\^{}(s-1), and

\begin{Verbatim}[fontsize=\small,breaklines=true,breakanywhere=true]
C=B(E,z;h_1,...,h_(s-1)),
B(L,t_L;h_1,...,h_s)=alpha I_m+Lift_E(C).                    (5)
\end{Verbatim}

The base trace is one by (3). If tr\_E(C)=1, then (4) gives lift trace
one. Since m is even, the trace of alpha I\_m as an L-matrix is m
Tr\_(E/L)(alpha)=0. Consequently

\begin{Verbatim}[fontsize=\small,breaklines=true,breakanywhere=true]
tr_L B(L,t_L;h_1,...,h_s)=1                                (6)
\end{Verbatim}

for every native recursion, by induction. There is no division by the
rank and no confusion between a native trace and its amplified trace.

The actual field proof from Block 166 applies with this degree-two extension
and normalized lift. At a nontrivial step its hypotheses are
n=2\^{}(s-1)q\_s, R=2\^{}(s-2)q\_(s-1), s(Lift(C))=C, and

\begin{Verbatim}[fontsize=\small,breaklines=true,breakanywhere=true]
n>2R,
-2R>-n-q_1.
\end{Verbatim}

The first strict inequality follows from h\_s\textgreater h\_(s-1).
These are the conditions used in BK Remark 2.2.7. The displayed beta
itself generates the totally ramified degree-2\^{}s field normalizing
the complete chain, and its final critical exponent is -q\_1. This
statement does not identify the final field with its auxiliary derived
field.

For the final length d write the native shorter matrices as

\begin{Verbatim}[fontsize=\small,breaklines=true,breakanywhere=true]
b_j=B(F,t;h_(d-j+1),...,h_d),       1<=j<=d.
\end{Verbatim}

They are the compatible native truncations of (5), because the outer
extension, uniformizer, and lift depend on the same retained largest h
at each stage. For \(1\le j<d\), in rank 2\^{}(j+1), put

\begin{Verbatim}[fontsize=\small,breaklines=true,breakanywhere=true]
gamma_j=b_j tensor I_2,
delta_j=b_(j+1)-gamma_j.
\end{Verbatim}

The defining-sequence calculation of Block 166 gives gamma\_j as a simple
approximation at the critical level r\_j=q\_(d-j) of b\_(j+1), and
delta\_j has grade -r\_j. Equation (6) gives the exact scalar identity

\begin{Verbatim}[fontsize=\small,breaklines=true,breakanywhere=true]
tr_F(gamma_j)=0,
tr_F(delta_j)=1.                                          (7)
\end{Verbatim}

Thus for every scalar u in U\_F\^{}1 on which its additive correction is
defined,

\begin{Verbatim}[fontsize=\small,breaklines=true,breakanywhere=true]
psi_(delta_j)(u I)=psi_F((u-1)tr_F(delta_j))
                 =psi_F(u-1)=1.                           (8)
\end{Verbatim}

It is (8), not a false trace-zero assertion, that drives the full
character construction.

\hypertarget{dyadic_character-the-odd-critical-level-rounding}{%
\subsection{3. The odd-critical-level
rounding}\label{dyadic_character-the-odd-critical-level-rounding}}

The critical levels r=q\_i are odd. Write r=2s+1, and set

\begin{Verbatim}[fontsize=\small,breaklines=true,breakanywhere=true]
m=floor(r/2)=s,        m+1=(r+1)/2.
\end{Verbatim}

The relevant BK statements use the floor, with no assumption that r or
the residue characteristic is odd or even:

\begin{itemize}
\tightlist
\item
  BK 3.1.9(ii) identifies H\^{}k(beta) and H\^{}k(gamma) for
  k\textgreater=floor(r/2)+1.
\item
  BK 3.3.18 gives the correction bijection at
  floor(r/2)\textless=m\textless=r-1.
\item
  BK 3.2.5 extends the resulting character at floor(r/2) to a full
  character, its fibers being the indicated determinant characters.
\end{itemize}

At our exact endpoint m=s, the correction character is legitimate. For
x,y in the radical power P\^{}(s+1),

\begin{Verbatim}[fontsize=\small,breaklines=true,breakanywhere=true]
delta*x*y belongs to P^(2(s+1)-r)=P^1.
\end{Verbatim}

The level-one additive matrix character is trivial on P\^{}1, so the
multiplicative cross term contributes one. Therefore psi\_delta is a
character on U\^{}(s+1), and hence on the recursive subgroup used below.
This checks the borderline odd-level rounding directly. Replacing m by
r/2 as if it were an integer would be incorrect.

The same check applies at the native minimal base level q, which is odd:
its forced additive character begins at floor(q/2)+1=(q+1)/2, and
2*(q+1)/2-q=1.

\hypertarget{dyadic_character-finite-character-extension-at-maximal-rank}{%
\subsection{4. Finite character extension at maximal
rank}\label{dyadic_character-finite-character-extension-at-maximal-rank}}

We reuse the elementary extension procedure with its compatibility
condition explicit. If K/F is a finite totally ramified field, and chi
is a smooth character of U\_K\^{}a trivial on U\_K\^{}a intersect
U\_F\^{}1, define a character on U\_K\^{}a U\_F\^{}1 by

\begin{Verbatim}[fontsize=\small,breaklines=true,breakanywhere=true]
chi_0(ku)=chi(k).
\end{Verbatim}

It is well defined by the intersection condition. Choose M\textgreater=a
with chi\textbar U\_K\^{}M=1. Then chi\_0 is a character of a subgroup
of the finite abelian group \(U_K^1/U_K^M\). Extend it by
adjoining generators and choosing the required complex roots, and pull
it back. The resulting chi\_K extends chi and is trivial on U\_F\^{}1.

All freely chosen character extensions below take place at maximal
native rank, where the centralizer is K itself. We do not apply this
scalar-intersection argument directly inside a nonmaximal matrix
centralizer, which would require a separate determinant-image check.

\hypertarget{dyadic_character-base-full-character-with-trace-one}{%
\subsection{5. Base full character with trace
one}\label{dyadic_character-base-full-character-with-trace-one}}

At native rank two, let E=F{[}alpha{]}, with q=q\_d and a=floor(q/2)+1.
The forced character on U\_E\^{}a is

\begin{Verbatim}[fontsize=\small,breaklines=true,breakanywhere=true]
phi(v)=psi_F(Tr_(E/F)(alpha(v-1))).
\end{Verbatim}

For scalar u in U\_E\^{}a intersect U\_F\^{}1, equation (3) gives

\begin{Verbatim}[fontsize=\small,breaklines=true,breakanywhere=true]
phi(u)=psi_F((u-1)Tr_(E/F)(alpha))=psi_F(u-1)=1.
\end{Verbatim}

Section~4 constructs chi\_E on U\_E\^{}1 extending phi and trivial on
U\_F\^{}1. Define theta\_1 on \(H^1=U_E^1\) U\_(A\_1)\^{}a
by chi\_E on the first factor and psi\_alpha on the second. The
restrictions agree on the intersection. BK 3.2.1 gives the actual full
simple character, with theta\_1\textbar U\_F\^{}1=1.

This uses the true trace value one and only the stated conductor
condition on psi\_F. It does not require an artificial zero-trace
primitive parameter in a quadratic extension.

\hypertarget{dyadic_character-native-rank-induction-via-full-beta-determinant-transfer}{%
\subsection{6. Native-rank induction via full beta-determinant
transfer}\label{dyadic_character-native-rank-induction-via-full-beta-determinant-transfer}}

Assume theta\_j is a full simple character for native b\_j in rank
2\^{}j, with theta\_j\textbar U\_F\^{}1=1. Let K\_j=F{[}b\_j{]}. Since
K\_j is maximal, the character chi\_j=theta\_j\textbar U\_(K\_j)\^{}1 is
its full beta-determinant character.

Transfer theta\_j to the multiplicity-two realization gamma\_j=b\_j
tensor I\_2 in rank 2\^{}(j+1), using BK 3.6.14. Denote the full
transferred character by Theta\_j. BK 3.6.13(ii) preserves the entire
beta-determinant character on U\_(K\_j)\^{}1 at level m0=0. Hence for
every literal scalar u in U\_F\^{}1,

\begin{Verbatim}[fontsize=\small,breaklines=true,breakanywhere=true]
Theta_j(u I)=chi_j(det_(K_j)(u I_2))
            =chi_j(u^2)=theta_j(u^2)=1.                    (9)
\end{Verbatim}

This statement includes nonminimal b\_j. The determinant exponent two is
essential and has not been suppressed.

Now set beta=b\_(j+1), gamma=gamma\_j, delta=delta\_j, r=r\_j,
m=floor(r/2), and

\begin{Verbatim}[fontsize=\small,breaklines=true,breakanywhere=true]
D=H^(m+1)(beta)=H^(m+1)(gamma).
\end{Verbatim}

The equality and endpoint m were verified in Section~3. BK 3.3.18
constructs the actual recursive simple character

\begin{Verbatim}[fontsize=\small,breaklines=true,breakanywhere=true]
vartheta=(Theta_j restricted to D)*psi_delta.
\end{Verbatim}

For u in D intersect U\_F\^{}1, equations (8) and (9) prove

\begin{Verbatim}[fontsize=\small,breaklines=true,breakanywhere=true]
vartheta(u)=1.                                             (10)
\end{Verbatim}

The new field K=F{[}beta{]} is maximal and totally ramified in native
rank 2\^{}(j+1). Its chain has K-period one, so

\begin{Verbatim}[fontsize=\small,breaklines=true,breakanywhere=true]
D intersect U_K^1=U_K^(m+1),
D intersect U_F^1=U_F^ceil((m+1)/2^(j+1)).
\end{Verbatim}

Thus (10) is exactly the compatibility required by Section~4. Extend
vartheta\textbar U\_K\^{}(m+1) to chi\_K on U\_K\^{}1 trivial on
U\_F\^{}1. BK 3.2.5 constructs the full simple character by

\begin{Verbatim}[fontsize=\small,breaklines=true,breakanywhere=true]
theta_(j+1)(kd)=chi_K(k)*vartheta(d),
k in U_K^1, d in D.
\end{Verbatim}

That proposition supplies normalization and well-definedness; its proof
is a character-extension argument and imposes no exclusion of residue
characteristic two. The new theta\_(j+1) is trivial on U\_F\^{}1.
Induction proves (1) in every native rank 2\^{}d.

More generally the character argument needs only integral base field
trace and integral traces for the native corrections. The dyadic tower
establishes their exact value one, while the odd-prime tower of Block 166
establishes zero.

\hypertarget{dyadic_character-canonical-beta-extension-and-prescribed-tame-center}{%
\subsection{7. Canonical beta extension and prescribed tame
center}\label{dyadic_character-canonical-beta-extension-and-prescribed-tame-center}}

Let K/F be the final totally ramified degree-N parameter field and let
eta be the Heisenberg representation for the constructed theta. The
compact type group is

\begin{Verbatim}[fontsize=\small,breaklines=true,breakanywhere=true]
J=O_K^times J^1.
\end{Verbatim}

Here \(\OO_K^\times=U_K^1\) because the residue field is
F\_2. Since U\_K\^{}1 is contained in J\^{}1, we have

\begin{Verbatim}[fontsize=\small,breaklines=true,breakanywhere=true]
J=J^1,       J/J^1=F_2^times={1}.                           (11)
\end{Verbatim}

Thus the actual canonical beta extension on J equals eta; there is no
residue-character choice and no change of that extension. It acts
trivially on the central U\_F\^{}1 by (1). Equation (11) does not
identify any externally prescribed dyadic quadratic phases.

Let omega\_F be any smooth quasicharacter trivial on U\_F\^{}1. Since
\(\OO_F^\times=U_F^1\), it is determined by its value
omega\_F(t). The actual maximal type extends to K\^{}times J.
Concretely, (K\^{}times J)/J is infinite cyclic, generated by a
K-uniformizer. An intertwiner for its conjugation action on eta gives an
extension Lambda\_0. Denote its central character by omega\_0.

Choose a complex scalar s with

\begin{Verbatim}[fontsize=\small,breaklines=true,breakanywhere=true]
s^N=omega_F(t)/omega_0(t).
\end{Verbatim}

Multiply Lambda\_0 by the character of K\^{}times J trivial on J and
taking the K-uniformizer to s. Because t has valuation N in K, this
changes its scalar by s\^{}N and gives the requested central
quasicharacter on all F\^{}times. The value at every element of J,
including the canonical beta extension, is unchanged. This does not
assume that a chosen K-uniformizer raised to N equals t.

The retained maximal-simple-type theorem now gives the actual
irreducible admissible supercuspidal pi=c-Ind\_(K\^{}times
J)\^{}GL\_N(F) Lambda. The noncompact scalar is fixed before any H5
comparison. An incompatible immutably prescribed extension is not
overwritten; as in Block 166, its value must be free or already compatible.

\hypertarget{dyadic_character-what-this-closes-and-what-it-does-not}{%
\subsection{8. What this closes and what it does
not}\label{dyadic_character-what-this-closes-and-what-it-does-not}}

This closes the full-character and actual tame-central-character
construction for the explicit dyadic arithmetic tower. The failure of
trace zero at p=2 is repaired by the correct integral trace invariant.
Both the odd-level floor and the multiplicity-two determinant transfer
are accounted for explicitly.

This note does not compute \(J^1/H^1\) or prove that it
contains a particular dyadic source packet. It does not establish a
quadratic refinement, canonical phase matching, nontrivial stabilizer
action, general transport graph, source admission, or H5. In particular,
an actual supercuspidal host with the required central character is not
asserted to be a completed realization of the original dyadic packet
without those separate proofs.

\hypertarget{dyadic_character-exact-source-checks-and-input-pins}{%
\subsection{9. Exact source checks and input
pins}\label{dyadic_character-exact-source-checks-and-input-pins}}

The following final Block 166 files were read and reused without editing:

\begin{itemize}
\tightlist
\item
  \path{attachment166/central/TRACE_ZERO_CHARACTER_INDUCTION166.md}
  SHA256
  \path{99cb6e80c419e0bf7325e2b2b3fbaba1391903448a3b6a0c4cd9f3944ed406a1}.
\item
  \path{attachment166/tower/TRACE_ZERO_AS_TOWER166.md} SHA256
  \path{49a18fb733f0209b3e38118f5132839a060961edd0bdc0211e1f7b441a4e5653}.
\item
  \path{attachment166/tower/RECURSIVE_TOWER166.md} SHA256
  \path{6f39a077df18eb03f9fab2b6e426529dbe14cf4f7e1cfce95ddb0fcd187d7b33}.
\item
  \path{attachment166/audit/TOWER_CENTRAL_REVIEW166.md} SHA256
  \path{4bc71dd050ddb321a253105c8a4ee7a86185af255532d5605caa7a9de383e26b}.
\end{itemize}

Primary retained source: Bushnell--Kutzko, \emph{The admissible dual of
GL(N) via compact open subgroups}, text SHA256
\path{d69db0827256707a34cd2562af68b9de11fc889ab3f43535ec47b620d58ca9c8}.
The current dyadic check reread Definitions~3.1.7--8 and
Proposition 3.1.9(ii) (printed pp. 90--91), Definition~3.2.1 and
Proposition 3.2.5 (printed pp. 98--100), and Corollary 3.3.18
(printed pp. 112--113). Their level thresholds are floor(r/2)+1 or
floor(r/2), as used above, and their statements do not exclude p=2. The
transfer's preservation of the full beta-determinant character is
BK 3.6.13(ii), with existence Theorem~3.6.14 (printed pp. 136--137), already
directly checked in Block 166 and reread for this application.

The character-extension and correction proofs use neither division by
two nor an odd-order averaging argument. The exact dyadic trace
calculation and the borderline correction-character check are given in
Sections~2--3. No new publication lookup, external
arithmetic-induction/descent constructor, or inverse LLC is used.

\clearpage
\section{Exact dyadic capacity and numerical-profile obstruction}\label{app:dyadic_capacity}

Status: ordinary proof from the supplied actual defining sequence and
the retained primary Bushnell--Kutzko results. No Lean or numerical
computation is used. This note calculates the quotient of an actual
simple type. It does not prove universal H5 or identify an explicitly
constructed host with the source-defined arithmetic candidate.

\hypertarget{dyadic_capacity-exact-hypotheses-and-result}{%
\subsection{1. Exact hypotheses and
result}\label{dyadic_capacity-exact-hypotheses-and-result}}

Let F=F\_2((t)). Fix d\textgreater=1 and increasing positive odd
integers

\begin{Verbatim}[fontsize=\small,breaklines=true,breakanywhere=true]
h_1<...<h_d,       q_i=3h_i,       n=2^d.
\end{Verbatim}

Suppose the actual arithmetic constructor supplies a full-degree totally
ramified parameter beta on the complete period-n chain, with actual
approximating fields of degrees

\begin{Verbatim}[fontsize=\small,breaklines=true,breakanywhere=true]
2,4,...,2^d
\end{Verbatim}

and ambient critical levels

\begin{Verbatim}[fontsize=\small,breaklines=true,breakanywhere=true]
nu_j=2^(d−j) q_(d−j+1),       1<=j<=d.
\end{Verbatim}

These are the supplied dyadic specialization of the recursive tower's
field/profile conclusions. The calculation applies equally to the
original polynomial tower of Block166 and to an Artin--Schreier variant
once that variant's field/profile conclusions are established. No
trace-zero assertion is needed here.

The exact actual Heisenberg quotient is

\begin{Verbatim}[fontsize=\small,breaklines=true,breakanywhere=true]
dim_F2(J^1/H^1)=2^d−2,
dim eta=2^((2^d−2)/2).
\end{Verbatim}

Consequently its symplectic dimension is at least the 2d dimensions of
the whole d-plane source packet exactly when d\textgreater=3. In rank~8,
d=3, the quotient has dimension 6 and eta has dimension 8: the source
carrier fits exactly, with no complementary Heisenberg factor needed.

\begin{longtable}[]{@{}
  >{\raggedright\arraybackslash}p{(\columnwidth - 10\tabcolsep) * \real{0.1667}}
  >{\raggedright\arraybackslash}p{(\columnwidth - 10\tabcolsep) * \real{0.1667}}
  >{\raggedright\arraybackslash}p{(\columnwidth - 10\tabcolsep) * \real{0.1667}}
  >{\raggedright\arraybackslash}p{(\columnwidth - 10\tabcolsep) * \real{0.1667}}
  >{\raggedright\arraybackslash}p{(\columnwidth - 10\tabcolsep) * \real{0.1667}}
  >{\raggedright\arraybackslash}p{(\columnwidth - 10\tabcolsep) * \real{0.1667}}@{}}
\toprule\noalign{}
\begin{minipage}[b]{\linewidth}\raggedright
d
\end{minipage} & \begin{minipage}[b]{\linewidth}\raggedright
Ambient rank
\end{minipage} & \begin{minipage}[b]{\linewidth}\raggedright
Actual quotient dimension
\end{minipage} & \begin{minipage}[b]{\linewidth}\raggedright
eta dimension
\end{minipage} & \begin{minipage}[b]{\linewidth}\raggedright
Source carrier dimension
\end{minipage} & \begin{minipage}[b]{\linewidth}\raggedright
Capacity
\end{minipage} \\
\midrule\noalign{}
\endhead
\bottomrule\noalign{}
\endlastfoot
1 & 2 & 0 & 1 & 2 & Fails \\
2 & 4 & 2 & 2 & 4 & Fails \\
3 & 8 & 6 & 8 & 8 & Exact \\
4 & 16 & 14 & 128 & 16 & Sufficient \\
general & 2\^{}d & 2\^{}d−2 & \(2^{2^{d-1}-1}\) & 2\^{}d &
Sufficient iff d\textgreater=3 \\
\end{longtable}

This improves the leading-layer-only sufficient threshold
d\textgreater=4 to the exact whole-profile threshold d\textgreater=3. It
also exposes a genuine low-rank limitation of this proposed
whole-carrier realization.

\hypertarget{dyadic_capacity-floors-and-ceilings-in-the-actual-bk-orders}{%
\subsection{2. Floors and ceilings in the actual BK
orders}\label{dyadic_capacity-floors-and-ceilings-in-the-actual-bk-orders}}

Write P=rad(A), and let B\_j be the A-centralizer of the degree-2\^{}j
approximation. Its complete-chain grade has F\_2-dimension

\begin{Verbatim}[fontsize=\small,breaklines=true,breakanywhere=true]
b_j=2^(d−j),       b_0=n.
\end{Verbatim}

Put

\begin{Verbatim}[fontsize=\small,breaklines=true,breakanywhere=true]
a_j=ceil(nu_j/2),       t_j=floor(nu_j/2)+1.
\end{Verbatim}

The exact BK definitions, 3.1.7--8, use a\_j as the J cutoff and t\_j as
the H cutoff. Thus a\_j=t\_j for an odd critical level, while
t\_j=a\_j+1 for an even critical level.

The increasing h\_i give

\begin{Verbatim}[fontsize=\small,breaklines=true,breakanywhere=true]
nu_j>2 nu_(j+1),       1<=j<d.
\end{Verbatim}

This strict separation is useful: the centralizer approximation estimate
from BK 3.1.5 puts each newly introduced centralizer into the preceding
one modulo a term above the preceding half-level. In the notation of
these stages, the relevant error level is
nu\_j−nu\_(j+1)\textgreater nu\_j/2. Therefore the successive grade
dimensions really are b\_j; there is no unproved assumption that the
approximation fields are literally nested inside the final field.

Applying the ring recursion gives the J-grade staircase:

\begin{itemize}
\tightlist
\item
  below a\_d the grade dimension is b\_d=1;
\item
  on a\_(j+1)\textless=k\textless a\_j it is b\_j;
\item
  from a\_1 upward it is b\_0=n.
\end{itemize}

The H staircase uses the corresponding t\_j. Hence the exact
contribution of transition j to \(J^1/H^1\) is

\begin{Verbatim}[fontsize=\small,breaklines=true,breakanywhere=true]
(b_(j−1)−b_j) · 1_(nu_j even).
\end{Verbatim}

One can also obtain this by subtracting the exact logarithmic indices of
H and J in 1+P. It is an equality of the group index
\textbar \(J^1/H^1\)\textbar, not a guess from the visible
leading matrix support. BK 3.4.1 makes the resulting commutator pairing
nondegenerate.

For j\textless d the level nu\_j has a factor of 2 and is even. The final
level nu\_d=q\_1 is odd. Therefore

\begin{Verbatim}[fontsize=\small,breaklines=true,breakanywhere=true]
dim_F2 J^1/H^1
  =sum_(j=1)^(d−1) (2^(d−j+1)−2^(d−j))
  =2^d−2.
\end{Verbatim}

The omitted final increment would have dimension 1. It cannot be treated
as a symplectic plane or silently restored: at that odd level the two
defining lattice cutoffs are exactly equal. Its disappearance is also
consistent with the even dimension required of a nondegenerate
alternating space.

The source packet requires dimension 2d. At d=1,2, 2\^{}d−2\textless2d.
At d=3 equality holds. If the inequality holds for d\textgreater=3, then

\begin{Verbatim}[fontsize=\small,breaklines=true,breakanywhere=true]
2^(d+1)−2−2(d+1) = (2^d−2−2d)+(2^d−2)>0,
\end{Verbatim}

so it continues to hold for every larger d.~This proves the exact
threshold.

\hypertarget{dyadic_capacity-conductors-still-match-despite-the-lost-final-contribution}{%
\subsection{3. Conductors still match despite the lost final
contribution}\label{dyadic_capacity-conductors-still-match-despite-the-lost-final-contribution}}

This section reuses the same retained formal-degree and self-pair
conductor identities already audited in Blocks 161 and 165. It is not
needed to prove the dimension formula, but shows why numerical conductor
checks do not expose the low-rank capacity failure.

Put Delta\_j=b\_(j−1)−b\_j. The exact grade count gives

\begin{Verbatim}[fontsize=\small,breaklines=true,breakanywhere=true]
I=log_2[(1+P):J^1]
  =sum_j Delta_j ceil(nu_j/2)−(n−1).
\end{Verbatim}

For this profile only the final level is odd and Delta\_d=1. Therefore

\begin{Verbatim}[fontsize=\small,breaklines=true,breakanywhere=true]
2I=sum_j Delta_j nu_j+1−2(n−1).
\end{Verbatim}

The Heisenberg dimension exponent is D/2, where D=n−2. The retained
formal-degree identity for a full-degree, totally ramified, torsion-one
maximal type gives

\begin{Verbatim}[fontsize=\small,breaklines=true,breakanywhere=true]
a(Pi^vee×Pi)=n^2−1+(n−1)+2I+D.
\end{Verbatim}

Substitution yields exact cancellation of the odd-cutoff correction:

\begin{Verbatim}[fontsize=\small,breaklines=true,breakanywhere=true]
a(Pi^vee×Pi)=n^2−1+sum_j Delta_j nu_j
           =2^(2d)−1+sum_(i=1)^d 3·4^(i−1) h_i.
\end{Verbatim}

The ordinary conductor is n+nu\_1, so Swan(Pi)=2\^{}(d−1)·3h\_d, and the
torsion number is one. These are the usual three numerical invariants of
the stipulated independent dyadic Heisenberg tensor source.

Thus all three numerical invariants can match while eta is too small for
the whole source carrier at ranks 2 and 4. This is a capacity failure, not
a conductor mismatch. It does not by itself refute any original H1--H5
statement.

\hypertarget{dyadic_capacity-rank2-failure-is-forced-by-swan-and-torsion-alone}{%
\subsection{4. Rank 2 failure is forced by Swan and torsion
alone}\label{dyadic_capacity-rank2-failure-is-forced-by-swan-and-torsion-alone}}

Let Pi be any actual supercuspidal representation of GL2(F) with torsion
number one and odd Swan exponent N=3h\_1. BK 6.2.5 forces its maximal
parameter field K/F to be totally ramified of degree~2 and its hereditary
order to have complete period 2. The ordinary conductor formula then
fixes its native level as N.

Every parameter beta of a maximal simple type in Pi has v\_K(beta)=−N,
which is odd. Since the residue extension is trivial, BK 1.4.14--15 say
beta is minimal over F. The two order definitions are therefore

\begin{Verbatim}[fontsize=\small,breaklines=true,breakanywhere=true]
J_order=O_K+P^ceil(N/2),
H_order=O_K+P^(floor(N/2)+1).
\end{Verbatim}

The exponents coincide because N is odd. Hence \(J^1=H^1\),
eta is one-dimensional, and no two-dimensional symplectic source plane
can embed into this quotient.

This proves failure for every maximal simple type of every candidate
with these two invariants. It is stronger than failure of one explicit
matrix formula. It remains a statement about attachment inside the
actual maximal-simple-type Heisenberg carrier, not a proof that no other
arithmetic realization relation could work.

\hypertarget{dyadic_capacity-rank4-failure-is-also-forced-by-swan-and-torsion}{%
\subsection{5. Rank 4 failure is also forced by Swan and
torsion}\label{dyadic_capacity-rank4-failure-is-also-forced-by-swan-and-torsion}}

Let Pi be any actual supercuspidal representation of GL4(F) with torsion
number one and Swan exponent N=6h\_2, where h\_2 is odd. Its maximal
parameter field is totally ramified of degree~4, its order has period 4,
and its native level is N.

Here v\_2(N)=1. The defining-sequence/minimality argument of Block166
forces the terminal minimal approximation field E/F to have degree~2.
There is no possible intermediate degree between 2 and 4. The scalar
terminal case is excluded by 4 not dividing N, and a minimal
full-degree~4 parameter is excluded because gcd(N,4)=2. Thus the actual
degree pattern is forced to be 4→2.

Write c=−k\_0(beta,A), so 1\textless=c\textless N. Select the critical
approximation furnished by the construction in BK 2.4.1 and the remark
preceding 2.4.7. Because the original field is maximal, this
approximation can be selected with a simple maximal derived stratum in
B=End\_E(F\^{}4)≅M2(E). Equivalently one can replace its derived stratum
by an equivalent simple representative and use BK 2.2.2--3. The derived
field has degree~2 and ramification index 2 over E. A simple stratum
{[}B\_order,c,c−1,alpha{]} is minimal. Since the complete B-chain has
period 2 and the derived field is totally ramified, BK 1.4.14 forces

\begin{Verbatim}[fontsize=\small,breaklines=true,breakanywhere=true]
gcd(c,2)=1.
\end{Verbatim}

Therefore c is odd. It cannot be made even while preserving this degree
pattern and a genuine maximal derived field.

Let j=(c+1)/2. The actual orders are

\begin{Verbatim}[fontsize=\small,breaklines=true,breakanywhere=true]
J_order=O_K+Q^j+P^(N/2),
H_order=O_K+Q^j+P^(N/2+1).
\end{Verbatim}

The final low transition contributes nothing, while the high transition
has dimension 4−2=2. Thus every such Pi has

\begin{Verbatim}[fontsize=\small,breaklines=true,breakanywhere=true]
dim_F2 J^1/H^1=2,       dim eta=2.
\end{Verbatim}

The desired two-plane source carrier has dimension 4, so it cannot be a
sector of this eta. Requiring the self-pair conductor additionally fixes
c=3h\_1 through

\begin{Verbatim}[fontsize=\small,breaklines=true,breakanywhere=true]
a(Pi^vee×Pi)=15+2N+c,
\end{Verbatim}

but is unnecessary for the capacity failure.

\hypertarget{dyadic_capacity-rank8-the-three-invariants-do-not-force-the-good-profile}{%
\subsection{6. Rank 8: the three invariants do not force the good
profile}\label{dyadic_capacity-rank8-the-three-invariants-do-not-force-the-good-profile}}

The full 2→4→8 tower has quotient dimension 6 and enough capacity.
Nevertheless Swan, torsion number and self-pair conductor alone do not
force that degree pattern. Here is an actual skipped-degree arithmetic
construction with the same three numerical invariants and a smaller
quotient.

Fix increasing odd h\_1\textless h\_2\textless h\_3 and put

\begin{Verbatim}[fontsize=\small,breaklines=true,breakanywhere=true]
q=3h_3,       N=4q=12h_3,       c=h_1+4h_2.
\end{Verbatim}

Then c is odd and 0\textless c\textless N; indeed c\textless5h\_3, so
even 2c\textless N.

Construct the separable totally ramified quadratic extension

\begin{Verbatim}[fontsize=\small,breaklines=true,breakanywhere=true]
E=F(alpha),       alpha^2+alpha=t^(−q).
\end{Verbatim}

Put a=(q+1)/2 and z=t\^{}a alpha. Then z is a uniformizer and

\begin{Verbatim}[fontsize=\small,breaklines=true,breakanywhere=true]
z^2+t^a z+t=0.
\end{Verbatim}

The element alpha has E-valuation −q and is minimal over F. On E\^{}4,
gamma=alpha I\_4 is minimal of ambient level N in the complete
period-eight order A.

Now construct the separable Eisenstein quartic extension over E

\begin{Verbatim}[fontsize=\small,breaklines=true,breakanywhere=true]
L=E(e),       e^4+z e+z=0,
\end{Verbatim}

and let C be multiplication by e\^{}(−c) on L as a four-dimensional
E-space. Since c is odd, e\^{}(−c) is minimal, generates L/E, and
defines a simple maximal derived stratum {[}B\_order,c,c−1,C{]}.

Let P0 extract the coefficient of 1 on each E-coordinate in the basis
1,z. The normalized outer corestriction from the Artin--Schreier
right-inverse calculation has

\begin{Verbatim}[fontsize=\small,breaklines=true,breakanywhere=true]
s(P0)=t/z^2.
\end{Verbatim}

Consequently the explicit F-linear lift

\begin{Verbatim}[fontsize=\small,breaklines=true,breakanywhere=true]
b=(z^2/t) C P0
\end{Verbatim}

satisfies s(b)=C and has valuation −c.~Define the actual rank-eight
matrix

\begin{Verbatim}[fontsize=\small,breaklines=true,breakanywhere=true]
beta=alpha I_4+(z^2/t) C P0.
\end{Verbatim}

BK 2.2.3 applies because gamma is simple at level c\textless N and the
derived stratum is simple maximal. It proves that F{[}beta{]}/F is
totally ramified of degree~8, normalizes the complete chain, has leading
level N, and has critical exponent −c.~Its defining degree pattern is
8→2, skipping degree~4.

Since N is even and c is odd, the actual Heisenberg quotient has
dimension

\begin{Verbatim}[fontsize=\small,breaklines=true,breakanywhere=true]
(8−4)·1+(4−1)·0=4,
\end{Verbatim}

so eta has dimension 4. Yet its invariants are

\begin{Verbatim}[fontsize=\small,breaklines=true,breakanywhere=true]
Swan=N=12h_3,
torsion=1,
a(Pi^vee×Pi)=63+4N+3c
             =63+48h_3+12h_2+3h_1,
\end{Verbatim}

exactly the same as the full three-stage profile. Any maximal simple
character and its standard type completion give an actual supercuspidal
Pi with these values. The required full three-plane source carrier would
have dimension 8, larger than this eta.

This is an arithmetic counterexample to the proposed implication from
those three numerical invariants to adequate rank-eight Heisenberg
capacity. It is not an H5 counterexample, does not identify the skipped
host with the source-defined classical candidate, and does not claim
that it preserves the supplied full source characters.

\hypertarget{dyadic_capacity-which-repairs-have-actually-been-tested}{%
\subsection{7. Which repairs have actually been
tested}\label{dyadic_capacity-which-repairs-have-actually-been-tested}}

\hypertarget{dyadic_capacity-replacing-the-displayed-type-or-parameter-at-fixed-swan-and-torsion}{%
\subsubsection{Replacing the displayed type or parameter at fixed Swan
and
torsion}\label{dyadic_capacity-replacing-the-displayed-type-or-parameter-at-fixed-swan-and-torsion}}

At ranks 2 and 4 this fails by Sections~4--5 for every maximal simple type
satisfying those invariants. A new matrix search or a different critical
coefficient cannot manufacture the missing quotient dimensions.

At rank~8 the full tower supplies the good profile, but the three
invariants alone do not force it inside a separately selected candidate.
Section~6 proves the failure of that proposed inference with an explicit
actual host. A finer endo-class or ramification comparison would be
needed for that same-candidate assertion.

\hypertarget{dyadic_capacity-skipping-approximation-degrees-or-changing-critical-parity}{%
\subsubsection{Skipping approximation degrees or changing critical
parity}\label{dyadic_capacity-skipping-approximation-degrees-or-changing-critical-parity}}

At rank~2 there is no degree to skip. At rank~4 the level excludes both a
scalar minimal approximation and a full-degree minimal approximation.
The remaining quadratic derived step has odd critical exponent by
minimality. Thus neither skipping nor declaring that exponent even is a
valid repair at fixed invariants.

At rank~8 a degree skip is genuinely possible, but it decreases capacity
in the explicit example above.

\hypertarget{dyadic_capacity-scalar-twists-shallow-or-deep}{%
\subsubsection{Scalar twists, shallow or
deep}\label{dyadic_capacity-scalar-twists-shallow-or-deep}}

A scalar twist of the existing actual inducing representation does not
increase its dimension. The exact Block164 pushout transport identifies
the same quotient J/H and leaves its commutator pairing unchanged. Thus
it cannot turn an insufficient eta into a whole source carrier.

For an actual parameter, replacing beta by beta+aI with a∈F leaves its
field, centralizer and commutator map unchanged. A shallow scalar
addition does not change the leading level. A dominating scalar addition
can change that level, but it contributes no new noncentral Heisenberg
directions: in the defining sequence the added scalar stage has equal H
and J orders. This agrees with the exact representation-level twist
transport. Enlarging a displayed numerical level is therefore not a
capacity repair.

This paragraph does not assert that every arbitrary scalar character
twist is represented by one globally specified additive parameter. The
already proved group/representation transport suffices for its negative
capacity conclusion.

\hypertarget{dyadic_capacity-adding-residue-multiplicity-or-increasing-ambient-rank}{%
\subsubsection{Adding residue multiplicity or increasing ambient
rank}\label{dyadic_capacity-adding-residue-multiplicity-or-increasing-ambient-rank}}

The torsion-one condition forces a full-degree totally ramified
parameter field and residue multiplicity one. A host with an unramified
or residual multiplicity factor does not preserve that invariant.
Increasing the ambient rank changes the required source rank. Neither
maneuver closes the current fixed-rank construction.

\hypertarget{dyadic_capacity-using-a-larger-arithmetic-carrier-inside-the-same-representation}{%
\subsubsection{Using a larger arithmetic carrier inside the same
representation}\label{dyadic_capacity-using-a-larger-arithmetic-carrier-inside-the-same-representation}}

A larger compact subgroup representation, a genuinely different
arithmetic quotient, or an induction-stage carrier might have enough
dimension. Such an enlargement is not ruled out here. It would need an
actual finite Heisenberg action carrying the prescribed full source data
and a proved recovery in the unchanged candidate. Artificially adding
vector-space copies or transporting abstract operators to an induced
space does not provide that arithmetic attachment or prove its kernel
compatibility. No such repair is constructed in this note.

\hypertarget{dyadic_capacity-quadratic-refinements-and-marked-centers}{%
\subsection{8. Quadratic refinements and marked
centers}\label{dyadic_capacity-quadratic-refinements-and-marked-centers}}

For d\textgreater=3, the dimension result gives a symplectic embedding
of the source quotient into the actual host quotient. With C trivial, a
C×-central pushout can then be matched by its alternating commutator,
and Stone--von Neumann supplies the corresponding source sector. This is
the same finite scalar-pushout attachment contract used previously.

At p=2, that statement must not be promoted to an identification of
arbitrary preselected native quadratic refinements. The C× pushout
permits rescaling lifts by i and can identify the pushouts of the
dihedral and quaternion extraspecial planes, even though their native
square functions differ. A chosen source embedding can preserve that
source's full finite central extension and its character values. It does
not thereby identify a separately marked target mu\_2 extension or its
Arf invariant.

In particular, symplectic-isotropic reduction respects a prescribed
quadratic form only when the reduced subspace is totally singular for
that form. The retained \path{DYADIC_AND_WEAK_ATTACHMENT159.md} gives the
exact quadratic/Arf obstruction. Additional source or target marked
centers, chosen lifts, quadratic refinements and nontrivial stabilizer
actions therefore require their own comparison. They are neither erased
from the task nor supplied by a dimension inequality.

\hypertarget{dyadic_capacity-scope-and-source-provenance}{%
\subsection{9. Scope and source
provenance}\label{dyadic_capacity-scope-and-source-provenance}}

The strongest proved conclusion is exact capacity for the specified
actual full defining profile, together with invariant low-rank failures
and an explicit rank-eight counterexample to numerical profile forcing.
No universal impossibility for original Theorem~16.2 is asserted. In
particular this note supplies no complete original H1--H5 counterexample
and does not decide whether that endpoint requires this precise
whole-maximal-type-carrier realization.

Primary source directly reread: Bushnell--Kutzko, \emph{The admissible
dual of GL(N) via compact open subgroups}, 1.4.14--15; 2.2.2--3;
2.4.1--2 and the maximal-field remark immediately preceding 2.4.7;
3.1.5, 3.1.7--9; 3.2.1--4; 3.4.1; 6.2.5 and its following paragraph.

Authorized retained text:

\begin{Verbatim}[fontsize=\small,breaklines=true,breakanywhere=true]
closure_recovery/extracted/Paper_XI_Continuation158_General_Sectors_and_Arit
hmetic_Constraints/scratch158/information158/input_cache/BK_source.txt
SHA256 d69db0827256707a34cd2562af68b9de11fc889ab3f43535ec47b620d58ca9c8
\end{Verbatim}

Internal inputs reused at their stated scope:

\begin{itemize}
\tightlist
\item
  \path{attachment166/tower/RECURSIVE_TOWER166.md}: actual defining degrees,
  levels, filtered right inverse and separated grade count; its field
  constructor permits p=2.
\item
  \path{attachment166/tower/TRACE_ZERO_AS_TOWER166.md}, Sections~1--2:
  normalized corestriction and explicit right inverse. Their algebraic
  formulas apply at p=2; the odd-prime trace-zero assertion is not used.
\item
  \path{attachment166/capacity/LEADING_LAYER_CAPACITY166.md}: the leading
  subspace argument and the minimal terminal degree calculation.
\item
  \path{base_realization165/polarization/PROFILE_RIGIDITY165.md} and
  \path{closure_recovery/full161/scratch161/placement161/ADJOINT_CONDUCTOR_TEST161.md}:
  retained formal-degree/self-pair conductor identities and
  normalization.
\item
  \path{input_pair164/TWIST_TRANSPORT.md}: exact character-twist/pushout
  transport.
\item
  \path{closure_recovery/extracted/Paper_XI_Continuation159_Type_Completion_and_Sector_Recovery/scratch159/coverage159/DYADIC_AND_WEAK_ATTACHMENT159.md}:
  distinction between symplectic and marked quadratic reduction.
\end{itemize}

No full copyrighted source is redistributed here. No new external
correspondence theorem or H5 premise is introduced.

\clearpage
\section{Whole finite extensions and dyadic marking criteria}\label{app:dyadic_markings}

Evidence: direct reading of the original dyadic contract and ordinary
finite group proofs. This note does not prove a dyadic arithmetic field
construction, its simple-character existence, nontrivial stabilizer
attachment, H5, or a Lean closure. It identifies exactly what a supplied
actual host and sector must preserve, and proves the marked finite
attachment step for C=\{1\}.

\hypertarget{dyadic_markings-the-original-datum-includes-more-than-the-alternating-form}{%
\subsection{1. The original datum includes more than the alternating
form}\label{dyadic_markings-the-original-datum-includes-more-than-the-alternating-form}}

The original text
\path{closure_attempt162/inputs/original_reading.txt} states in
Definition~11.1, starting line 4911, that at p=2 a restricted bridge
additionally supplies compatible quadratic refinements and an
isomorphism of the resulting Heisenberg central extensions lifting the
given symplectic isometry. Sections~13.1--13.9 define the refinements and
their Arf/Gauss signs.

The retained visual appendix is particularly explicit, at
lines 6892--6947: with center \{I,-I\} fixed, squares of lifts carry the
quadratic bit. Multiplying a lift by i changes the particular finite
central-extension realization. The appendix says that allowing complex
scalar rescaling \textbf{without recording this distinction} erases data
that the bridge is required to retain.

The finite record must therefore retain at least:

\begin{itemize}
\tightlist
\item
  the labeled symplectic quotient and its pairing;
\item
  the supplied quadratic refinement q;
\item
  the specified finite central extension, its central
  embedding/character, and any marked lifts or specified extension
  isomorphism;
\item
  the resulting normalized Gauss sign, and all declared phase operators;
\item
  the complete original character values retained by the scalar pushout.
\end{itemize}

The last item alone does not retain the preceding finite-subgroup/lift
marking unless that marking is explicitly carried along. The raw text
does not require literal embeddings of the old arithmetic matrix groups
into a new type. Nor does it define a unique quadratic refinement
extracted automatically from an unmarked C× scalar extension of every
arithmetic type.

\hypertarget{dyadic_markings-general-finite-central-extension-criterion}{%
\subsection{2. General finite central-extension
criterion}\label{dyadic_markings-general-finite-central-extension-criterion}}

Let V be an elementary abelian 2-group with ordered basis
v\_1,\ldots,v\_m. Let E\_s and E\_t be central extensions of V by the
same abelian group Z. Fix the identity on Z and a quotient isomorphism
f.~Choose lifts a\_i of v\_i in E\_s and b\_i of f(v\_i) in E\_t. Write

\begin{Verbatim}[fontsize=\small,breaklines=true,breakanywhere=true]
a_i²=s_i,  b_i²=t_i,  [a_i,a_j]=c_ij, [b_i,b_j]=d_ij.
\end{Verbatim}

Every element has a unique normal form

\begin{Verbatim}[fontsize=\small,breaklines=true,breakanywhere=true]
z a_1^(epsilon_1) ... a_m^(epsilon_m), epsilon_i in {0,1},
\end{Verbatim}

and analogously on the target. The multiplication is completely
determined by the displayed squares and commutators.

\textbf{Criterion.} A center-fixing isomorphism lifting f exists
precisely when

\begin{Verbatim}[fontsize=\small,breaklines=true,breakanywhere=true]
c_ij=d_ij for every i,j,
s_i/t_i belongs to Z² for every i.                     (1)
\end{Verbatim}

Indeed, every such isomorphism must take a\_i to lambda\_i b\_i for some
lambda\_i in Z, forcing lambda\_i²=s\_i/t\_i. Conversely these equations
make the assignment a\_i -\textgreater{} lambda\_i b\_i preserve all
relations, and the normal forms prove it is bijective.

If the \textbf{particular lifts} a\_i and b\_i must correspond exactly,
the condition is s\_i=t\_i, with lambda\_i=1. If a full chosen section
is marked, rather than only generating lifts, its multiplication
coefficients must also correspond:

\begin{Verbatim}[fontsize=\small,breaklines=true,breakanywhere=true]
a(v)a(w)=c_s(v,w)a(v+w),
b(fv)b(fw)=c_t(fv,fw)b(f(v+w)),
\end{Verbatim}

and exact section preservation requires c\_s(v,w)=c\_t(fv,fw) for every
pair. Transporting the complete section under an existing group
isomorphism retains these identities automatically. Preserving just its
squares need not retain an additional independently specified section.

\hypertarget{dyadic_markings-fixed-two-element-center}{%
\subsubsection{2.1. Fixed two-element
center}\label{dyadic_markings-fixed-two-element-center}}

For Z=\{+1,-1\}, Z²=\{1\}. The lift square is independent of the lift
within its fiber and defines

\begin{Verbatim}[fontsize=\small,breaklines=true,breakanywhere=true]
a(v)²=(-1)^(q_s(v)).
\end{Verbatim}

The identity (ab)²=a²b²{[}b,a{]} gives the quadratic refinement law.
Criterion (1) becomes the exact requirement

\begin{Verbatim}[fontsize=\small,breaklines=true,breakanywhere=true]
q_t(f(v))=q_s(v) for every v.                          (2)
\end{Verbatim}

For an unspecified symplectic isomorphism between equal-dimensional
nondegenerate spaces, such an f exists exactly when their Arf invariants
agree. For a preassigned f, equality of Arf invariants alone is
insufficient: (2) must hold for that f.

\hypertarget{dyadic_markings-four-element-and-complex-centers}{%
\subsubsection{2.2. Four-element and complex
centers}\label{dyadic_markings-four-element-and-complex-centers}}

For Z=mu\_4, multiplication of a lift by i changes its square by -1. The
unmarked extension criterion preserves square classes modulo Z²=mu\_2.
In particular, squares +1 and -1 do not distinguish unmarked mu\_4
extensions. A separately marked finite mu\_2-central extension subgroup
(such as the full D8 or Q8), or prescribed lifts, still distinguishes
them.

For Z=C× every square ratio has a square root. Thus matching commutators
is sufficient for an unmarked center-fixing extension isomorphism. This
fact does not delete a marked finite subgroup from the objects being
compared.

\hypertarget{dyadic_markings-exact-d8q8-check}{%
\subsection{3. Exact D8/Q8 check}\label{dyadic_markings-exact-d8q8-check}}

In dimension two over C put

\begin{Verbatim}[fontsize=\small,breaklines=true,breakanywhere=true]
X=[[0,1],[1,0]],       Z=[[1,0],[0,-1]].
\end{Verbatim}

Then X²=Z²=I, XZ=-ZX, and {[}X,Z{]}=-I. The group generated by X and Z
has order 8 and is D8. Its refinement on V=F2² is

\begin{Verbatim}[fontsize=\small,breaklines=true,breakanywhere=true]
q_D(x,y)=xy,       Arf(q_D)=0,
\end{Verbatim}

with square table (+1,+1,+1,-1) at (0,e,f,e+f), and Gauss sum +2.

The generators iX,iZ also anticommute, but each squares to -I and their
product also squares to -I. They generate Q8, with

\begin{Verbatim}[fontsize=\small,breaklines=true,breakanywhere=true]
q_Q(x,y)=xy+x+y,   Arf(q_Q)=1,
\end{Verbatim}

square table (+1,-1,-1,-1), and Gauss sum -2.

Both C× pushouts identify with the same ambient matrix group

\begin{Verbatim}[fontsize=\small,breaklines=true,breakanywhere=true]
P={z X^a Z^b : z in C×, a,b in {0,1}}.
\end{Verbatim}

Nevertheless the marked subgroups

\begin{Verbatim}[fontsize=\small,breaklines=true,breakanywhere=true]
H_D=<X,Z>,          H_Q=<iX,iZ>
\end{Verbatim}

are different. No isomorphism fixing the two-element center can map H\_D
onto H\_Q: it would preserve q and hence the Arf invariant, which
differs. There is also no center-fixing automorphism of P carrying the
marked H\_D onto H\_Q. Their unmarked pushouts are isomorphic, but the
marked objects (P,H\_D) and (P,H\_Q) are not.

Conversely a Q8 source embeds perfectly into the same P by sending its
generators to iX and iZ. This preserves every source square and its
negative Gauss sign. It does not identify those generators with the
independently marked D8 lifts X,Z. This is the distinction needed for
the arithmetic sector.

\hypertarget{dyadic_markings-marked-scalar-pushout-attachment-with-c1}{%
\subsection{4. Marked scalar-pushout attachment with
C=\{1\}}\label{dyadic_markings-marked-scalar-pushout-attachment-with-c1}}

Assume an actual type, its full character theta, and its Heisenberg
quotient W have already been constructed. Let P\_W be its actual C×
scalar pushout, acting on its actual Heisenberg representation eta\_W.
Let V\_src be a finite nondegenerate symplectic F2-space of dimension 2s,
with its complete supplied dyadic marked extension data. For several
breaks use their labeled orthogonal sum, balanced central product, and
supplied product refinement.

Assume the following precise finite contract:

\begin{enumerate}
\def\labelenumi{\arabic{enumi}.}
\tightlist
\item
  The source full-character pushout P\_src has quotient V\_src and the
  same represented scalar center as P\_W.
\item
  The supplied finite central extensions and their marked lifts are
  retained as subgroups/elements of P\_src, or as already specified maps
  into P\_src. Any central character kernel already factored out must be
  recorded; a scalar pushout cannot recover invisible elements of that
  kernel.
\item
  The host has enough symplectic capacity, dim W\textgreater=2s, and the
  new attachment into this host has no independently fixed conflicting
  quadratic marking.
\item
  The declared stabilizer is C=\{1\}; the graph has one vertex.
\end{enumerate}

Then a complete marked attachment exists in the \textbf{same actual
candidate}.

\textbf{Proof.} Choose a symplectic embedding f:V\_src -\textgreater{} W
and let V be its image. The restriction P\_V of P\_W and P\_src have the
same commutator. Criterion (1) with center C× produces an isomorphism

\begin{Verbatim}[fontsize=\small,breaklines=true,breakanywhere=true]
F:P_src -> P_V
\end{Verbatim}

fixing every scalar and lifting f.~Carry every supplied finite subgroup
H\_src to the actual subgroup F(H\_src) of P\_V, and every marked lift a
to F(a). This operation does not erase square data:

\begin{Verbatim}[fontsize=\small,breaklines=true,breakanywhere=true]
F(a)²=F(a²),        [F(a),F(b)]=F([a,b]).              (3)
\end{Verbatim}

Since F fixes the scalar center, a²=(-1)\^{}q(v) implies
F(a)²=(-1)\^{}q(v). The transported refinement is therefore exactly q,
not a new choice with the same polar form. All marked section products,
specified extension maps, and finite group relations are preserved by
the same F. Consequently the Arf invariant and normalized Gauss sign are
unchanged.

Every original h in a full character domain still maps to theta\_i(h),
because F fixes the entire scalar center. Thus the dyadic marking is
retained together with the full original character values.

Write W=V perpendicular-sum U. Finite Stone--von Neumann gives

\begin{Verbatim}[fontsize=\small,breaklines=true,breakanywhere=true]
eta_W isomorphic to eta_src tensor eta_U.
\end{Verbatim}

Choose a line in eta\_U and its projector e. It lies in the actual
Heisenberg image algebra, and the marked source subgroup acts on e
eta\_W by precisely its supplied irreducible Heisenberg representation.
Its lift squares remain the operators in (3). The existing sector
recovery theorem then recovers the unchanged inducing representation
Lambda, and the same arithmetic candidate. There is no C-equivariance
equation to solve because C=\{1\}. QED.

This realizes an actual finite marked subgroup in the scalar extension
of the actual type, with actual represented operators that are complex
scalars times type operators. It does not assert that every marked
source element is the image of an unscaled element of the old or new
arithmetic matrix group. That stronger literal-embedding requirement is
absent from the present contract.

It also does not assert that the newly marked subgroup equals a
separately specified canonical native subgroup of the host. The marked
subgroup is part of the result and must be retained in the artifact.
Omitting it and recording only the unmarked P\_W would reproduce the
loss displayed by D8/Q8.

\hypertarget{dyadic_markings-if-the-host-refinement-is-independently-fixed}{%
\subsection{5. If the host refinement is independently
fixed}\label{dyadic_markings-if-the-host-refinement-is-independently-fixed}}

Suppose instead that the actual host arrives with a particular
nonsingular quadratic refinement q\_W and its actual finite mu\_2
central extension H\_W. An attachment into that marked extension
requires a quadratic isometric embedding, not merely a symplectic
embedding. There is a sharp capacity test.

\textbf{Theorem.} Let dim V\_src=2s and dim W=2R. With embeddings
otherwise free, a quadratic isometric embedding

\begin{Verbatim}[fontsize=\small,breaklines=true,breakanywhere=true]
(V_src,q_src) -> (W,q_W)
\end{Verbatim}

exists if and only if R\textgreater=s and, in the equal-capacity case
R=s,

\begin{Verbatim}[fontsize=\small,breaklines=true,breakanywhere=true]
Arf(q_src)=Arf(q_W).                                  (4)
\end{Verbatim}

If R\textgreater s, either source Arf type embeds.

\textbf{Proof.} Necessity of dimension is immediate. An
equal-dimensional embedding is an isometry, forcing (4). If
R\textgreater s, choose a nondegenerate quadratic space U of
dimension 2(R-s) with

\begin{Verbatim}[fontsize=\small,breaklines=true,breakanywhere=true]
Arf(q_U)=Arf(q_W)+Arf(q_src).
\end{Verbatim}

Both Arf types exist in every positive even dimension, by original
Proposition 13.7. The orthogonal sum V\_src perpendicular-sum U then has
the same dimension and Arf invariant as W. Original Standard Fact 13.4
gives an isometry to W; restriction to V\_src gives the desired
embedding. QED.

After this embedding is chosen, the preimage of its image in H\_W has
exactly q\_src. The fixed-two-element-center criterion of Section~2 gives
the extension isomorphism. Generating lifts with the specified squares
can be matched. Additional fully prescribed section elements still
require their exact product relations, as described in Section~2.

Thus \textbf{one spare symplectic plane absorbs the Arf mismatch},
without changing the host's own refinement. This conclusion remains
conditional on the host's actual marked refinement/finite extension
being supplied or constructed. It does not manufacture a canonical
arithmetic quadratic datum from its dimension.

For a preassigned embedding f, the missing equation is instead

\begin{Verbatim}[fontsize=\small,breaklines=true,breakanywhere=true]
q_W(f(v))=q_src(v) for all v.                          (5)
\end{Verbatim}

Its discrepancy q\_W composed with f minus q\_src is a linear
functional. Strict capacity allows a different quadratic embedding when
it is free; it does not authorize replacing a fixed original map, fixed
depth placement, or fixed transport marking.

\hypertarget{dyadic_markings-exact-outcome-for-the-dyadic-construction}{%
\subsection{6. Exact outcome for the dyadic
construction}\label{dyadic_markings-exact-outcome-for-the-dyadic-construction}}

Under the retained one-vertex, C=\{1\}, full scalar-pushout realization
contract, an actual host with sufficient symplectic capacity does
support every supplied finite dyadic marking, provided the transported
finite subgroup, lifts, refinement, and phase are explicitly kept. The
proof is Section~4, rather than an assertion that commutators determine
the dyadic sign.

If a particular host-native quadratic extension is fixed in advance, use
Section~5. Strict capacity suffices for its quadratic embedding; equal
capacity requires the Arf equation, and a preassigned embedding requires
(5). If only the type quotient is known and a specifically canonical
host quadratic datum is requested, constructing and comparing that datum
remains an arithmetic obligation. The scalar-pushout proof must not
silently claim to have done it.

No part of this note handles nontrivial prescribed C actions. In that
case quadratic embedding, full extension action, lift markings, and
chosen honest phase operators must all be equivariant; dimension and Arf
equality alone do not settle those equations.

\hypertarget{dyadic_markings-evidence-and-author-self-review}{%
\subsection{Evidence and author
self-review}\label{dyadic_markings-evidence-and-author-self-review}}

Read directly: original Definition~11.1, Sections~13 and 15, and visual
appendix pages V--22/V--23 in the retained raw file; final Block166
\path{PROOF166.md}, Section~2, and \path{LEADING_LAYER_CAPACITY166.md}, Sections~2.1
and 4. The extension criterion and D8/Q8 calculation are complete finite
proofs. The quadratic embedding theorem uses the Arf classification
already explicitly supplied by the original paper. The sector conclusion
reuses the retained exact recovery theorem and does not assert new
arithmetic existence.

Self-review checked that complex lift rescaling is never treated as an
automorphism changing squares within a fixed marked extension, that the
marked subgroup remains in the output, that preassigned maps differ from
freely chosen embeddings, and that fixed native refinements require
their own actual input. No independent review of this note or new Lean
execution is claimed here.

\clearpage
\section{The actual enlarged rank-two inducing carrier}\label{app:rank2}

Status: ordinary proof. The construction below supplies an actual
arithmetic finite Heisenberg pair and exact recovery of an already fixed
representation. It does not change the previously audited
zero-dimensional BK quotient at rank two. No Lean verification or H5
comparison is claimed.

\hypertarget{rank2-exact-setting-and-result}{%
\subsection{1. Exact setting and
result}\label{rank2-exact-setting-and-result}}

Let F=F\_2((t)) and G=GL2(F). Let {[}A,q,0,beta{]} be a minimal simple
stratum with K=F{[}beta{]} totally ramified of degree~2 and A its
complete period-two hereditary order. Assume

\begin{Verbatim}[fontsize=\small,breaklines=true,breakanywhere=true]
q=2m+1>=3,       m>=1.
\end{Verbatim}

Write R=rad(A), \(U_A^a=1+R^a\) for a\textgreater=1, and
\(U_K^a=1+\mathfrak p_K^a\). Let theta be a full simple character.
Since q is odd and the residue field is F2, the actual maximal simple
type group and character groups coincide:

\begin{Verbatim}[fontsize=\small,breaklines=true,breakanywhere=true]
J=J^1=H^1=U_K^1 U_A^(m+1).
\end{Verbatim}

Fix the full actual extension Lambda of theta to

\begin{Verbatim}[fontsize=\small,breaklines=true,breakanywhere=true]
N=K×J,
\end{Verbatim}

and the already fixed irreducible supercuspidal candidate

\begin{Verbatim}[fontsize=\small,breaklines=true,breakanywhere=true]
Pi=c-Ind_N^G Lambda.
\end{Verbatim}

Lambda is one-dimensional here. Its value on a uniformizer of K,
including the noncompact scalar normalization, is part of the fixed
data.

Define the larger compact open arithmetic subgroup

\begin{Verbatim}[fontsize=\small,breaklines=true,breakanywhere=true]
P=U_K^1 U_A^m.
\end{Verbatim}

Then the following statements hold.

\begin{enumerate}
\def\labelenumi{\arabic{enumi}.}
\item
  J is normal of index 2 in P, and P∩N=J.
\item
  E=Ind\_J\^{}P theta is an actual irreducible two-dimensional
  P-representation.
\item
  There is an explicitly defined normal subgroup H\_new⊂J of index 2 such
  that P/H\_new≅F2² and theta\_new=theta\textbar H\_new has
  nondegenerate commutator pairing on that quotient. E is its Heisenberg
  representation.
\item
  E embeds into the restriction of the same Pi, as the subspace of
  inducing functions supported on NP.
\item
  With X=c-Ind\_J\^{}G theta, a uniformizer varpi\_K and
  s=Lambda(varpi\_K), there is an exact quotient

\begin{Verbatim}[fontsize=\small,breaklines=true,breakanywhere=true]
Pi ≅ X/(T−s)X,
(Tf)(g)=f(varpi_K g).
\end{Verbatim}

  Induction in stages identifies X with c-Ind\_P\^{}G E. This gives an
  exact same-candidate recovery using the retained noncompact scalar,
  not merely a nonzero cyclic-generation assertion.
\end{enumerate}

The new Heisenberg pair is an actual pair of compact subgroups of
GL2(F), but is not asserted to be the BK \(H^1/J^1\) pair of
a maximal simple type. Whether the original H2 realization relation
allows this enlargement is a separate obligation.

\hypertarget{rank2-the-compact-enlargement-is-a-group-with-normal-subgroup-of-index-two}{%
\subsection{2. The compact enlargement is a group, with normal subgroup
of index
two}\label{rank2-the-compact-enlargement-is-a-group-with-normal-subgroup-of-index-two}}

K× normalizes A and every power of its radical, so U\_K\^{}1 normalizes
U\_A\^{}m. Therefore P is a subgroup and is compact open. It is
contained in U\_A\^{}1 because m\textgreater=1 and
\(U_K^1\subset U_A^1\).

The standard radical calculation gives

\begin{Verbatim}[fontsize=\small,breaklines=true,breakanywhere=true]
[U_A^a,U_A^b]⊂U_A^(a+b),       a,b>=1.
\end{Verbatim}

Indeed, for x∈R\^{}a and y∈R\^{}b, the difference (1+x)(1+y)−(1+y)(1+x)
is in R\^{}(a+b), and multiplication by unit inverses preserves that
ideal.

Consequently

\begin{Verbatim}[fontsize=\small,breaklines=true,breakanywhere=true]
[U_A^m,U_K^1]⊂U_A^(m+1),
[U_A^m,U_A^(m+1)]⊂U_A^(2m+1)⊂U_A^(m+1).
\end{Verbatim}

These facts, and normalization by U\_K\^{}1, show that J is normal in P.
They also show

\begin{Verbatim}[fontsize=\small,breaklines=true,breakanywhere=true]
[P,J]⊂U_A^(m+1).                                      (1)
\end{Verbatim}

To calculate the index, use

\begin{Verbatim}[fontsize=\small,breaklines=true,breakanywhere=true]
U_A^m∩J=U_K^m U_A^(m+1).
\end{Verbatim}

For if ku∈U\_A\^{}m with k∈U\_K\^{}1 and u∈U\_A\^{}(m+1), then
\(k\in K^\times\cap U_A^m=U_K^m\). The reverse inclusion is immediate.
The complete period-two A-grade \(R^m/R^{m+1}\) has
F2-dimension 2, and its field grade \(\mathfrak p_K^m/\mathfrak p_K^{m+1}\) has
dimension 1. Thus

\begin{Verbatim}[fontsize=\small,breaklines=true,breakanywhere=true]
[P:J]=[U_A^m:U_K^m U_A^(m+1)]=2.                      (2)
\end{Verbatim}

Finally, if kj∈P with k∈K× and j∈J, then \(k\in Pj^{-1}\subset U_A^1\).
Hence k∈U\_K\^{}1 and kj∈J. Therefore

\begin{Verbatim}[fontsize=\small,breaklines=true,breakanywhere=true]
P∩N=J.                                                (3)
\end{Verbatim}

These arguments use the actual group intersections. They do not replace
P or J by a vector-space surrogate.

\hypertarget{rank2-irreducibility-and-the-conjugate-character-ratio}{%
\subsection{3. Irreducibility and the conjugate-character
ratio}\label{rank2-irreducibility-and-the-conjugate-character-ratio}}

For the minimal simple character, the primary BK intertwining theorem
gives

\begin{Verbatim}[fontsize=\small,breaklines=true,breakanywhere=true]
I_G(theta)=J K×J=N.
\end{Verbatim}

Here K× normalizes J and theta. In particular, for g∈P\textbackslash J,
(3) says g does not intertwine theta. Because J is normal in P, this
means that the two characters

\begin{Verbatim}[fontsize=\small,breaklines=true,breakanywhere=true]
theta,
theta^g(j)=theta(gjg^−1)
\end{Verbatim}

are distinct. The index-two induction E=Ind\_J\^{}P theta is therefore
irreducible. This also follows directly from the two distinct J-weight
lines, which are interchanged by g.

Define the actual character of J

\begin{Verbatim}[fontsize=\small,breaklines=true,breakanywhere=true]
lambda(j)=theta^g(j)/theta(j).
\end{Verbatim}

It is nontrivial. By (1),

\begin{Verbatim}[fontsize=\small,breaklines=true,breakanywhere=true]
lambda(j)=theta(gjg^−1j^−1),
\end{Verbatim}

and the argument of theta lies in U\_A\^{}(m+1). The full simple
character has exact restriction

\begin{Verbatim}[fontsize=\small,breaklines=true,breakanywhere=true]
theta(1+x)=psi_F(tr_F(beta x)),       x∈R^(m+1).
\end{Verbatim}

Since F has characteristic 2, its additive characters take values in
\{1,−1\}. Consequently

\begin{Verbatim}[fontsize=\small,breaklines=true,breakanywhere=true]
lambda(J)={1,−1}.                                    (4)
\end{Verbatim}

No assertion about shallow values of theta is used here. Theta may have
a larger finite image on J. Only its ratio with its actual conjugate is
proved to have order 2.

As g²∈J, conjugation by g² fixes the character theta. Thus
\(\lambda^g=\lambda^{-1}=\lambda\). It follows that

\begin{Verbatim}[fontsize=\small,breaklines=true,breakanywhere=true]
H_new=ker(lambda)
\end{Verbatim}

is normal in P. It has index 2 in J and index 4 in P.

\hypertarget{rank2-the-actual-finite-heisenberg-pair}{%
\subsection{4. The actual finite Heisenberg
pair}\label{rank2-the-actual-finite-heisenberg-pair}}

Use the function model of E with f(jp)=theta(j)f(p) and right
translation by P. Coordinates are f(1),f(g). Then for j∈J,

\begin{Verbatim}[fontsize=\small,breaklines=true,breakanywhere=true]
E(j)=diag(theta(j),theta^g(j))
    =theta(j) diag(1,lambda(j)),
\end{Verbatim}

while

\begin{Verbatim}[fontsize=\small,breaklines=true,breakanywhere=true]
E(g)=[[0,1],[theta(g²),0]].                            (5)
\end{Verbatim}

Therefore E(g)²=theta(g²)I. Choose j\_0∈J with lambda(j\_0)=−1, and put

\begin{Verbatim}[fontsize=\small,breaklines=true,breakanywhere=true]
Z=theta(j_0)^−1 E(j_0)=diag(1,−1),
X=E(g).
\end{Verbatim}

The exact matrix identity is

\begin{Verbatim}[fontsize=\small,breaklines=true,breakanywhere=true]
XZX^−1=−Z.                                           (6)
\end{Verbatim}

The projective image of J has order 2, generated by Z, while X is
off-diagonal and contributes an independent projective involution. Thus
the projective image of P is exactly C2×C2, with nontrivial commutator
(6). Its projective kernel is precisely H\_new. In particular

\begin{Verbatim}[fontsize=\small,breaklines=true,breakanywhere=true]
P/H_new≅F2².
\end{Verbatim}

On H\_new, E is the scalar theta\_new=theta\textbar H\_new. This
character is P-invariant: conjugation by J fixes theta and conjugation
by g changes theta by lambda, which is trivial on H\_new. The resulting
alternating form

\begin{Verbatim}[fontsize=\small,breaklines=true,breakanywhere=true]
(pH_new,p'H_new)↦theta_new([p,p'])
\end{Verbatim}

is nondegenerate, since its two displayed generators have commutator −1.
Therefore (H\_new,P,theta\_new) is an actual arithmetic finite
Heisenberg pair, and its actual irreducible carrier E has dimension 2.

All representations here have finite image on the compact group P: theta
has open kernel, the induction has finite dimension, and an open normal
core of that kernel acts trivially. Thus no infinite projective-image
issue is being suppressed.

This proof does not require rescaling the actual matrix g inside GL2(F).
Its represented square theta(g²) can be an arbitrary retained scalar.
Rescaling X is only allowed later in a C× pushout if that is the
attachment contract.

\hypertarget{rank2-explicit-embedding-into-the-same-fixed-candidate}{%
\subsection{5. Explicit embedding into the same fixed
candidate}\label{rank2-explicit-embedding-into-the-same-fixed-candidate}}

Realize Pi as functions F:G→C with

\begin{Verbatim}[fontsize=\small,breaklines=true,breakanywhere=true]
F(nx)=Lambda(n)F(x),
\end{Verbatim}

compact support modulo N, and right translation by G. For f∈E define
i(f) by

\begin{Verbatim}[fontsize=\small,breaklines=true,breakanywhere=true]
i(f)(np)=Lambda(n)f(p),       n∈N, p∈P,
\end{Verbatim}

and zero off NP. This is well-defined because N∩P=J and
Lambda\textbar J=theta. The set NP is a finite union of N-cosets, by
{[}P:J{]}=2, so the support is compact modulo N. The map i is
P-equivariant and injective.

Its image is exactly the two-dimensional P-stable subspace of Pi
consisting of inducing functions supported on NP. Thus the enlarged
carrier is already inside the given Pi. No new candidate is chosen, and
all noncompact values of Lambda remain fixed.

Since Pi is irreducible and i(E) is nonzero, its G-span is Pi. More
precisely, there is the canonical G-equivariant quotient

\begin{Verbatim}[fontsize=\small,breaklines=true,breakanywhere=true]
Q_E:c-Ind_P^G E→Pi,
Q_E(F)=sum_(Px∈P\G) Pi(x^−1)i(F(x)).                 (7)
\end{Verbatim}

The sum is finite and independent of representatives. If x is replaced
by px, then F(px)=E(p)F(x) and P-equivariance of i cancels the
additional Pi(p\^{}−1). Formula (7) is nonzero and hence onto. The next
section identifies its kernel exactly.

\hypertarget{rank2-exact-recovery-by-the-noncompact-hecke-eigenvalue}{%
\subsection{6. Exact recovery by the noncompact Hecke
eigenvalue}\label{rank2-exact-recovery-by-the-noncompact-hecke-eigenvalue}}

Let X=c-Ind\_J\^{}G theta, in the same left-covariant/right-translation
convention. Induction in stages gives an explicit isomorphism

\begin{Verbatim}[fontsize=\small,breaklines=true,breakanywhere=true]
c-Ind_P^G E → X,
F ↦ [x↦F(x)(1)].                                     (8)
\end{Verbatim}

The inverse sends f∈X to the E-valued function F(x)(p)=f(px). Both
covariance statements follow directly from the definitions, and the
finite index {[}P:J{]} ensures the support conditions agree.

Fix a uniformizer varpi=varpi\_K. Because K× normalizes theta, the
operator

\begin{Verbatim}[fontsize=\small,breaklines=true,breakanywhere=true]
(Tf)(x)=f(varpi x)                                   (9)
\end{Verbatim}

is a well-defined invertible G-equivariant operator on X. Put
s=Lambda(varpi), which is nonzero. Since k\_K=F2 and J contains
U\_K\^{}1,

\begin{Verbatim}[fontsize=\small,breaklines=true,breakanywhere=true]
N/J≅K×/U_K^1≅Z,
\end{Verbatim}

with generator varpi. Define

\begin{Verbatim}[fontsize=\small,breaklines=true,breakanywhere=true]
Q:X→Pi,
(Qf)(x)=sum_(ell∈Z) s^(−ell) f(varpi^ell x).          (10)
\end{Verbatim}

For each x the sum is finite, since f has compact support modulo the
open subgroup J. Equation (10) satisfies

\begin{Verbatim}[fontsize=\small,breaklines=true,breakanywhere=true]
(Qf)(varpi x)=s(Qf)(x),
(Qf)(jx)=theta(j)(Qf)(x),       j∈J,
\end{Verbatim}

and therefore the required N-covariance. Its support is compact modulo
N. It commutes with G-right translation and satisfies

\begin{Verbatim}[fontsize=\small,breaklines=true,breakanywhere=true]
QT=sQ.                                               (11)
\end{Verbatim}

It is onto: a function in Pi has finite support modulo N, and on each
occupied N-coset its value can be lifted to a function supported on one
J-coset. Extend those finitely many choices using theta.

To compute the kernel, decompose X as the algebraic direct sum of its
subspaces supported on the left N-cosets in G. Each such summand is the
space of finitely supported sequences on N/J≅Z, with the fixed theta
covariance. The operator T is a bilateral shift. On that summand, (10)
is evaluation of its Laurent polynomial at the nonzero scalar s, up to
the harmless convention of replacing the Laurent variable by its
inverse. The kernel is therefore the principal ideal generated by T−s.
The elementary identity

\begin{Verbatim}[fontsize=\small,breaklines=true,breakanywhere=true]
ker(C[z,z^−1]→C, z↦s)=(z−s)
\end{Verbatim}

proves, summand by summand, that

\begin{Verbatim}[fontsize=\small,breaklines=true,breakanywhere=true]
ker Q=(T−s)X.
\end{Verbatim}

All vectors occupy finitely many such summands, so no completion or
closure is needed. Thus

\begin{Verbatim}[fontsize=\small,breaklines=true,breakanywhere=true]
Pi ≅ X/(T−s)X ≅ (c-Ind_P^G E)/(T−s)(c-Ind_P^G E),    (12)
\end{Verbatim}

where T is transported through (8).

Under (8), formula (10) agrees with (7): on the standard copy of E
supported on P, it gives exactly i, and G-equivariance then gives
equality everywhere. Thus (12) identifies the kernel of the canonical
recovery map, not merely the kernel of an unrelated surjection.

The same description can be phrased as evaluation of the type Hecke
algebra at the character determined by Lambda. Its intertwining support
is N, each supported intertwining space is one-dimensional, and the
retained extension Lambda normalizes the Laurent generator. The explicit
formulas above avoid any ambiguity about convolution orientation.

The scalar s is essential. E alone, being compact data, does not
determine it. Changing s can change the unramified normalization of the
candidate. The construction retains the original Lambda and s, so it
recovers the exact original Pi rather than an unspecified member of its
inertial class.

\hypertarget{rank2-full-scalar-pushout-attachment-and-remaining-marking-obligations}{%
\subsection{7. Full scalar-pushout attachment and remaining marking
obligations}\label{rank2-full-scalar-pushout-attachment-and-remaining-marking-obligations}}

Form the C× pushout of the actual pair (H\_new,P,theta\_new). It is a
Heisenberg extension of the actual symplectic plane P/H\_new. Any
supplied source plane with the same chosen alternating form has a
center-preserving pushout isomorphism into it. Every source
full-character value is then preserved in the center, as in the retained
full-character attachment theorem. Stone--von Neumann identifies its
source carrier with the actual two-dimensional E.

Together with the embedding i and exact quotient (12), this provides an
actual arithmetic carrier and recovery inside the same fixed candidate.
It closes the carrier-capacity failure of the old rank-two eta by
changing the compact arithmetic pair, not by increasing the old eta or
discarding its dimension calculation.

At p=2, a C× pushout does not automatically identify a separately marked
native mu\_2 extension or quadratic refinement. The source finite
central extension can be retained through its chosen embedding, while a
required comparison with the target's preselected square function
remains an additional question. No arbitrary native Arf invariant is
asserted here.

Nontrivial stabilizer operators, transport graphs, and original H2's
requirement that a target be a particular BK type-side quotient are
likewise separate. In particular, the theorem proves that the new pair
is an actual arithmetic finite Heisenberg pair, but does not rename it
as the unchanged maximal-simple-type pair. No new hypothesis is silently
added to original Theorem~16.2.

\hypertarget{rank2-scope-smallest-boundary-and-sources}{%
\subsection{8. Scope, smallest boundary, and
sources}\label{rank2-scope-smallest-boundary-and-sources}}

The inequality q\textgreater=3 matters. If q=1, then m=0 and the
displayed enlargement is not the group used above. Over residue field F2
the complete Iwahori unit group equals U\_A\^{}1, and a naive level-zero
replacement does not give the asserted index-two enlargement. No q=1
claim is made. The source levels q=3h with positive odd h satisfy the
required q\textgreater=3.

The proof uses the actual one-dimensional Lambda because the maximal
parameter field has full degree and residue field F2. It does not apply
the same scalar argument automatically to a higher-rank inducing
representation.

Primary BK passages: 3.1.7--8 (odd minimal H/J equality); 3.2.1--2 (deep
restriction and K× normalization); 3.3.2 and 5.1.1/5.1.8 (intertwining);
and 6.2.1--5 (the fixed maximal-type extension and compact induction).
The subgroup, matrix, Mackey-support, and Laurent-evaluation
calculations are proved explicitly above. No external correspondence
theorem is used in this enlargement.

Retained primary text:

\begin{Verbatim}[fontsize=\small,breaklines=true,breakanywhere=true]
closure_recovery/extracted/Paper_XI_Continuation158_General_Sectors_and_Arit
hmetic_Constraints/scratch158/information158/input_cache/BK_source.txt
SHA256 d69db0827256707a34cd2562af68b9de11fc889ab3f43535ec47b620d58ca9c8
\end{Verbatim}

The previously reviewed negative statement concerned eta of the original
BK \(J^1/H^1\). That quotient remains zero. The new positive
statement concerns a different, explicitly constructed arithmetic pair
inside the same ambient representation, with exact
noncompact-normalization-aware recovery.

\clearpage
\section{The actual enlarged rank-four inducing carrier}\label{app:rank4}


\hypertarget{rank4-result-and-boundary}{%
\subsection{Result and boundary}\label{rank4-result-and-boundary}}

For the actual rank-four Artin--Schreier candidate of Block167, one can
enlarge the compact group J1 to an actual compact subgroup P with

\begin{Verbatim}[fontsize=\small,breaklines=true,breakanywhere=true]
J1 normal in P,  [P:J1]=2,
E=Ind_(J1)^P(eta) irreducible of dimension 4.
\end{Verbatim}

The representation E occurs inside the restriction of the same actual
supercuspidal candidate pi. There is an explicit normal subgroup H0 of P
with {[}P:H0{]}=16 such that P/H0 is elementary abelian and the
commutator form supplied by the actual character on H0 is nondegenerate.
Thus the carrier has an actual four-dimensional binary symplectic
quotient. No abstract transfer of matrix operators is used to create
this action.

The finite image of P has scalar center theta(H0), which can be larger
than \{+1,-1\}. It is literally an extraspecial group if this scalar
center has order two. That additional assertion is not assumed or proved
here. The construction supplies an actual finite Heisenberg central
extension with nondegenerate projective quotient F\_2\^{}4. Matching an
independently prescribed quadratic phase is also a separate assertion.

\hypertarget{rank4-actual-arithmetic-groups-and-notation}{%
\subsection{1. Actual arithmetic groups and
notation}\label{rank4-actual-arithmetic-groups-and-notation}}

Let F=F\_2((t)), and take the audited rank-four constructor with odd
0\textless h\_1\textless h\_2. Write

\begin{Verbatim}[fontsize=\small,breaklines=true,breakanywhere=true]
q_i=3h_i,  n=2q_2,  a=n/2=q_2,
r=q_1=2m+1,  m=(r-1)/2>=1.
\end{Verbatim}

We have a\textgreater r, so n\textgreater2r. Let A be the complete
period-four hereditary order, with radical P\_A. Let gamma be the actual
degree-two outer minimal approximation, E=F{[}gamma{]},
B=End\_E(F\^{}4)=M\_2(E), B0=A cap B, and Q=rad(B0). The actual final
parameter is beta=gamma+delta, and

\begin{Verbatim}[fontsize=\small,breaklines=true,breakanywhere=true]
s(delta)=c,
\end{Verbatim}

where c is the minimal degree-two Artin--Schreier parameter over E, of
B0-level r. The final field K=F{[}beta{]} is totally ramified of degree
four; no embedding E subset K is asserted.

Use \(U_A^j=1+P_A^j\) and \(U_B^j=1+Q^j\).
BK 3.1.7--9 give the actual groups

\begin{Verbatim}[fontsize=\small,breaklines=true,breakanywhere=true]
H=H1(beta)=U_K^1 U_B^(m+1) U_A^(a+1),
J=J1(beta)=U_K^1 U_B^(m+1) U_A^a.
\end{Verbatim}

In particular J/H has dimension two over F\_2. Let theta be the actual
full simple character constructed for beta, and let eta be its unique
Heisenberg representation on J, of dimension two. The argument below
works with this supplied actual theta; it does not replace its values.

Choose the outer minimal simple character theta\_gamma from its defining
sequence. On H cap U\_A\^{}(m+1), the recursion says

\begin{Verbatim}[fontsize=\small,breaklines=true,breakanywhere=true]
theta=theta_gamma psi_delta.
\end{Verbatim}

Moreover H is contained in H1(gamma)=U\_B\^{}1 U\_A\^{}(a+1). The
character theta\_gamma is invariant under U\_B\^{}1: on U\_B\^{}1 it
factors through the E-determinant, and on U\_A\^{}(a+1) it is the
gamma-additive character.

\hypertarget{rank4-the-actual-final-field-has-the-needed-graded-line}{%
\subsection{2. The actual final field has the needed graded
line}\label{rank4-the-actual-final-field-has-the-needed-graded-line}}

For every integer j\textgreater=1 and k in K cap P\_A\^{}j, the
minimal-approximation commutator estimate gives

\begin{Verbatim}[fontsize=\small,breaklines=true,breakanywhere=true]
k=b+e,  b in Q^j,  e in P_A^(j+n-r).              (1)
\end{Verbatim}

Indeed {[}gamma,k{]}=-{[}delta,k{]} belongs to P\_A\^{}(j-r), and the
minimal corestriction exact sequence for gamma solves its commutator
with an error of grade j+n-r. The remaining difference centralizes gamma
and therefore belongs to B. This is the usual minimal case of BK 1.4.10.

Applying s to {[}beta,k{]}=0 and using (1) gives

\begin{Verbatim}[fontsize=\small,breaklines=true,breakanywhere=true]
[c,b] in Q^(j+n-2r).
\end{Verbatim}

Hence the leading B-grade of b commutes with the leading grade of c. Let
R be the complete period-two E-chain shift, with R\^{}2=z I, where z is
the chosen uniformizer of E. In the constructor's basis

\begin{Verbatim}[fontsize=\small,breaklines=true,breakanywhere=true]
c congruent to R^(-r) modulo Q^(1-r).
\end{Verbatim}

Because r is odd, the centralizing subspace on every B-grade j is the
one-dimensional line spanned by R\^{}j. On the other hand K is a totally
ramified maximal field normalizing the complete A-chain, so (K cap
P\_A\^{}j)/(K cap P\_A\^{}(j+1)) has dimension one over F\_2. Estimate
(1) is injective on the leading grade, so its image is exactly that
centralizing line.

This proves the required graded assertion for the actual K, without
replacing it by the derived model E{[}c{]} or assuming nested fields.

\hypertarget{rank4-normality-and-index-of-the-enlarged-group}{%
\subsection{3. Normality and index of the enlarged
group}\label{rank4-normality-and-index-of-the-enlarged-group}}

Set

\begin{Verbatim}[fontsize=\small,breaklines=true,breakanywhere=true]
P=J U_B^m.
\end{Verbatim}

We first prove this is a group. The group U\_B\^{}m normalizes
U\_B\^{}(m+1) and every U\_A\^{}j. For k in U\_K\^{}1, (1) approximates
k by an element of U\_B\^{}1 with error in U\_A\^{}(n-r+1). Therefore

\begin{Verbatim}[fontsize=\small,breaklines=true,breakanywhere=true]
[U_B^m,U_K^1]
   subset U_B^(m+1) U_A^(m+n-r+1)
   subset H,
\end{Verbatim}

where the last inclusion follows from a\textgreater r=2m+1. Thus
U\_B\^{}m normalizes both H and J, and P is an actual compact subgroup
with J normal in P.

The quotient \(U_B^m/U_B^{m+1}\) has dimension two. By
Section~2 the image of J cap U\_B\^{}m on this grade is exactly its
centralizing line, of dimension one. For the converse inclusion needed
here, the B-approximation of a K-element of grade m differs from it by
grade m+n-r\textgreater a; that error lies in U\_A\^{}a subset J. Thus
the whole centralizing line really occurs in the intersection, not only
in an abstract associated graded space. Consequently

\begin{Verbatim}[fontsize=\small,breaklines=true,breakanywhere=true]
[P:J]=[U_B^m:J cap U_B^m]=2.                    (2)
\end{Verbatim}

For a concrete representative choose D=diag(1,0) in B and put

\begin{Verbatim}[fontsize=\small,breaklines=true,breakanywhere=true]
x=R^m D,  u=1+x.
\end{Verbatim}

Its leading grade is outside the centralizing line, so u represents the
nontrivial coset of P/J. Since \(u^2=1+x^2\) in
characteristic two, we have u\^{}2 in U\_B\^{}(2m) subset U\_B\^{}(m+1)
subset H.

\hypertarget{rank4-the-actual-character-has-an-orbit-of-length-two}{%
\subsection{4. The actual character has an orbit of length
two}\label{rank4-the-actual-character-has-an-orbit-of-length-two}}

Define the character of H

\begin{Verbatim}[fontsize=\small,breaklines=true,breakanywhere=true]
lambda(h)=theta(u h u^(-1))/theta(h).
\end{Verbatim}

For h in H the preceding normality argument gives {[}u,h{]} in H cap
U\_A\^{}(m+1). Since theta\_gamma is invariant under u,

\begin{Verbatim}[fontsize=\small,breaklines=true,breakanywhere=true]
lambda(h)=psi_delta([u,h]) in {+1,-1}.          (3)
\end{Verbatim}

Thus lambda has order dividing two, and its values are actual values of
the supplied simple character. It is J-invariant, since both theta and
its u-conjugate are J-invariant.

It is nontrivial. Take y=R\^{}(m+1)D and h=1+y in U\_B\^{}(m+1). The
leading commutator in (3), after corestriction, has value

\begin{Verbatim}[fontsize=\small,breaklines=true,breakanywhere=true]
psi_E(tr_(M_2(E)/E)(c[x,y]))=-1.
\end{Verbatim}

Here the commutator remainder has B-grade at least 3m+1, which becomes
positive after multiplication by c of grade -2m-1. On the leading
grades, writing X=R\^{}m diag(A,B) and Y=R\^{}(m+1)diag(C,D), one
obtains

\begin{Verbatim}[fontsize=\small,breaklines=true,breakanywhere=true]
tr(R^(-2m-1)[X,Y])=(A+B)(C+D) in F_2.
\end{Verbatim}

Both sums equal one for the chosen D=diag(1,0). The normalized psi\_E is
nontrivial on the residue field and trivial on its maximal ideal. For
this AS normalization one can check the sign directly from
psi\_E(1)=psi\_F(Tr\_(E/F)(alpha))=psi\_F(1)=-1.

Hence theta\^{}u differs from theta, eta\^{}u differs from eta, and
Clifford theory applied to (2) gives

\begin{Verbatim}[fontsize=\small,breaklines=true,breakanywhere=true]
E_carrier=Ind_J^P(eta) irreducible, dim(E_carrier)=4.       (4)
\end{Verbatim}

\hypertarget{rank4-an-actual-nondegenerate-symplectic-quotient-of-dimension-four}{%
\subsection{5. An actual nondegenerate symplectic quotient of dimension
four}\label{rank4-an-actual-nondegenerate-symplectic-quotient-of-dimension-four}}

Put H0=ker(lambda) inside H, and theta0=theta restricted to H0. In the
two-coset induction model, h in H acts by

\begin{Verbatim}[fontsize=\small,breaklines=true,breakanywhere=true]
diag(theta(h) I_2, theta^u(h) I_2).
\end{Verbatim}

Thus it acts scalarly precisely when h belongs to H0. We check that H0
is the entire scalar preimage and that its quotient is symplectic.

First, lambda is invariant under u. Indeed u\^{}2 belongs to J and theta
is J-invariant, so \(\lambda^u=\lambda^{-1}=\lambda\) by (3).
Consequently H0 is normal in P, and theta0 is P-invariant.

Second, lambda is trivial on U\_A\^{}(a+1). There the correction formula
is valid, gamma commutes with u, and the delta contribution to a
commutator has grade at least a+1+m-r=a-m\textgreater0. Therefore it
lies in the kernel of the additive character.

Write J=H U\_A\^{}a. Since lambda is a character and J-invariant, it
kills {[}H,H{]} and {[}H,U\_A\^{}a{]}. It also kills
\([U_A^a,U_A^a]\), which lies in U\_A\^{}(2a). Hence
lambda({[}J,J{]})=1. It kills every square in J as well: squares of H
have lambda-value one, squares of U\_A\^{}a lie in U\_A\^{}(2a), and the
mixed commutator contributions have just been handled.

Finally lambda({[}u,H{]})=1 by u-invariance, and {[}u,U\_A\^{}a{]}
belongs to U\_A\^{}(a+m) subset U\_A\^{}(a+1), so lambda({[}u,J{]})=1.
Also lambda(u\^{}2)=1 because u commutes with u\^{}2. It follows that

\begin{Verbatim}[fontsize=\small,breaklines=true,breakanywhere=true]
[P,P] subset H0,  g^2 in H0 for every g in P,
[P:H0]=[P:J][J:H][H:H0]=2*4*2=16.
\end{Verbatim}

Therefore P/H0 is an actual elementary abelian group of dimension four.
If an element of J acts scalarly on the carrier, its first eta-block is
scalar, so it belongs to H by the original nondegenerate Heisenberg
quotient J/H. It then belongs to H0. An element outside J exchanges the
two coset summands and cannot act scalarly. Thus H0 is precisely the
scalar preimage.

The pairing

\begin{Verbatim}[fontsize=\small,breaklines=true,breakanywhere=true]
(gH0,kH0) |-> theta0([g,k])
\end{Verbatim}

is alternating and well defined. Since the projective quotient has
exponent two, its values lie in \{+1,-1\}. It is nondegenerate: an
element in its radical commutes with the entire irreducible carrier
image and therefore acts scalarly by Schur's lemma, forcing it into H0.

This is the actual four-dimensional binary Heisenberg quotient sought
for capacity. Its operators are supplied by elements of P.

\hypertarget{rank4-the-carrier-occurs-in-the-same-supercuspidal-candidate}{%
\subsection{6. The carrier occurs in the same supercuspidal
candidate}\label{rank4-the-carrier-occurs-in-the-same-supercuspidal-candidate}}

Let pi=c-Ind\_N\^{}(GL\_4(F))(Lambda) be the existing actual candidate,
with N=K\^{}times J. The residue field is F\_2, so the compact extended
type restricts to eta on J. Since P is contained in U\_A\^{}1,

\begin{Verbatim}[fontsize=\small,breaklines=true,breakanywhere=true]
N cap P=J.
\end{Verbatim}

Indeed an element k j in P with k in K\^{}times has k in U\_A\^{}1,
hence k in U\_K\^{}1 subset J. The identity double-coset term in the
restriction of compact induction to P is therefore

\begin{Verbatim}[fontsize=\small,breaklines=true,breakanywhere=true]
Ind_J^P(eta)=E_carrier.
\end{Verbatim}

It is realized concretely by the functions supported on N P. Thus the
four-dimensional carrier occurs in the same pi, with its already fixed
full type character and noncompact extension scalar. No new arithmetic
candidate, beta extension, or inverse correspondence is used.

\hypertarget{rank4-exact-limits-finite-image-versus-literal-extraspecial-group}{%
\subsection{7. Exact limits: finite image versus literal extraspecial
group}\label{rank4-exact-limits-finite-image-versus-literal-extraspecial-group}}

The actual finite image is P/ker(theta0), because H0 is the scalar
preimage. Its center is the scalar cyclic 2-group theta0(H0), and its
commutator subgroup has order two. Its quotient by its center is
F\_2\^{}4.

If theta0(H0) has order two, this finite image is extraspecial of order
32. If its scalar center is larger, quotienting arbitrary central
elements need not preserve its commutator: in a cyclic 2-group the
order-two subgroup can lie inside the squares. Thus a literal
extraspecial quotient is not justified solely by the calculations above.

The stronger statement proved is the one with full actual character: an
actual compact subgroup, actual induced carrier, actual scalar subgroup
and nondegenerate actual commutator pairing. Passing to the usual full
scalar pushout retains all theta0-values. It does not prove that an
independently supplied quadratic refinement, stabilizer action, or
source phase agrees with this carrier. Those requirements must be
checked separately. No H5 or universal claim~66 closure is asserted.

\hypertarget{rank4-provenance}{%
\subsection{Provenance}\label{rank4-provenance}}

Uses the audited dyadic parameter/character construction, BK 1.4.10,
BK 3.1.7--9, and the defining-sequence character formula. The source
normalization and all field/chain data are the actual ones in
\path{dyadic167/tower/DYADIC_AS_TOWER167.md}. This note changes neither that
proof nor its recorded exact computation. Clifford induction and the
identity compact-induction double-coset term are used in their ordinary
finite-index forms, with hypotheses proved above.

\clearpage
\section{Rank-four normalization and induction in stages}\label{app:rank4_stages}

This is a corollary of the enlarged-carrier construction in
Appendix~\ref{app:rank4}.

Use that note's actual groups and notation:

\begin{Verbatim}[fontsize=\small,breaklines=true,breakanywhere=true]
n=2a,  r=2m+1,  a>r,
H subset J=J1(beta) normal in P=J U_B^m,
[P:J]=2,  N=K^times J,
pi=c-Ind_N^G(Lambda),  Lambda|J=eta.
\end{Verbatim}

The field K normalizes P. Consequently the following are actual groups
and actual smooth representations:

\begin{Verbatim}[fontsize=\small,breaklines=true,breakanywhere=true]
M=K^times P,
Sigma=Ind_N^M(Lambda),
dim Sigma=4,
Sigma|P=Ind_J^P(eta),
pi is isomorphic to c-Ind_M^G(Sigma).
\end{Verbatim}

The representation Sigma is irreducible, and M is compact modulo the
center of G. Thus the enlarged carrier has an exact arithmetic inducing
datum for the same candidate, not merely an occurrence in its
restriction.

\hypertarget{rank4_stages-normalization-by-the-actual-parameter-field}{%
\subsection{1. Normalization by the actual parameter
field}\label{rank4_stages-normalization-by-the-actual-parameter-field}}

The units O\_K\^{}times equal U\_K\^{}1 because the residue field is
F\_2. They belong to J and hence normalize P. It remains to check a
uniformizer k of K. This element has A-grade one and normalizes A.

The actual-field approximation estimate proved in Section~2 of the
carrier note gives

\begin{Verbatim}[fontsize=\small,breaklines=true,breakanywhere=true]
k=b+e,  b in Q^1,  e in P_A^(n-r+1).
\end{Verbatim}

The leading grade of b is the nonzero centralizing line of the primitive
derived parameter. Over F\_2 it is exactly the leading grade of R, the
period-two E-chain shift. Therefore

\begin{Verbatim}[fontsize=\small,breaklines=true,breakanywhere=true]
b=R v,  v in 1+Q^1.
\end{Verbatim}

In particular b is invertible and normalizes both the complete E-chain
and the complete interleaved F-chain. We can write

\begin{Verbatim}[fontsize=\small,breaklines=true,breakanywhere=true]
k=b(1+epsilon),  epsilon in P_A^(n-r).
\end{Verbatim}

For u in U\_B\^{}m, the usual filtration commutator bound gives

\begin{Verbatim}[fontsize=\small,breaklines=true,breakanywhere=true]
(1+epsilon)u(1+epsilon)^(-1)u^(-1)
   in U_A^(n-r+m).
\end{Verbatim}

This group is contained in U\_A\^{}a subset J, because

\begin{Verbatim}[fontsize=\small,breaklines=true,breakanywhere=true]
n-r+m=2a-m-1>a.
\end{Verbatim}

Conjugation by b preserves both relevant filtrations. Since b normalizes
U\_B\^{}m, it follows that

\begin{Verbatim}[fontsize=\small,breaklines=true,breakanywhere=true]
k U_B^m k^(-1) subset U_B^m J=P.
\end{Verbatim}

Also kJk\^{}(-1)=J, as K\^{}times normalizes its actual simple-character
groups. Thus kPk\^{}(-1) is contained in P and contains J. Its index
over J is two, equal to {[}P:J{]}, so it equals P. This proves
normalization by k, its inverse, and hence all K\^{}times.

No equality or inclusion between the fields E and K was used. The
element b belongs to the actual outer centralizer B and only
approximates k to the stated precision.

\hypertarget{rank4_stages-the-two-inducing-groups-have-index-two}{%
\subsection{2. The two inducing groups have index
two}\label{rank4_stages-the-two-inducing-groups-have-index-two}}

Since K\^{}times normalizes P, the product M=K\^{}times P is a group.
Both P/J and the conjugation action of K\^{}times on it are defined; P/J
has order two, whose automorphism group is trivial. Therefore

\begin{Verbatim}[fontsize=\small,breaklines=true,breakanywhere=true]
[K^times,P] subset J.
\end{Verbatim}

Together with J normal in P and K\^{}times normalizing J, this proves
N=K\^{}times J normal in M. The intersection K\^{}times cap P is
U\_K\^{}1, which is contained in J, since P subset U\_A\^{}1.
Consequently

\begin{Verbatim}[fontsize=\small,breaklines=true,breakanywhere=true]
[M:N]=[P:J]=2.
\end{Verbatim}

The quotient \(K^\times/F^\times\) is compact because K/F is a
finite totally ramified extension, and P is compact. Hence M is compact
modulo F\^{}times. It is open in G because it contains the open subgroup
J.

\hypertarget{rank4_stages-the-enlarged-inducing-representation-is-irreducible}{%
\subsection{3. The enlarged inducing representation is
irreducible}\label{rank4_stages-the-enlarged-inducing-representation-is-irreducible}}

Let u in P represent its nontrivial coset modulo J. The carrier note
proved that theta\^{}u differs from theta, and therefore eta\^{}u
differs from eta. It follows that Lambda\^{}u differs from Lambda, since
their restrictions to J already differ. Applying finite-index Clifford
theory to N normal in M gives

\begin{Verbatim}[fontsize=\small,breaklines=true,breakanywhere=true]
Sigma=Ind_N^M(Lambda) irreducible,
dim Sigma=2 dim Lambda=4.
\end{Verbatim}

There is one N--P double coset in M, and N cap P=J. Thus restriction to
P gives the actual four-dimensional Heisenberg carrier:

\begin{Verbatim}[fontsize=\small,breaklines=true,breakanywhere=true]
Sigma|P is isomorphic to Ind_J^P(eta).
\end{Verbatim}

Its central character on F\^{}times is unchanged from Lambda, because
F\^{}times is central and lies in N.

\hypertarget{rank4_stages-exact-recovery-by-induction-in-stages}{%
\subsection{4. Exact recovery by induction in
stages}\label{rank4_stages-exact-recovery-by-induction-in-stages}}

Induction in stages for these open subgroups gives

\begin{Verbatim}[fontsize=\small,breaklines=true,breakanywhere=true]
c-Ind_M^G(Sigma)
   is isomorphic to c-Ind_M^G(Ind_N^M Lambda)
   is isomorphic to c-Ind_N^G(Lambda)
   =pi.
\end{Verbatim}

The second induction is finite-index, so ordinary and compact induction
from N to M agree. This is an exact equality of the resulting
isomorphism classes and applies to the same fixed Lambda, including its
noncompact scalar choice. It introduces no new candidate and no H5
premise.

\hypertarget{rank4_stages-scope}{%
\subsection{Scope}\label{rank4_stages-scope}}

This establishes the arithmetic inducing datum and exact recovery for
the enlarged carrier. It leaves unchanged the distinction between the
original two-dimensional BK quotient J/H and the new four-dimensional
actual quotient P/H0. It also leaves unchanged the source-phase and
literal extraspecial-center limitations recorded in the carrier note.

Dependency pin: \path{RANK4_ENLARGED_CARRIER167.md} SHA256
\path{1517e26138776fe0f32d74a174d309ec5dd6a0f572a979473d8af3d20b7a4155}.

\clearpage
\section{Exact enlarged-carrier recovery and original-contract comparison}\label{app:lowrank_contract}

Evidence: a bounded contract audit and elementary conditional recovery
lemmas. This note does not prove that the proposed compact enlargement
is a group, that its induced representation is Heisenberg, or that the
source markings attach to it. Those are the arithmetic inputs to be
verified by the separate rank-two and rank-four constructions.

\textbf{Verdict:} an enlarged compact carrier is not, merely by being
Heisenberg or an equivalent type, a new instance of the original H2
quotient. It can serve as a common assembly carrier for the already
supplied H2 targets if its actual occurrence, marked attachment, and
same-candidate recovery are proved and retained. The recovery mechanism
differs from the finite sector idempotent used when the packet lies
inside the original inducing carrier.

\hypertarget{lowrank_contract-what-original-h2-literally-requires}{%
\subsection{1. What original H2 literally
requires}\label{lowrank_contract-what-original-h2-literally-requires}}

Original Section~10 fixes an actual maximal simple character theta, its
actual groups H\_theta¹ subset J\_theta¹ subset J\_theta, and

\begin{Verbatim}[fontsize=\small,breaklines=true,breakanywhere=true]
V_theta=J_theta¹/H_theta¹.
\end{Verbatim}

Its Heisenberg representation is the particular eta\_theta containing
theta. Definition~11.1 supplies an equivariant symplectic isometry from a
source envelope to that type-side quotient, together with the central
character, and at p=2 the compatible quadratic refinements and
central-extension lift. Standard Fact 10.2 separately supplies the
canonical beta extension of this simple-character datum.

Consequently the following implications are invalid without more proof:

\begin{Verbatim}[fontsize=\small,breaklines=true,breakanywhere=true]
Xi is a finite Heisenberg representation
   therefore Xi=eta_theta' for a maximal simple character theta';

(P,Xi) is equivalent to a known arithmetic type
   therefore it has that type's J¹/H¹ or canonical beta extension.
\end{Verbatim}

An equivalence of generated representation categories does not identify
the particular finite group quotient in Definition~11.1. If a fixed
conductor and torsion profile forces every maximal simple-character
carrier in the candidate to have dimension one or two, constructing a
larger Xi in that candidate does not refute the dimension theorem. Xi is
another arithmetic carrier.

There is nevertheless an important distinction between \textbf{replacing
an H2 target} and \textbf{assembling already supplied H2 targets}. The
accepted sector construction itself uses a finite sector that need not
be a whole maximal simple-character carrier. Its old actual break
targets remain unchanged and are attached to that sector. An enlarged
P-carrier can play the analogous assembly role, provided the complete
dictionary and recovery below are proved. Calling it merely an arbitrary
vector space would understate that result; calling it a new canonical
maximal simple type would overstate it.

\hypertarget{lowrank_contract-precise-conditional-setup}{%
\subsection{2. Precise conditional
setup}\label{lowrank_contract-precise-conditional-setup}}

Let G=GL\_n(F), Z=F×I\_n, and let the already constructed candidate be

\begin{Verbatim}[fontsize=\small,breaklines=true,breakanywhere=true]
pi=c-Ind_N^G Lambda,
\end{Verbatim}

with its fixed actual extended type and fixed central character omega.
Let J be a compact open normal subgroup of N such that

\begin{Verbatim}[fontsize=\small,breaklines=true,breakanywhere=true]
eta=Lambda|J is irreducible,       N/J is infinite cyclic.
\end{Verbatim}

In the proposed F2-residue, full-degree totally ramified setting, J is
the actual compact type group J¹=J, not a new group called J merely to
fit this notation. Let P be an actual compact open subgroup of G
satisfying

\begin{Verbatim}[fontsize=\small,breaklines=true,breakanywhere=true]
J subset P,          N intersect P=J.
\end{Verbatim}

Define the genuine finite-dimensional compact induction

\begin{Verbatim}[fontsize=\small,breaklines=true,breakanywhere=true]
Xi=Ind_J^P eta.                                       (1)
\end{Verbatim}

Neither irreducibility nor Heisenberg structure of Xi is assumed to
follow from its dimension. They must be proved separately when needed.

If the known eta intertwining satisfies I\_G(eta)=N, then Xi is
irreducible: the finite Mackey endomorphism formula gives

\begin{Verbatim}[fontsize=\small,breaklines=true,breakanywhere=true]
End_P(Ind_J^P eta)=End_J(eta)=C,
\end{Verbatim}

because no nontrivial double coset in P lies in I\_G(eta).
Finite-dimensional smooth representations of compact P factor through a
finite group and are semisimple, so scalar endomorphisms imply
irreducibility. This criterion requires its actual group and
intertwining hypotheses; dimension alone is not a substitute.

\hypertarget{lowrank_contract-xi-occurs-in-the-same-pi-and-can-be-an-equivalent-compact-type}{%
\subsection{3. Xi occurs in the same pi, and can be an equivalent
compact
type}\label{lowrank_contract-xi-occurs-in-the-same-pi-and-can-be-an-equivalent-compact-type}}

The subspace of the actual compact-induction model of pi consisting of
functions supported on NP is stable under P. Mackey restriction
identifies that subspace with

\begin{Verbatim}[fontsize=\small,breaklines=true,breakanywhere=true]
Ind_(N intersect P)^P(Lambda|_(N intersect P))=Xi.     (2)
\end{Verbatim}

Thus Xi is an actual compact P-subrepresentation of the same pi, with no
change of arithmetic candidate or noncompact scalar.

Induction in stages also gives the exact G-isomorphism

\begin{Verbatim}[fontsize=\small,breaklines=true,breakanywhere=true]
c-Ind_P^G Xi isomorphic to c-Ind_J^G eta.             (3)
\end{Verbatim}

If (J,eta) is already a type for a stated Bernstein block and Xi is
irreducible, then (P,Xi) detects the same irreducible representations:

\begin{Verbatim}[fontsize=\small,breaklines=true,breakanywhere=true]
Hom_P(Xi,V) isomorphic to Hom_J(eta,V).
\end{Verbatim}

This is the precise equivalent-type statement that follows from (1). It
does not assert that Xi is the eta of another maximal simple character,
or attach a canonical beta extension to Xi. The known canonical
extension in Lambda remains the one originally supplied.

\hypertarget{lowrank_contract-exact-recovery-before-fixing-the-center-a-quotient}{%
\subsection{4. Exact recovery before fixing the center: a
quotient}\label{lowrank_contract-exact-recovery-before-fixing-the-center-a-quotient}}

Choose u in N mapping to a generator of N/J, and retain the actual
matrix Lambda(u). Use the model

\begin{Verbatim}[fontsize=\small,breaklines=true,breakanywhere=true]
X=c-Ind_J^G eta,
f(jg)=eta(j)f(g),
\end{Verbatim}

with G acting by right translation. Define the invertible G-equivariant
shift operator

\begin{Verbatim}[fontsize=\small,breaklines=true,breakanywhere=true]
(Tf)(g)=Lambda(u)^(-1)f(ug).                          (4)
\end{Verbatim}

Normality of J and the actual intertwining relation for Lambda(u) ensure
that Tf has the same J-covariance. Its normalization explicitly retains
the fixed noncompact extension; changing Lambda(u) changes the relation
below.

Define

\begin{Verbatim}[fontsize=\small,breaklines=true,breakanywhere=true]
(Qf)(g)=sum_(k in Z) Lambda(u)^(-k)f(u^k g).         (5)
\end{Verbatim}

The sum is finite on each argument and has the required compact support
modulo N. It is N-covariant and therefore lies in pi. It is surjective
and satisfies Q(Tf)=Q(f). Decomposing support into N/J-orbits reduces
its kernel to the elementary augmentation kernel on finitely supported
sequences on Z. That kernel consists exactly of consecutive differences.
Hence

\begin{Verbatim}[fontsize=\small,breaklines=true,breakanywhere=true]
X/(T-1)X isomorphic to pi.                           (6)
\end{Verbatim}

Through (3), this is a recovery of pi from the enlarged compact carrier
plus the original noncompact relation. It does not select the correct
relation from the finite Heisenberg packet alone.

In general this quotient cannot be replaced by a G-equivariant split
inclusion before central normalization. A nonzero vector of X has finite
support modulo compact J. The valuation of its determinant support is a
finite subset of Z. A scalar uniformizer z=tI\_n shifts these valuations
by n; consequently X has no nonzero eigenvector for z. Every vector of
pi is such an eigenvector, with eigenvalue omega(t). Therefore no
nonzero G-map pi -\textgreater{} X exists. In particular (6) does not
split.

This is why the Block159 sector summand projector cannot simply be
quoted: there the finite sector was inside Lambda and its stabilizer
contained the noncompact center. Here Xi is larger than Lambda and lies
in pi\textbar P, with P compact. These are different induction problems.

\hypertarget{lowrank_contract-after-retaining-the-center-a-finite-hecke-idempotent}{%
\subsection{5. After retaining the center: a finite Hecke
idempotent}\label{lowrank_contract-after-retaining-the-center-a-finite-hecke-idempotent}}

There is also an exact summand version once the already fixed central
character is imposed. Put J\_Z=ZJ and P\_Z=ZP. Assume, as in the
proposed totally ramified full-degree setting, that

\begin{Verbatim}[fontsize=\small,breaklines=true,breakanywhere=true]
N/J_Z is finite cyclic of order ell
\end{Verbatim}

(here ell=n). Let Lambda\_0=Lambda\textbar J\_Z. Because Xi occurs in
the same pi, its action on Z intersect P agrees with omega. It therefore
extends to P\_Z by letting Z act through omega; call this Xi\_omega.
Induction in stages gives

\begin{Verbatim}[fontsize=\small,breaklines=true,breakanywhere=true]
Y=c-Ind_(P_Z)^G Xi_omega
  isomorphic to c-Ind_(J_Z)^G Lambda_0.              (7)
\end{Verbatim}

Using the same normalized shift (4), now on Y, one has T\^{}ell=1.
Indeed u\^{}ell belongs to J\_Z and the J\_Z covariance cancels
Lambda(u)\^{}ell. The finite operator

\begin{Verbatim}[fontsize=\small,breaklines=true,breakanywhere=true]
e_Lambda=(1/ell) sum_(j=0)^(ell-1) T^j               (8)
\end{Verbatim}

is a G-equivariant idempotent whose image is exactly the fixed pi. To
see this directly, N-covariant functions embed in J\_Z-covariant
functions because N/J\_Z is finite. On that embedded pi the finite
version of (5) equals ell times the identity. Its average is precisely
(8), and its image consists exactly of the N-covariant functions.

Equivalently, Ind\_(J\_Z)\^{}N Lambda\_0 is Lambda tensored with the
regular representation of N/J\_Z. Formula (8) selects the Lambda summand.
It uses the already fixed Lambda(u), including its root choice, so it
preserves the same candidate. Other root choices select the other
extensions in this finite induction. They are not interchangeable input
data.

This is a finite Hecke/induction projector after central-character
normalization. It is not the original finite image-algebra projector
onto a sector of eta. A final artifact must retain (7), (8), and the
normalization Lambda(u), rather than describing recovery merely as ``the
carrier is an equivalent type.''

\hypertarget{lowrank_contract-the-exact-remaining-admissibility-checks}{%
\subsection{6. The exact remaining admissibility
checks}\label{lowrank_contract-the-exact-remaining-admissibility-checks}}

For either proposed low-rank enlargement, the new proof must supply all
of the following concrete inputs:

\begin{enumerate}
\def\labelenumi{\arabic{enumi}.}
\tightlist
\item
  \textbf{Actual compact group:} prove that the proposed P is a
  subgroup, is compact open, contains the actual J, and satisfies P
  intersect N=J. Compute {[}P:J{]} and the resulting Xi dimension.
\item
  \textbf{Actual carrier:} verify eta=Lambda\textbar J, its actual
  intertwining if the irreducibility criterion is used, and the Mackey
  realization (2).
\item
  \textbf{Actual finite Heisenberg structure:} exhibit the scalar
  preimage in P and the resulting elementary abelian projective
  quotient. Compute its represented commutator and prove nondegeneracy.
  An irreducible carrier of dimension two or four need not automatically
  have this structure.
\item
  \textbf{Original attachments:} attach the unchanged supplied actual H2
  break targets, retaining their full source character values and dyadic
  finite extension/lift/square markings. The precise
  transported-versus-native distinction remains the one proved in
  \path{DYADIC_MARKINGS167.md}.
\item
  \textbf{Recovery data:} retain the actual Lambda, central character
  and noncompact generator normalization, and prove the isomorphisms and
  quotient/idempotent above for those exact groups.
\end{enumerate}

When these are supplied, the result is stronger than occurrence of an
unrelated finite vector space: it is a marked arithmetic carrier in the
same candidate, equivalent through explicit induction and recovering
that candidate exactly. For C=\{1\} and one vertex, it can repair the
failed low-rank \textbf{assembly-carrier} realization under the retained
scalar-pushout contract. The original H2 targets remain their original
simple-character quotients; they are not renamed Xi.

If the intended target instead requires the source packet itself to
equal J\_theta'¹/H\_theta'¹ for a maximal simple character in this same
candidate, the enlargement does not solve it. The retained low-rank
dimension obstruction still applies. A separate theorem identifying Xi
with such an eta would be needed, and under the fixed obstructed profile
it cannot be inferred from type equivalence.

No H5 condition is added or discharged. If H5 is subsequently supplied,
it must concern this fixed recovered candidate. Nontrivial stabilizer
actions, canonical native quadratic comparison, and arbitrary transport
graphs remain outside this bounded argument.

\hypertarget{lowrank_contract-evidence-and-self-review}{%
\subsection{Evidence and self-review}\label{lowrank_contract-evidence-and-self-review}}

Read directly: original Section~10 and Definition~11.1 in the raw
manuscript; the Block159 \path{NONINVARIANT_SECTOR_RECONSTRUCTION159.md};
current \path{PROOF167.md} and \path{TARGET167.md}; the retained dyadic marking
theorem. The Mackey, induction, shift-quotient, and central-normalized
projector arguments above are proved conditionally for the displayed
groups, rather than assumed for the proposed congruence formula. No
actual low-rank P computation or new Lean execution is claimed in this
contract note.


\clearpage
\bibliographystyle{amsplain}
\bibliography{Paper_XI_Traditional_Final}

\end{document}
